% Adapted from the author's public Erdos Paper IV v1.23 series_template.tex.
% Layout provenance is recorded in 02_validation; the Markdown is semantic authority.
\documentclass[spanish,11pt,a4paper]{article}
\usepackage[margin=1in]{geometry}
\usepackage{fontspec}
\setmainfont{lmroman10-regular.otf}[BoldFont=lmroman10-bold.otf,ItalicFont=lmroman10-italic.otf,BoldItalicFont=lmroman10-bolditalic.otf]
\setsansfont{lmsans10-regular.otf}[BoldFont=lmsans10-bold.otf]
\setmonofont{lmmono10-regular.otf}[Scale=MatchLowercase]
\usepackage{amsmath,amssymb}
\providecommand{\mathscr}{\mathcal}
\usepackage[spanish,es-noshorthands,es-tabla]{babel}
\usepackage{microtype,parskip}
\usepackage{longtable,booktabs,array,calc}
\usepackage{xcolor,xurl}
\usepackage{hyperref,bookmark}
\usepackage{etoolbox,needspace}
\setcounter{secnumdepth}{-\maxdimen}
\setlength{\emergencystretch}{3em}
\providecommand{\tightlist}{\setlength{\itemsep}{0pt}\setlength{\parskip}{0pt}}
\newcounter{none}
\AtBeginEnvironment{longtable}{\Needspace{8\baselineskip}\small}
\pretocmd{\section}{\Needspace{5\baselineskip}}{}{}
\pretocmd{\subsection}{\Needspace{4\baselineskip}}{}{}
\pretocmd{\subsubsection}{\Needspace{4\baselineskip}}{}{}
\urlstyle{same}
\hypersetup{pdftitle={Descomposiciones en triángulos II: Absorción robusta para triángulos tipados},pdfauthor={Juan Pablo Traverso Gianini},
pdflang={es-ES},colorlinks=true,linkcolor=blue,urlcolor=blue,citecolor=blue}
\title{Descomposiciones en triángulos II: Absorción robusta para triángulos tipados}
\author{Juan Pablo Traverso Gianini\\
\small Investigador independiente, Santiago, Chile\\
\small \href{mailto:jtraverso@gmail.com}{jtraverso@gmail.com}\\
\small ORCID: 0009-0003-6068-4096}
\date{}
\begin{document}
\raggedbottom
\maketitle
\textbf{Paper I de la serie:} \emph{Descomposiciones en triángulos I: Formalización del teorema de Nash--Williams en Lean} {[}1{]}.

\textbf{Paper II de la serie:} \emph{Descomposiciones en triángulos II: Absorción robusta para triángulos tipados} (este artículo).

\textbf{Preprint 1.0 --- 20 de octubre de 2026.}

\textbf{DOI:} https://doi.org/10.5281/zenodo.23502224.

\subsection{Resumen}\label{resumen}

Estudiamos descomposiciones en triángulos que deben utilizar multiconjuntos prescritos de etiquetas de las partes. Las condiciones necesarias de divisibilidad registran tanto los conteos globales de tipos de aristas como el perfil de grados en cada vértice. Demostramos un teorema de descomposición para grafos típicos bajo una condición finita de tres pasos en los enlaces y la condición de pareja prohibida (3.2) y deducimos los resultados para cuatro partes casi completas utilizados en Paper I. Un teorema separado trata grafos semidensos de dos partes con triángulos mixtos. Las pruebas combinan correcciones reales que respetan los perfiles locales, certificados enteros acotados y un absorbedor fijo elegido antes del residuo. La división de vértices y las correcciones de paridad, equilibrio, carga y torsión extienden la construcción más allá de los enlaces de tres pasos; un catálogo cíclico fijo trata las familias con leyes de conservación. Junto con una clasificación finita exacta, estas construcciones establecen la propiedad para grafos típicos de las 765 familias admisibles de cuatro etiquetas, salvo permutación. Presentamos los argumentos de construcción y su ensamblaje cuantitativo, distinguiendo los datos de los certificados finitos de los argumentos universales que los utilizan.

\textbf{Palabras clave:} descomposición en triángulos; grafos con partición; absorción; retículos enteros; teorema de Nash--Williams.

\subsection{1. Introducción}\label{introducciuxf3n}

Asignemos etiquetas a los vértices de un grafo y fijemos los multiconjuntos de tres etiquetas permitidos para sus triángulos. ¿Cuándo podemos dividir sus aristas en triángulos de esos tipos? Incluso con dos etiquetas, las restricciones van más allá de las condiciones usuales de que todos los grados sean pares y el número de aristas sea divisible por tres. También restringen el número de aristas entre partes y los grados de cada vértice hacia cada parte. Si el grafo tiene estructura interna adicional, pueden aparecer obstrucciones que esos conteos no detectan.

El problema aparece en las ramas extremales de la demostración del teorema de Nash--Williams de Delcourt y Postle {[}2{]}. En ellas se obtienen grafos residuales con una partición de sus vértices y se aplica el Teorema 13.16 de ese artículo. Para demostrar dicho teorema, los autores remiten a su manuscrito \emph{Refined absorption for partitioned hypergraphs and beyond} {[}3{]}, que figura como referencia 17 en {[}2{]}. Al preparar el desarrollo de origen no hemos encontrado una versión pública de ese manuscrito. Por tanto, no identificamos los resultados que siguen con un enunciado que no hemos podido consultar.

Paper I {[}1{]} sigue la división en casos de Delcourt y Postle. Aquí proporcionamos resultados de descomposición para los grafos residuales que necesita esa prueba, junto con una formulación más general para familias de tipos que satisfacen condiciones adecuadas. Los resultados compartidos por los dos artículos son los siguientes.

{\begin{longtable}[]{@{}
  >{\raggedright\arraybackslash}p{(\linewidth - 4\tabcolsep) * \real{0.3333}}
  >{\raggedright\arraybackslash}p{(\linewidth - 4\tabcolsep) * \real{0.3333}}
  >{\raggedright\arraybackslash}p{(\linewidth - 4\tabcolsep) * \real{0.3333}}@{}}
\toprule\noalign{}
\begin{minipage}[b]{\linewidth}\raggedright
Resultado de este artículo
\end{minipage} & \begin{minipage}[b]{\linewidth}\raggedright
Grafo o datos de partida
\end{minipage} & \begin{minipage}[b]{\linewidth}\raggedright
Aplicación en Paper I
\end{minipage} \\
\midrule\noalign{}
\endhead
\bottomrule\noalign{}
\endlastfoot
Corolario 4.2 & Grafos de cuatro partes casi completos y casi equilibrados para \(F_1\)--\(F_3\) & Ramas extremales tipadas; \(F_3\) da el corolario de cuatro colores \\
Teorema C & Grafos semidensos de dos partes para \(F_0\), alejados de dos modelos equilibrados en al menos una parte & Rama que pasa por el Teorema 14.6 de Delcourt--Postle \\
Lema SC & Un certificado firmado, vecinos comunes y suficientes bisagras físicas & Absorbedor terminal de la rama no extremal \\
\end{longtable}
}

La ronda de cobertura protegida y su iteración se desarrollan en Paper I. Ni esa aplicación ni los tres resultados de la tabla requieren una clasificación completa de las familias de tipos.

El método de absorción procede de la absorción refinada {[}4{]}. La formulación para grafos típicos usa pesos regulares sobre los triángulos y una condición de extensión mediante vecindades comunes, en el contexto de los problemas de descomposición etiquetada estudiados por Keevash {[}5{]}. Enunciamos estas condiciones directamente para grafos; no tomamos como hipótesis implícita un teorema sobre un complejo etiquetado sin especificar. La contribución que desarrollamos es la construcción de los certificados y absorbedores requeridos para las familias indicadas, junto con las deducciones que permiten aplicarlos a los grafos residuales de Paper I.

Comenzamos con las definiciones y una condición finita sobre los enlaces. La sección 4 deduce el resultado para grafos casi completos y presenta las tres familias utilizadas en Paper I. La sección 5 separa las hipótesis distintas que necesita \(F_0\). La sección 6 contiene la construcción de absorción a partir de certificados firmados. La sección 7 explica las obstrucciones de retículo físico y de pesos. La sección 8 realiza físicamente el núcleo con leyes de conservación y deduce su aplicación a las diecisiete familias adicionales. La Sección 9 enuncia el teorema para todas las familias. El Apéndice G enuncia los resultados auxiliares y su alcance; el Apéndice H demuestra sus construcciones y estimaciones cuantitativas.

Resultados anteriores de descomposición abordan varias partes de este marco. Barber, Kühn, Lo, Osthus y Taylor {[}11{]} estudian descomposiciones en cliques de grafos multipartitos equilibrados densos y su relación con completar cuadrados latinos; Montgomery {[}12{]} obtiene descomposiciones fraccionarias en cliques de grafos multipartitos equilibrados densos. Esos resultados tratan cliques arcoíris, mientras que nuestras familias pueden repetir una etiqueta de parte y permitir varios tipos de triángulos. La absorción iterada convierte descomposiciones aproximadas en exactas en el trabajo de Barber, Kühn, Lo y Osthus {[}13{]} y en el marco de diseños de hipergrafos de Glock, Kühn, Lo y Osthus {[}14{]}. Los teoremas de diseños de Keevash {[}5,15{]} proporcionan el marco más amplio de regularidad, extensión y retículos. Son antecedentes del método, no consecuencias de este artículo. La Sección 2.3 precisa los límites de nuestra comparación con {[}5{]}.

Los argumentos de empaquetamiento pertenecen a la tradición del nibble de Rödl {[}16{]}, Frankl--Rödl {[}17{]} y Pippenger--Spencer {[}18{]}. Haxell--Rödl {[}19{]} comparan tamaños de empaquetamientos enteros y fraccionarios en grafos densos; esa comparación global no proporciona por sí sola la cota local del residuo utilizada aquí. Ehard--Glock--Joos {[}20{]} dan un teorema ponderado de emparejamientos seudoaleatorios con sus propias hipótesis de grados, codegrados y pesos. No lo invocamos sin esas hipótesis: los Apéndices H.7--H.8 proporcionan las estimaciones locales y los parámetros específicos utilizados en este artículo.

\subsubsection{1.1. Cómo leer la demostración}\label{cuxf3mo-leer-la-demostraciuxf3n}

El texto principal enuncia los resultados, demuestra las deducciones para grafos casi completos y el lema del absorbedor finito, y explica la realización física y la clasificación. Los Apéndices A--F contienen las construcciones extensas. Comparten un principio: la reserva y su catálogo completo se fijan antes de conocer el residuo. Por eso se controla la carga de todos los elementos del catálogo y no sólo la de los que utiliza un residuo particular.

\Needspace{21\baselineskip}
{\begin{longtable}[]{@{}
  >{\raggedright\arraybackslash}p{(\linewidth - 4\tabcolsep) * \real{0.3333}}
  >{\raggedright\arraybackslash}p{(\linewidth - 4\tabcolsep) * \real{0.3333}}
  >{\raggedright\arraybackslash}p{(\linewidth - 4\tabcolsep) * \real{0.3333}}@{}}
\toprule\noalign{}
\begin{minipage}[b]{\linewidth}\raggedright
Tarea
\end{minipage} & \begin{minipage}[b]{\linewidth}\raggedright
Argumento
\end{minipage} & \begin{minipage}[b]{\linewidth}\raggedright
Aplicación
\end{minipage} \\
\midrule\noalign{}
\endhead
\bottomrule\noalign{}
\endlastfoot
Corregir pesos reales conservando perfiles locales & Apéndices A.2--A.3 & Teorema A y ensamblaje por componentes \\
Convertir certificados enteros en absorbedores privados & Apéndices A.4--A.6 & Enlaces de tres pasos, BAL y familias conservadas \\
Refinar y eliminar cargas globales & Apéndice B & Teorema A \\
Obtener certificados físicos y pesos exactos en el grafo semidenso & Apéndice C & Teorema C \\
Dividir enlaces desconectados y corregir invariantes & Apéndices D--E & Teorema D \\
Construir un catálogo conservado de carga total acotada & Apéndice F y sección 8 & Diecisiete casos conservados \\
Demostrar las construcciones auxiliares y sus constantes & Apéndice H & B.3--B.4, C.5a, G.2 y G.4 \\
\end{longtable}
}

Las letras de los apéndices son independientes de los Teoremas A--D. En ellos \(F\) denota la misma familia que \(\mathcal F\) en el texto principal, y sus ecuaciones llevan el prefijo correspondiente. Las tablas finitas están en las fuentes congeladas, ubicadas en {[}8, Secciones 7--8{]}; el resultado para todas las familias utiliza datos explícitos de certificados.

Las versiones inglesa y española contienen el texto principal y los seis apéndices completos, con la misma numeración y las mismas expresiones matemáticas.


\Needspace{9\baselineskip}
\subsection{2. Tipos, divisibilidad y condiciones sobre el grafo}\label{tipos-divisibilidad-y-condiciones-sobre-el-grafo}

Todos los grafos son finitos, simples y no dirigidos. Escribimos \[V(G)=V_0\sqcup\cdots\sqcup V_{t-1},\qquad n_i=|V_i|,\qquad n=|V(G)|,\] y ponemos \(p(v)=i\) cuando \(v\in V_i\). Un \emph{tipo} es un multiconjunto de tres etiquetas de \(\{0,\ldots,t-1\}\). Fijamos una familia \(\mathcal F\) de tipos. Un par \(\{i,j\}\), con repetición permitida, es \emph{admisible} si está contenido en algún elemento de \(\mathcal F\). Un \emph{triángulo de \(\mathcal F\)} es un triángulo real de \(G\) cuyo multiconjunto de etiquetas pertenece a \(\mathcal F\). Una \emph{descomposición en \(\mathcal F\)} es una partición de \(E(G)\) en esos triángulos.

\subsubsection{2.1. Dos niveles de divisibilidad entera}\label{dos-niveles-de-divisibilidad-entera}

Para un tipo \(\tau\), sea \(b_\tau\) su vector de conteos de aristas, indexado por pares no ordenados de etiquetas. Las etiquetas repetidas se cuentan con sus multiplicidades. Por ejemplo, el tipo \(001\) tiene una arista de tipo \(00\) y dos de tipo \(01\).

Si \(i\) aparece en \(\tau\), sea \(a_{\tau,i}\) el vector que cuenta las otras dos etiquetas vistas desde un vértice con etiqueta \(i\). Definimos el \emph{retículo global de tipos} y los \emph{retículos de perfiles locales} por \[\begin{gathered}
\mathcal L_E=\langle b_\tau:\tau\in\mathcal F\rangle_{\mathbb Z},
\\
\mathcal L_i=\langle a_{\tau,i}:\tau\in\mathcal F,\ i\in\tau\rangle_{\mathbb Z}.
\end{gathered}
\tag{2.1}\] Diremos que \(G\) es \emph{integralmente divisible para \(\mathcal F\)} si su vector de conteos de aristas pertenece a \(\mathcal L_E\) y \[\bigl(|N_G(v)\cap V_j|\bigr)_{j=0}^{t-1}\in\mathcal L_{p(v)}
\quad\text{para todo }v\in V(G).
\tag{2.2}\] Enunciaremos por separado la exigencia de que todas las aristas tengan tipos admisibles. Ambas condiciones, (2.1) y (2.2), son necesarias: basta sumar las contribuciones de los triángulos de una descomposición, globalmente y en un vértice fijo, respectivamente.

El \emph{retículo físico de triángulos} conserva las aristas individuales del grafo: \[\Lambda_{\mathcal F}(G)
=\left\langle\mathbf1_{E(T)}:T\text{ es un triángulo de }\mathcal F\text{ en }G\right\rangle_{\mathbb Z}
\subseteq\mathbb Z^{E(G)}.
\tag{2.3}\] Un \emph{certificado firmado} de un vector de aristas \(x\) es una identidad \[x=\sum_T c_T\mathbf1_{E(T)},\qquad c_T\in\mathbb Z.
\tag{2.4}\] Su coste es \(\sum_T|c_T|\). La divisibilidad integral recoge conteos de tipos y perfiles locales; (2.4) da una representación mediante triángulos reales. Son condiciones distintas. La sección 7 muestra una obstrucción elemental que las separa.

\subsubsection{2.2. Pesos regulares y extensión}\label{pesos-regulares-y-extensiuxf3n}

Para \(\omega,c>0\), el grafo tiene \emph{pesos de triángulos \((\omega,c)\)-regulares} si existe un peso \(y(T)\) para cada triángulo de \(\mathcal F\) tal que \[\begin{gathered}
\frac{\omega}{n}\le y(T)\le\frac1{\omega n},\\
\left|\sum_{T\ni e}y(T)-1\right|\le c
\quad(e\in E(G)).
\end{gathered}
\tag{2.5}\] La cota inferior se exige para todos los triángulos físicos permitidos, no solo para aquellos a los que una descomposición fraccional elegida asigna peso positivo.

Para \(h\in\mathbb N\), el grafo tiene \emph{extensión \((\omega,h)\)} si se cumple lo siguiente. Para todo \(S\subseteq V(G)\) con \(|S|\le h\), todo \(A\subseteq S\) y toda etiqueta \(i\) que forme un par admisible con cada etiqueta de \(p(A)\), \[\left|\{v\in V_i\setminus S:A\subseteq N_G(v)\}\right|\ge\omega n.
\tag{2.6}\] Excluir \(S\) permite elegir sucesivamente vértices auxiliares distintos. Obsérvese que (2.6) exige vecinos comunes para conjuntos de tamaño acotado, no solo para pares.

Diremos que una familia tiene la \emph{propiedad de descomposición en grafos típicos} si \[\forall\gamma>0\ \exists h\in\mathbb N,\omega_0>0\ 
\forall\omega\in(0,\omega_0]\ \exists c>0,N\in\mathbb N
\tag{2.7}\] tales que todo grafo integralmente divisible para \(\mathcal F\), de orden \(n\ge N\), cuyas aristas sean admisibles y que satisfaga \(n_i\ge\gamma n\) para cada \(i\), la regularidad (2.5) y la extensión (2.6), admite una descomposición en \(\mathcal F\). Esta definición da nombre a una propiedad precisa; afirmar que una familia la cumple requiere una demostración.

\subsubsection{2.3. Relación con el teorema de diseños etiquetados}\label{relaciuxf3n-con-el-teorema-de-diseuxf1os-etiquetados}

La propiedad (2.7) sigue el esquema de regularidad, extensión y retículos de Keevash {[}5, Teorema 3.1{]}. Ese teorema exige un complejo etiquetado adaptado finito, un sistema elemental de moléculas, pertenencia a su retículo efectivo de descomposición, pesos regulares y extensión con las restricciones de tipos especificadas. Impone además sus propias desigualdades de tamaño y parámetros. La correspondencia siguiente distingue la codificación propuesta de las hipótesis demostradas en este artículo.

{\begin{longtable}[]{@{}
  >{\raggedright\arraybackslash}p{(\linewidth - 2\tabcolsep) * \real{0.5000}}
  >{\raggedright\arraybackslash}p{(\linewidth - 2\tabcolsep) * \real{0.5000}}@{}}
\toprule\noalign{}
\begin{minipage}[b]{\linewidth}\raggedright
Elemento de (2.7)
\end{minipage} & \begin{minipage}[b]{\linewidth}\raggedright
Correspondencia y comprobación necesaria
\end{minipage} \\
\midrule\noalign{}
\endhead
\bottomrule\noalign{}
\endlastfoot
Tipos de triángulos permitidos & Se usa una molécula etiquetada finita por tipo y se conservan las etiquetas repetidas como posiciones distintas. Hay que especificar la acción de reetiquetado, el complejo adaptado y el sistema elemental de coordenadas de aristas. \\
Pesos (2.5) & Se distribuye cada peso entre sus representaciones etiquetadas y se comprueban la normalización y todas las cargas por coordenada de tipo de {[}5, Definición 2.37{]}. Las multiplicidades finitas afectan las constantes. \\
Extensión (2.6) & Las elecciones sucesivas de vecinos comunes dan inmersiones de patrones compatibles acotados. Para invocar {[}5, Definiciones 2.9--2.10{]} hay que identificar cada patrón restringido exigido y traducir sus parámetros de rango y oferta. \\
Divisibilidad (2.1)--(2.2) & Son condiciones agregadas de tipos y perfiles de vértices. Su identificación con la pertenencia al retículo físico efectivo de descomposición es una afirmación separada. \\
\end{longtable}
}

La última distinción importa incluso antes de elegir parámetros. La identificación disponible del retículo del grafo completo se aplica a un blowup completo permitido con al menos cinco vértices en cada clase. Un grafo típico de (2.7) no tiene por qué ser completo. Por ello, esa identificación por sí sola no permite sustituir la hipótesis de retículo de {[}5{]} por (2.1)--(2.2). Las construcciones por familias que siguen proporcionan certificados físicos y absorbedores; esta comparación bibliográfica no las sustituye.

Usamos, por tanto, (2.7) como una propiedad de grafos demostrada aquí para las familias enunciadas. No afirmamos haber probado la equivalencia de todas sus hipótesis con las de {[}5, Teorema 3.1{]}, ni que nuestro enunciado de existencia sea más fuerte que ese teorema general. Las condiciones adicionales de tres pasos y pareja prohibida delimitan nuestra construcción. Las aportaciones que corresponde comparar son las construcciones explícitas por familias, su ensamblaje cuantitativo y las interfaces necesarias en Paper I.


\Needspace{9\baselineskip}
\subsection{3. Una condición finita sobre la familia de tipos}\label{una-condiciuxf3n-finita-sobre-la-familia-de-tipos}

Fijemos cuatro etiquetas. Para una etiqueta raíz \(r\), su enlace tiene las etiquetas como vértices y una arista \(ab\) cuando \(\{r,a,b\}\in\mathcal F\). Se conservan los lazos, pues las etiquetas repetidas son tipos legítimos. Un vértice del enlace es activo si es incidente con alguna arista.

La familia tiene \emph{enlaces de tres pasos} si, para cada \(r\) y cada par de vértices activos \(a,b\) de su enlace, existen etiquetas \(v,w\) tales que \[\{r,a,v\},\quad\{r,v,w\},\quad\{r,w,b\}\in\mathcal F.
\tag{3.1}\] Los extremos pueden coincidir. Se exige un paseo, no un camino simple. Pedimos además que cada etiqueta tenga una pareja prohibida: \[\forall i\ \exists j\quad\{i,j\}\text{ no es admisible}.
\tag{3.2}\] Esta condición restringe el alcance del teorema siguiente; no se afirma que sea necesaria para toda descomposición. En la prueba proporciona una cota de grado: la clase prohibida de un vértice no contiene vecinos suyos. La Sección B.6 utiliza esta cota en el empaquetamiento inicial con cargas próximas a uno.

\subsubsection{Teorema A (grafos típicos con enlaces de tres pasos)}\label{teorema-a-grafos-tuxedpicos-con-enlaces-de-tres-pasos}

Sea \(\mathcal F\) una familia de multiconjuntos de tres elementos sobre cuatro etiquetas que satisface (3.1) y (3.2). Entonces \(\mathcal F\) tiene la propiedad de descomposición en grafos típicos (2.7).

El teorema va más allá de los grafos casi completos: (2.5) y (2.6) no exigen que todos los pares admisibles estén presentes. A la inversa, no podemos omitir estas hipótesis por el solo hecho de disponer de muchos dispositivos pequeños de intercambio.

\subsubsection{3.1. Organización de la construcción}\label{organizaciuxf3n-de-la-construcciuxf3n}

La prueba tiene una parte algebraica y otra que realiza las identidades dentro del grafo. La primera debe conservar tanto los totales de aristas por tipo como los perfiles locales. Resolver únicamente las ecuaciones globales dejaría una demanda que podría no ser realizable en cada vértice. La segunda parte representa las identidades mediante triángulos permitidos y elige vértices auxiliares con carga acotada.

La corrección real se aplica al defecto de carga de la familia ponderada de triángulos que se tiene en ese momento. La divisibilidad del grafo y los perfiles de esa familia ponderada sitúan el defecto en el espacio compatible requerido. Un inverso acotado sobre el grafo completo admisible proporciona allí una corrección. Para trasladarla al grafo real se usa la condición de extensión, teniendo en cuenta todos los vértices y aristas ya reservados. Los certificados enteros y el absorbedor de la familia resuelven después el residuo divisible. Ese absorbedor debe fijarse antes de seleccionar el residuo. En particular, no bastan certificados separados de coste acotado: sus vértices auxiliares podrían acumular una carga excesiva. La construcción elige un único catálogo físico y una reserva para todos los residuos permitidos y controla la incidencia en cada vértice. La Sección 8 demuestra este paso de realización para el núcleo con leyes de conservación.

El orden de dependencias es, por tanto, \[\begin{gathered}\text{construcción para la familia}\ \longrightarrow
\text{reserva protegida y absorbedor}\ \\\longrightarrow
\text{empaquetamiento y residuo}\ \longrightarrow
\text{absorción}.
\end{gathered}\tag{3.3}\] El paso del empaquetamiento al residuo no permite elegir un absorbedor nuevo. Este orden de cuantificadores forma parte de la construcción.

\textbf{Demostración del Teorema A: estructura y parámetros.} El Apéndice A construye primero un inverso real sobre el grafo completo de pares admisibles. La corrección tiene tres etapas: conteos globales por tipo, perfiles locales y el vector de aristas restante con perfiles nulos. El traslado al grafo real utiliza conjuntos de vecinos comunes acotados. Los certificados enteros se obtienen por separado; un inverso real no bastaría para establecer divisibilidad entera ni descomponibilidad binaria.

El Apéndice B utiliza la identidad del paseo de tres pasos para identificar cada retículo local con el retículo de suma par en sus etiquetas activas. Se conserva el grupo cociente global, incluida su torsión. Se realiza en el grafo un refinador ordinario de piezas pares y un registro finito fijo elimina las cargas globales restantes. El catálogo resultante tiene piezas acotadas e incidencia total acotada sobre \textbf{todos} sus elementos. El Apéndice A incorpora después un absorbedor privado para cada elemento del catálogo, controlando el grado acumulado. Así se obtiene un único absorbedor para todo subgrafo divisible de la reserva.

Finalmente, el Apéndice B.6 elige la tolerancia de corrección y la escala de la reserva antes del grafo. Se fijan la reserva y el absorbedor, se recalculan los pesos regulares después de eliminar sus aristas descomponibles y se aplican los lemas de empaquetamiento y cobertura para dejar un residuo divisible dentro de la reserva. El absorbedor ya fijado cierra la descomposición. Las desigualdades explícitas de los Apéndices A.6 y B.5--B.6 garantizan que cada construcción cabe en los recursos que deja la anterior. Los resultados auxiliares de empaquetamiento, cobertura, refinamiento ordinario y emparejamiento finito se utilizan con sus hipótesis expresas; no suponen la conclusión del Teorema A.


\Needspace{9\baselineskip}
\subsection{4. Grafos casi completos}\label{grafos-casi-completos}

El \emph{grado admisible faltante} de un vértice \(v\) es el número de vértices \(u\ne v\) para los que \(\{p(u),p(v)\}\) es admisible pero \(uv\notin E(G)\). En un grafo casi completo este número estará acotado, para cada vértice, por un entero \(d\le\delta n\).

Para etiquetas \(a,b\) y un tipo \(\tau\), definimos \[q_\tau(a,b)=\frac1n
\sum_{\substack{c:\{a,b,c\}=\tau}}
\bigl(n_c-\operatorname{mult}_c\{a,b\}\bigr).
\tag{4.1}\] La suma recorre etiquetas, no encajes ordenados. Cuenta, de forma normalizada, los posibles terceros vértices en el grafo completo admisible, excluyendo los extremos.

\subsubsection{Teorema B (deducción para grafos casi completos)}\label{teorema-b-deducciuxf3n-para-grafos-casi-completos}

Fijemos una familia \(\mathcal F\) sobre \(t\) etiquetas con la propiedad de descomposición en grafos típicos. Para cada \(\gamma,\omega_1>0\) existen \(\varepsilon,\delta>0\) y \(N\in\mathbb N\) tales que se cumple lo siguiente. Sea \(G\) un grafo de orden \(n\ge N\), con aristas únicamente admisibles, \(n_i\ge\gamma n\) para cada \(i\) y grado admisible faltante a lo sumo \(d\le\delta n\). Supongamos que \(G\) es integralmente divisible para \(\mathcal F\) y que existen pesos de tipos \(w_\tau\) que satisfacen \[\begin{gathered}
\omega_1\le w_\tau\le\omega_1^{-1}\quad(\tau\in\mathcal F),\\
\left|\sum_{\tau\in\mathcal F}w_\tau q_\tau(a,b)-1\right|\le\varepsilon
\quad\text{para cada par admisible }\{a,b\}.
\end{gathered}
\tag{4.2}\] Entonces \(G\) admite una descomposición en \(\mathcal F\).

\emph{Demostración.} Aplicamos (2.7) con \(\gamma\) y obtenemos \(h\) y \(\omega_0\). Elegimos \[\begin{gathered}
\omega=\min\{\omega_1,\omega_0,\gamma/2\},\\
\varepsilon=\frac{c}{4(1+|\mathcal F|/\omega_1)},\\
\delta=\min\left\{\frac\varepsilon4,\frac\gamma{4(h+1)}\right\},
\end{gathered}
\tag{4.3}\] donde \(c>0\) y \(N_0\) son la tolerancia de regularidad y el umbral obtenidos para este \(\omega\). Tomamos \[N\ge\max\left\{N_0,1,
\left\lceil\frac{4(t+h+1)}{\min\{\varepsilon,\gamma\}}\right\rceil\right\}.
\tag{4.4}\]

Para comprobar la extensión, fijamos \(S,A,i\) como en (2.6). A lo sumo \(|S|\) vértices de \(V_i\) quedan excluidos por \(S\), y cada elemento de \(A\) excluye a lo sumo otros \(d\). En consecuencia, el número de opciones es al menos \[n_i-|S|-|A|d\ge\gamma n-h-hd\ge\gamma n/2\ge\omega n.
\tag{4.5}\] La desigualdad central se sigue de (4.3)--(4.4).

Para una arista \(uv\) y un tipo \(\tau\), sea \(t_\tau(u,v)\) el número de terceros vértices reales de ese tipo. Las aristas admisibles que faltan en \(u\) y \(v\), junto con la exclusión de los extremos, dan la cota \[|t_\tau(u,v)-nq_\tau(p(u),p(v))|\le2d+2t\le\varepsilon n.
\tag{4.6}\] El término adicional \(2t\) conserva la cota uniforme de conteo usada en la construcción. Asignamos a cada triángulo físico de tipo \(\tau\) el peso \(w_\tau/n\). Las cotas sobre \(w_\tau\) implican las de (2.5), pues \(\omega\le\omega_1\). Por (4.2) y (4.6), \[\left|\sum_{T\ni uv}\frac{w_{p(T)}}n-1\right|
\le\varepsilon\left(1+\sum_{\tau\in\mathcal F}w_\tau\right)
\le\varepsilon(1+|\mathcal F|/\omega_1)=c/4\le c.
\tag{4.7}\] Así, el mismo grafo satisface regularidad, extensión, divisibilidad integral y las condiciones sobre el tamaño de las partes. Aplicamos (2.7).

\subsubsection{4.1. Pesos de tipos equilibrados}\label{pesos-de-tipos-equilibrados}

Un peso positivo \(w_\tau\) para cada tipo es \emph{equilibrado} si, para cada par admisible, \[\sum_{\substack{c:\{a,b,c\}\in\mathcal F}}w_{\{a,b,c\}}=t.
\tag{4.8}\] Esta es una condición finita sobre la familia. No se deduce automáticamente de la divisibilidad integral de un grafo.

\subsubsection{Corolario 4.1}\label{corolario-4.1}

Supongamos que \(t\ge1\), que \(\mathcal F\) tiene la propiedad de descomposición en grafos típicos y que (4.8) admite una solución positiva. Existen \(\eta,\delta>0\) y \(N\) tales que todo grafo integralmente divisible para \(\mathcal F\), de orden \(n\ge N\), con aristas únicamente admisibles, \[(1/t-\eta)n\le n_i\le(1/t+\eta)n,
\tag{4.9}\] y grado admisible faltante a lo sumo \(d\le\delta n\), admite una descomposición en \(\mathcal F\).

\emph{Demostración.} Como la familia es finita y los pesos son positivos, podemos elegir \(\omega_1>0\) con \(\omega_1\le w_\tau\le\omega_1^{-1}\) para cada tipo. Aplicamos el Teorema B con \(\gamma=1/(2t)\) y este \(\omega_1\). Si \(\varepsilon,N_1,\delta\) son sus constantes, tomamos \[\begin{gathered}
\eta=\min\{\varepsilon/(2t),1/(2t)\},\\
N\ge\max\{N_1,\lceil8t/\varepsilon\rceil\}.
\end{gathered}
\tag{4.10}\] Todas las partes tienen entonces tamaño al menos \(n/(2t)\). Reindexando la suma de (4.2) por las etiquetas que completan el par y usando (4.8), obtenemos \[\left|\sum_\tau w_\tau q_\tau(a,b)-1\right|
\le\sum_{c:\{a,b,c\}\in\mathcal F}
w_{\{a,b,c\}}(\eta+2/n)
=t\eta+2t/n\le\varepsilon.
\tag{4.11}\] Se aplica el Teorema B.

\subsubsection{4.2. Las familias utilizadas en Paper I}\label{las-familias-utilizadas-en-paper-i}

Con etiquetas \(0,1,2,3\), definimos \[\begin{aligned}
\mathcal F_1&=\{002,003,112,113,022,122,033,133\},\\
\mathcal F_2&=\{002,003,112,113,023,123\},\\
\mathcal F_3&=\{012,013,023,123\}.
\end{aligned}
\tag{4.12}\] Cada par admisible tiene exactamente dos etiquetas que lo completan. Por tanto, \(w_\tau=2\) satisface (4.8). La parte activa de cada enlace raíz es un triángulo o un camino de tres vértices con un lazo en cada extremo. La primera forma aparece en todas las raíces de \(F_3\) y en las raíces cero y uno de \(F_2\); los demás enlaces tienen la segunda forma. Ambas admiten un paseo de tres pasos entre cualquier par de vértices, incluidos extremos coincidentes: en el camino, los lazos de los extremos aportan el paso impar; en el triángulo se puede recorrer el triángulo cerrado o avanzar y retroceder por una arista. Se cumple así (3.1). En \(F_1\) y \(F_2\), cero y uno son parejas prohibidas entre sí; en \(F_1\), dos y tres son parejas prohibidas, mientras que en \(F_2\) están prohibidos los pares \(22\) y \(33\). En \(F_3\) están prohibidos todos los pares repetidos. Esto verifica (3.2) sin imponer ninguna condición sobre el grafo.

\subsubsection{Corolario 4.2 (resultado equilibrado de cuatro partes)}\label{corolario-4.2-resultado-equilibrado-de-cuatro-partes}

Para cada \(\mathcal F\in\{\mathcal F_1,\mathcal F_2,\mathcal F_3\}\) existen \(\eta,\delta>0\) y \(N\) tales que todo grafo de orden \(n\ge N\) con aristas únicamente admisibles, grado admisible faltante a lo sumo \(d\le\delta n\), tamaños de partes en \([(1/4-\eta)n,(1/4+\eta)n]\) y divisibilidad integral para \(\mathcal F\) admite una descomposición en \(\mathcal F\).

\emph{Demostración.} El Teorema A da la propiedad para grafos típicos. El peso \(2\) y el Corolario 4.1 dan la conclusión.

Para \(\mathcal F_1\) y \(\mathcal F_2\), Paper I utiliza también la forma idealizada: cada parte tiene tamaño al menos \((1-\varepsilon)n/4\) y cada vértice tiene al menos \(|V_j|-\varepsilon n\) vecinos en cada parte admisible \(V_j\). Para \(\varepsilon>0\) suficientemente pequeño, el Corolario 4.2 cubre este caso. En efecto, sumando las cotas inferiores de las otras tres partes se obtiene la cota superior \((1/4+3\varepsilon/4)n\). Al sumar los grados faltantes hacia las cuatro partes se obtienen a lo sumo \(4\varepsilon n\) vecinos admisibles faltantes. Basta elegir \(\varepsilon\) de modo que \(\varepsilon\le\eta\) y \(4\varepsilon\le\delta\). Solo pueden aparecer aristas de tipos permitidos y la compatibilidad entera sigue siendo una hipótesis de esta forma idealizada.

La aplicación a \(F_3\) es la reducción de cuatro colores. Requiere comprobar las condiciones de los retículos global y local en el grafo residual: una coloración propia no basta para asegurar divisibilidad. Paper I realiza esas comprobaciones y une la descomposición resultante al empaquetamiento de limpieza.


\Needspace{9\baselineskip}
\subsection{5. Un teorema semidenso separado}\label{un-teorema-semidenso-separado}

Sea \(\mathcal F_0=\{001,011\}\). Un triángulo de \(\mathcal F_0\) tiene dos vértices en una parte y uno en la otra. En este caso la divisibilidad integral equivale a \[d_G(v)\text{ par para todo }v,\qquad e_{01}=2(e_{00}+e_{11}).
\tag{5.1}\] Los generadores de tipos de aristas son \((1,2,0)\) y \((0,2,1)\), en el orden \(00,01,11\). En un vértice de la parte cero, los generadores locales son \((1,1)\) y \((0,2)\); su retículo es el conjunto de pares cuya suma es par. En la parte uno se obtiene la condición correspondiente. Esto prueba la equivalencia.

Para un conjunto \(D\), consideremos todas las particiones \(D=X\sqcup(D\setminus X)\) con \(|X|=\lfloor|D|/2\rfloor\). Sus dos \emph{modelos equilibrados} son el grafo bipartito completo entre ambas partes y la unión disjunta de las dos cliques. Diremos que \(G[D]\) está \emph{\(\alpha\)-cerca} de un modelo equilibrado si se puede obtener uno de esos grafos mediante a lo sumo \(\alpha|D|^2\) modificaciones de aristas no ordenadas.

\subsubsection{Teorema C (resultado semidenso de dos partes)}\label{teorema-c-resultado-semidenso-de-dos-partes}

Para cada \(\alpha>0\) existen \(\varepsilon_0>0\) y \(N\) tales que se cumple lo siguiente. Sea \(G\) integralmente divisible para \(\mathcal F_0\), de orden \(n\ge N\), con una partición \(V_0\sqcup V_1\) que satisface \[|V_i|\ge(1-\varepsilon_0)n/2\quad(i=0,1),
\tag{5.2}\] y, para cada \(v\in V_i\), \[(1/4-\varepsilon_0)n\le |N_G(v)\cap V_i|\le(1/4+\varepsilon_0)n,
\qquad |N_G(v)\cap V_{1-i}|\ge(1/2-\varepsilon_0)n.
\tag{5.3}\] Si al menos uno de \(G[V_0]\) y \(G[V_1]\) no está \(\alpha\)-cerca de ninguno de los dos modelos equilibrados, entonces \(G\) admite una descomposición en \(\mathcal F_0\).

El grado interno de (5.3) es aproximadamente la mitad del tamaño de su parte. No estamos en el régimen casi completo admisible del Teorema B. La exclusión de ambos modelos en una de las partes es una hipótesis sustantiva: proporciona la geometría necesaria para la construcción de \(F_0\) y no puede reemplazarse por la divisibilidad agregada.

Contar pares ordenados en lugar de modificaciones no ordenadas duplica la distancia. Por tanto, la condición anterior corresponde a una cota de desacuerdo ordenado \(2\alpha|D|^2\), con el mismo \(\alpha\). Mantener explícito este factor evita cambiar la región excepcional al aplicar el teorema.

\textbf{Demostración del Teorema C: el argumento geométrico separado.} El Apéndice C parte de la exclusión de los dos modelos equilibrados en la parte adecuada. Un argumento de estabilidad para paseos cortos proporciona muchos paseos de longitudes tres y cuatro en ella. Estos permiten construir certificados firmados sobre triángulos mixtos físicos y bisagras explícitas; los conteos conservan las cuatro raíces fijas. Así se dispone de la construcción terminal entera con constantes que dependen sólo del parámetro de no extremalidad.

El resultado fraccional se construye por separado. En el grafo completo de enlaces cruzados, un inverso de incidencia interna corrige simultáneamente las cargas de aristas internas de ambas partes y las cargas cruzadas. Eliminar los pares cruzados ausentes produce una perturbación controlada de ese operador; una corrección convergente da pesos exactos, acotados y no negativos en el grafo real. La estimación de eliminación del Apéndice C.6 permite repetir esta construcción después de reservar el absorbedor. Una segunda construcción de pesos, seguida de la cobertura, deja el residuo divisible para dicho absorbedor. El Apéndice C.6 fija todas las escalas y el umbral de orden. El argumento no utiliza el Teorema A ni el manuscrito no disponible {[}3{]}.

\subsection{6. De certificados firmados a absorbedores finitos}\label{de-certificados-firmados-a-absorbedores-finitos}

Esta sección trata triángulos ordinarios. La construcción no requiere no extremalidad ni una descomposición fraccional.

Para \(u,v,a,b\) distintos, una \emph{bisagra} entre los triángulos \(uva\) y \(uvb\) es un conjunto de aristas \(H\subseteq E(G)\), disjunto de ambos triángulos, tal que \[H\cup\{ua,va\}\quad\text{y}\quad H\cup\{ub,vb\}
\tag{6.1}\] admiten descomposiciones en triángulos. Su orden es el número de vértices incidentes con sus aristas. Las cuatro raíces pertenecen a ese soporte: cubrir una arista del ángulo sin usar la base excluida \(uv\) requiere aristas de la bisagra en sus extremos. Una vez fijadas las raíces, una bisagra de orden \(i\) tiene \(i-4\) vértices libres.

Exigimos que cada par de triángulos que comparten una arista tenga, para algún entero \(i\) con \(6\le i\le C_0\), al menos \(\beta n^{i-4}\) bisagras distintas de orden \(i\), contadas como conjuntos físicos de aristas. Llamaremos a esta condición una \emph{oferta de bisagras \((\beta,C_0)\)}.

\subsubsection{Lema SC (los certificados firmados producen absorbedores)}\label{lema-sc-los-certificados-firmados-producen-absorbedores}

Sea \(G\) de orden \(n\), sean \(\alpha,\beta>0\) y sean \(C_0,K,z\) enteros no negativos. Supongamos que cada par de vértices distintos tiene al menos \(\alpha n\) vecinos comunes y que \(G\) tiene una oferta de bisagras \((\beta,C_0)\). Sea \(P\subseteq E(G)\) con un certificado mediante triángulos reales \[\mathbf1_P=\sum_j c_j\mathbf1_{E(T_j)},\qquad c_j\in\mathbb Z,
\qquad \sum_j|c_j|\le K.
\tag{6.2}\] Supongamos que \(V(P)\subseteq Z\) y \(|Z|\le z\). Ponemos \[m=\binom{C_0+1}{2},\qquad B=z+3K+4KC_0+1.\] Si \[B<\alpha n,\qquad 2^mB<\beta n,
\tag{6.3}\] existe un conjunto de aristas \(A\subseteq E(G)\) tal que \(A\) y \(A\cup P\) admiten descomposiciones en triángulos y \[A\cap E(G[Z])=\varnothing,\qquad |A|\le4K(m+2).
\tag{6.4}\] En particular, \(A\) es disjunto de \(P\) en aristas.

\emph{Demostración.} Fijamos las cuatro raíces de una bisagra y un conjunto prohibido \(U\) que no contiene raíces. Para cada \(w\in U\), hay a lo sumo \(n^{i-5}\) elecciones de vértices adicionales y \(2^m\) conjuntos de aristas que pueden dar una bisagra de orden \(i\) que pase por \(w\). Así, a lo sumo \(2^m|U|n^{i-5}\) bisagras tocan \(U\). La segunda desigualdad de (6.3) dejará una bisagra disponible siempre que \(|U|\le B\).

Expandimos cada coeficiente de (6.2) en \(|c_j|\) ocurrencias firmadas de su triángulo. Hay a lo sumo \(K\) ocurrencias, cada una con tres posiciones de arista. Incluimos todos los vértices de estos triángulos en el conjunto usado inicial, aunque su contribución se cancele después.

Para una ocurrencia \(abc\), elegimos un vecino común nuevo \(x\) de \(a,b\) y luego una bisagra \(H\) entre \(abc\) y \(abx\) cuyos vértices no raíz sean nuevos. En una descomposición de \(H\cup\{ac,bc\}\), sean \(acy\) y \(bcz\) los triángulos que cubren \(ac\) y \(bc\). Las aristas excluidas de los triángulos obligan a que \(y,z\) estén fuera de \(\{a,b,c,x\}\). Además, \(y\ne z\), pues de lo contrario los dos triángulos compartirían \(cy\). Definimos \[W_{ab}=\{ax,bx\},\quad W_{ac}=\{ay,cy\},\quad W_{bc}=\{bz,cz\},
\qquad Y=H\setminus\{ay,cy,bz,cz\}.
\tag{6.5}\] Al retirar \(acy\) y \(bcz\) queda una descomposición de \(Y\). La otra descomposición de la bisagra prueba que \[Y\,\dot\cup\,W_{ab}\,\dot\cup\,W_{ac}\,\dot\cup\,W_{bc}
=H\cup\{ax,bx\}
\tag{6.6}\] también es descomponible. Cada ángulo completa su arista correspondiente en un triángulo. Ninguna arista de este dispositivo de ocurrencia une dos de sus raíces del certificado.

Para cada arista física \(e\), la identidad (6.2) dice que el número de posiciones positivas sobre \(e\) es igual al de posiciones negativas más \(\mathbf1_P(e)\). Emparejamos biyectivamente las posiciones negativas con posiciones positivas de la misma arista, dejando exactamente una posición positiva sin pareja por cada arista de \(P\).

Para cada par emparejado sobre \(uv\), sean \(a_-,a_+\) sus ápices. Elegimos una bisagra nueva \(J\) entre \(uva_-\) y \(uva_+\). Sus dos descomposiciones usan, respectivamente, el ángulo negativo y el positivo. Estas elecciones son disjuntas en aristas de la construcción previa. La base y las cuatro aristas de los ángulos están excluidas de \(J\). El único par restante entre sus raíces une \(a_-\) con \(a_+\); ningún dispositivo de ocurrencia contiene ese par, pues los ápices proceden de ocurrencias distintas. Tampoco lo contiene una bisagra de emparejamiento anterior: cada posición se empareja una sola vez y cada ocurrencia tiene tres ápices distintos. Las demás intersecciones se evitan eligiendo vértices nuevos.

Sea \(A\) la unión de todos los conjuntos \(Y\), todos los ángulos y todas las bisagras de emparejamiento. Sus dos descomposiciones se describen en la siguiente tabla.

{\begin{longtable}[]{@{}
  >{\raggedright\arraybackslash}p{(\linewidth - 2\tabcolsep) * \real{0.5000}}
  >{\raggedright\arraybackslash}p{(\linewidth - 2\tabcolsep) * \real{0.5000}}@{}}
\toprule\noalign{}
\begin{minipage}[b]{\linewidth}\raggedright
Descomposición de \(A\)
\end{minipage} & \begin{minipage}[b]{\linewidth}\raggedright
Descomposición de \(A\cup P\)
\end{minipage} \\
\midrule\noalign{}
\endhead
\bottomrule\noalign{}
\endlastfoot
Las ocurrencias positivas usan (6.6); las negativas usan \(Y\). & Las ocurrencias negativas usan (6.6); las positivas usan \(Y\). \\
Las bisagras de emparejamiento usan sus ángulos negativos. & Las bisagras de emparejamiento usan sus ángulos positivos emparejados. \\
Cada ángulo positivo ya se usó con su ocurrencia. & Cada ángulo positivo sin pareja completa su arista de \(P\). \\
\end{longtable}
}

Cada ángulo se usa exactamente una vez en cada columna. Esto prueba ambas afirmaciones de descomposición.

Hay a lo sumo \(3K\) raíces del certificado, \(K\) bisagras de ocurrencia y \(3K\) bisagras de emparejamiento. Cada bisagra tiene a lo sumo \(C_0\) vértices de soporte, incluidas sus raíces. Contando también \(Z\), el conjunto usado tiene tamaño a lo sumo \(z+3K+4KC_0\). La unidad adicional de \(B\) cubre la elección en curso. La primera desigualdad de (6.3) permite elegir cada vecino común nuevo, y el conteo de bisagras del comienzo de la prueba permite elegir cada bisagra nueva.

Cada ocurrencia aporta a lo sumo \(m+2\) aristas, y cada bisagra de emparejamiento, a lo sumo \(m\). Por tanto, \[|A|\le K(m+2)+3Km\le4K(m+2).\] Finalmente, cada arista añadida tiene un extremo fuera de \(Z\). Para un dispositivo de ocurrencia esto se sigue de la elección de vértices nuevos y de la exclusión de las aristas entre raíces del certificado. En una bisagra de emparejamiento, el único par de raíces que podría estar enteramente en \(Z\) es su base excluida. Esto prueba (6.4).

El mismo emparejamiento de posiciones tiene una versión tipada: todos los dispositivos de ocurrencia y todas las descomposiciones de bisagras deben utilizar tipos triangulares permitidos. Esa es la versión utilizada en la construcción semidensa separada del Teorema C. Para una arista de demanda interna de \(F_0\), se elige su vecino común auxiliar en la parte opuesta. Para una arista de demanda cruzada se utiliza una etiqueta de completación permitida por \(F_0\). Las sustituciones de bisagras son las configuraciones mixtas enumeradas en C.4. La colocación privada de C.5a y C.6 utiliza estas elecciones tipadas, no el Lema SC de triángulos ordinarios sin modificaciones.

Para las bisagras de ocho vértices usadas en Paper I, las cotas son \[B=z+35K+1,\qquad 2^{36}B<\beta n,\qquad |A|\le152K,
\tag{6.7}\] junto con \(B<\alpha n\). Paper I proporciona el certificado firmado y las estimaciones de vecinos comunes y bisagras. El Lema SC realiza la conversión. Construye un absorbedor para un \(P\) especificado; un absorbedor universal para una colección de posibles residuos requiere una construcción simultánea adicional. Ese cuantificador no forma parte de este lema.

\subsection{7. Alcance y obstrucciones}\label{alcance-y-obstrucciones}

\subsubsection{7.1. La divisibilidad agregada no siempre da un certificado físico}\label{la-divisibilidad-agregada-no-siempre-da-un-certificado-fuxedsico}

Para \(r\ge3\), ponemos \[a_\pm=2r^2\pm2r+1,\qquad b=2r^2+3,
\qquad G_r=(K_{a_+}\sqcup K_{a_-})\vee(K_b\sqcup K_b).
\tag{7.1}\] Etiquetamos con cero las dos cliques de la izquierda y con uno las dos de la derecha. Consideramos \(\mathcal F_0=\{001,011\}\). Este grafo es integralmente divisible para \(\mathcal F_0\), pero no admite una descomposición en \(\mathcal F_0\).

En efecto, los tamaños de las cuatro cliques son impares, por lo que todos los grados de su unión completa entre lados son pares. Escribamos \(u=2r^2\). Los conteos de aristas internas y cruzadas son \[e_{00}=u^2+3u,\qquad e_{11}=u^2+5u+6,
\qquad e_{01}=(2u+2)(2u+6)=2(e_{00}+e_{11}).
\tag{7.2}\] Así se cumple (5.1). Sea \(W\) el conjunto de aristas cruzadas entre la clique de tamaño \(a_+\) y una clique fija de tamaño \(b\). Todo triángulo permitido contiene cero o dos aristas de \(W\), pues su arista interna está dentro de una de las cuatro cliques. Pero \(|W|=a_+b\) es impar. Por tanto, el funcional de paridad \[x\longmapsto\sum_{e\in W}x_e\pmod2
\tag{7.3}\] se anula en cada generador de triángulo físico y no en \(\mathbf1_{E(G_r)}\). El indicador de aristas queda fuera del retículo físico (2.3), y la descomposición es imposible.

Esto prueba la distinción precisa que necesitamos aquí. No usamos este cálculo, por sí solo, para afirmar que el ejemplo satisface todas las hipótesis fraccionales o de bisagras de otro teorema. Esas afirmaciones más fuertes requieren sus propias verificaciones.

\subsubsection{7.2. Conteos de bisagras y raíces fijas}\label{conteos-de-bisagras-y-rauxedces-fijas}

La normalización de la sección 6 procede del número de vértices libres. Para dos triángulos fijos que comparten una arista, toda bisagra física de orden \(i\) contiene sus cuatro raíces. El número de esos conjuntos de aristas está acotado por una constante que depende de \(i\) multiplicada por \(n^{i-4}\). En particular, una cota inferior de orden \(n^{i-2}\) para \(i\) fijo no puede cumplirse para todo \(n\) suficientemente grande cuando existe tal par de triángulos.

Una discrepancia de normalización y un contraejemplo a un enunciado corregido son cuestiones distintas. Ninguna determina el contenido del manuscrito que figura como {[}17{]} en {[}2{]}. Nuestras aplicaciones usan la normalización y las hipótesis explícitas que hemos enunciado. Paper I, Apéndice A.2, utiliza bisagras de ocho vértices con oferta proporcional a \(n^4\), exactamente el caso \(i=8\) de la sección 6.

\subsubsection{7.3. Familias más allá de los enlaces de tres pasos}\label{familias-muxe1s-alluxe1-de-los-enlaces-de-tres-pasos}

Los enlaces de tres pasos son una condición suficiente para la construcción del Teorema A. Que falle esa condición no demuestra que un grafo carezca de descomposición. Por ejemplo, para la familia formada por el único tipo tricolor \(\{012\}\), un enlace raíz es bipartito y no hay un paseo cerrado impar en un vértice activo del enlace. Tratar esa familia con este método requiere otra construcción terminal.

El desarrollo más amplio utiliza división de clases de etiquetas (ruta PAR) y correcciones equilibradas que conservan información de carga y torsión (ruta BAL). Son teoremas de construcción, no una definición que supone la descomponibilidad. Estas construcciones cubren 748 representantes y la construcción del núcleo con leyes de conservación proporciona los diecisiete restantes. El censo exhaustivo y el transporte a familias admisibles arbitrarias están ahora demostrados en Lean. La sección 9 enuncia el teorema con hipótesis intrínsecas y explica qué establece la verificación finita. El censo no es necesario para el Corolario 4.2 ni el Teorema C.

El Apéndice D demuestra la reducción por división de vértices. El Apéndice E desarrolla las construcciones por componentes: paridad y equilibrio locales, la carga global adicional cuando existe y la corrección de torsión en el orden adecuado. El Apéndice F construye el catálogo abstracto fijo utilizado en la sección 8. Los certificados finitos de los componentes siguen siendo parte de los datos de la clasificación; su ubicación y la cobertura comprobada se registran en {[}8{]}.

\subsubsection{7.4. Los pesos positivos por tipo son una hipótesis separada}\label{los-pesos-positivos-por-tipo-son-una-hipuxf3tesis-separada}

Consideremos \(\mathcal F=\{001,012\}\) sobre cuatro etiquetas con proporciones de clases iguales. La ecuación (4.8) para el par admisible \(00\) obliga a \(w_{001}=4\), y la ecuación para \(02\) obliga a \(w_{012}=4\). La ecuación para \(01\) exigiría entonces \(w_{001}+w_{012}=4\), una contradicción. Ni siquiera las ecuaciones finitas de equilibrio tienen siempre solución. Este ejemplo explica la hipótesis sobre pesos del Teorema B y el Corolario 4.1; no es un contraejemplo al Teorema D, pues no utiliza la etiqueta tres.

\subsection{8. Realización física del núcleo con leyes de conservación}\label{realizaciuxf3n-fuxedsica-del-nuxfacleo-con-leyes-de-conservaciuxf3n}

El Apéndice F construye el catálogo abstracto utilizado a continuación. Su constante absoluta es \(C=C_m\) de (F.10); las cotas (F.12)--(F.16) proporcionan las estimaciones del grafo de entrada y de carga del catálogo completo necesarias para la realización física.

\subsubsection{Una reserva física fija}\label{una-reserva-fuxedsica-fija}

Para un conjunto de arcos \(D\), escribimos \(d_D(v)=d_D^+(v)+d_D^-(v)\). Un conjunto es balanceado si sus grados de entrada y salida coinciden en cada vértice. Para una colección finita \(\mathcal P\) de conjuntos de arcos, su carga de catálogo es

\[\ell_{\mathcal P}(v)=\sum_{P\in\mathcal P}d_P(v).\]

Se cuentan todas las piezas del catálogo, no sólo las elegidas para un residuo.

Sea \(C\) la constante absoluta de la construcción del catálogo abstracto y sea

\[q=4096(C+1)\cdot21.\]

\textbf{Lema de realización física.} Sea \(G\) un grafo de \(n\) vértices con una partición \(V_0,V_1,V_2\). Supongamos que cada parte tiene al menos \(\omega n\) vértices y que todo conjunto de hasta dos vértices cuyas etiquetas difieren de \(i\) tiene al menos \(\omega n\) vecinos comunes en \(V_i\). Sea \(X\) un conjunto de arcos apoyado en \(G\), orientado cíclicamente de \(V_i\) a \(V_{i+1}\), con índices módulo tres. Supongamos

\[0<t\le1,\qquad 1\le t^2n,\qquad32qt\le\omega,\]

\[|X|\le t^2n^2,\qquad d_X(v)\le2t^2n\quad(v\in V(G)).\]

Existen una reserva \(R\) orientada cíclicamente y apoyada en \(G\), y un catálogo finito \(\mathcal P\), elegidos antes de cualquier \(L\subseteq X\) balanceado, tales que

\[R\cap X=\varnothing,\qquad d_R(v)\le3qtn,\]

\[P\subseteq R\cup X,\quad |P|\le80,\quad P\text{ balanceado}
\qquad(P\in\mathcal P),\]

\[\begin{gathered}
|\mathcal P|\le(49C+16)|X|,\\
\ell_{\mathcal P}(v)\le3qtn.
\end{gathered}\]

Para todo \(L\subseteq X\) balanceado, alguna subcolección de \(\mathcal P\) particiona \(R\cup L\) en conjuntos disjuntos en arcos.

La cota corresponde al grado dirigido total, no al máximo de los grados de entrada y salida por separado. Como la orientación de \(X\) es cíclica, su grado total coincide con el del grafo no dirigido subyacente.

\emph{Demostración.} Aplicamos la construcción del catálogo abstracto a \(X\). Obtenemos un conjunto de puertos, aplicaciones prescritas de los puertos originales a sus raíces físicas, una reserva auxiliar \(R_0\) y un catálogo \(\mathcal P_0\). Escribimos \(I_0\) para la copia levantada de la entrada original. Cada pieza está contenida en \(R_0\cup I_0\) y tiene a lo sumo veinte arcos; \(R_0\) e \(I_0\) son disjuntos; y, para cada \(L\subseteq X\) balanceado, una subcolección particiona \(R_0\) unido a la copia levantada de \(L\). Las piezas son balanceadas o son canales no utilizados cuyos dos puertos extremos tienen la misma raíz física prescrita. Estos canales no tienen por qué ser balanceados en el conjunto de puertos.

Para un arco de catálogo \(e\), sea \(w(e)\) el número de piezas que lo contienen. Para un puerto \(a\), definimos

\[s(a)=\sum_{e\ni a}w(e).\]

La construcción abstracta da \(|\mathcal P_0|\le(49C+16)|X|\) y \(\sum_e w(e)\le(980C+48)|X|\). La carga de catálogo de un puerto antiguo es a lo sumo \(1257\); la de un vértice nuevo, a lo sumo \(7C\sqrt{|X|}\). La carga total ya fijada en una raíz física \(v\) es a lo sumo \(1257d_X(v)\). Por tanto, tomando

\[T=W=qtn,\qquad M=qt^2n^2,\]

tenemos \(s(a)\le W\), \(\sum_a s(a)\le M\), y carga fijada a lo sumo \(T\) en cada vértice físico. Usamos aquí la elección de \(q\), \(\sqrt{|X|}\le tn\), \(t\le1\) y \(1\le t^2n\).

Elegimos las imágenes de los puertos restantes, una a una, en sus partes prescritas. Las de los puertos fijados no se modifican. Para un puerto libre, elegimos un vértice con carga actual menor que \(T\). A lo sumo \(M/T\) vértices pueden tener carga al menos \(T\); como \(M<\omega nT\), existe un candidato en la parte requerida. Asignar el puerto añade a lo sumo \(W\). La aplicación resultante \(f\) tiene carga a lo sumo \(T+W\) en cada vértice físico. No es necesario que sea inyectiva.

Conservamos como arcos físicos originales los arcos del catálogo que pertenecen a \(I_0\). Para cada arco restante \(e=ab\), elegimos un camino

\[f(a),x_e,y_e,z_e,f(b)\]

de longitud cuatro. Si \(a\) tiene etiqueta \(i\), sus etiquetas sucesivas son \(i,i+1,i+2,i,i+1\). Los caminos deben evitar todas las aristas subyacentes de \(X\) y todos los caminos ya elegidos. Cobramos cada camino con su multiplicidad real \(w(e)\) en el catálogo fijo. Las cargas de los extremos se asignan antes de elegir los vértices interiores.

El lema de selección ponderada de caminos se aplica con grado prohibido \(d=2t^2n\), pues

\[\frac{4M}{T}+4+2(d+T+2W)<\omega n.\]

En efecto, el lado izquierdo es a lo sumo \((12+6q)tn<32qtn\le\omega n\). El mismo lema acota la carga conjunta de los arcos conservados y de todos los caminos ponderados por \(T+2W=3qtn\). Es necesario contar las ocurrencias en el catálogo completo: acotar sólo la unión de los caminos no controlaría la carga de catálogo.

Reemplazamos cada arco de catálogo no conservado por su camino, usando esta misma elección de \(f\) y de caminos para todos los residuos. La frontera firmada en vértices del reemplazo de \(ab\) es \(e_{f(b)}-e_{f(a)}\). Una pieza abstracta balanceada sigue siendo balanceada después del reemplazo, aunque \(f\) identifique puertos. Un canal no utilizado también pasa a ser balanceado: sus dos extremos tienen la misma raíz física. Cada pieza resultante tiene a lo sumo \(4\cdot20=80\) arcos. Todos tienen la orientación cíclica requerida.

Sea \(A\) la unión de los arcos del catálogo abstracto y sea \(I=A\cap I_0\). La partición del residuo vacío muestra que \(R_0\subseteq A\). Como todas las piezas están apoyadas en la unión disjunta \(R_0\cup I_0\), resulta \(A\setminus I=R_0\). Definimos \(R\) como la unión de los caminos que reemplazan \(R_0\), y \(\mathcal P\) como el conjunto de las piezas reemplazadas. Así, \(R\) está apoyada en \(G\) y evita \(X\).

Los bloques de reemplazo de arcos abstractos distintos son disjuntos en aristas. Esto se sigue de la selección de caminos, de que evitan \(X\) y de que los arcos de entrada conservados tienen imágenes físicas distintas. El reemplazo transporta entonces cada partición abstracta a una partición física disjunta en arcos. Las raíces prescritas identifican la copia levantada de cada \(L\) balanceado exactamente con \(L\); su unión es, por tanto, \(R\cup L\). Ninguna elección nueva depende de \(L\).

Finalmente, la estimación ponderada acota \(\ell_{\mathcal P}(v)\) por \(3qtn\); identificar piezas iguales no puede aumentar esa suma. El reemplazo tampoco aumenta el número de piezas. Aplicando la partición al residuo vacío, el grado \(d_R(v)\) queda acotado por la misma carga de catálogo. \(\square\)

\subsubsection{Regreso a las diecisiete familias}\label{regreso-a-las-diecisiete-familias}

Cada una de las diecisiete familias conservadas tiene una división explícita de vértices cuyas componentes son un núcleo conservado y las componentes no nucleares certificadas por separado. En el núcleo, los tipos triangulares de la familia dividida son exactamente las ternas con tres colores distintos. La oferta de vecinos comunes se toma de la propia familia de la componente del núcleo. No se infiere aumentando esa familia.

El lema de realización física, seguido de la construcción del absorbedor terminal para piezas acotadas, da un certificado universal de absorbedores de tres colores con un exponente absoluto \(L_C\) en las condiciones de escala siguientes. Pasando al núcleo físico inducido y transportando de vuelta su absorbedor obtenemos el certificado de la componente del núcleo. Para las componentes no nucleares, que son finitas, tomamos una escala común \(2^{-L_N}\) de sus certificados existentes. Aumentar el exponente de escala sólo refuerza la hipótesis de dispersión. Podemos usar, por tanto,

\[L=\max\{L_C,L_N\}\]

en todas las componentes ocupadas. El transporte por división de vértices produce un certificado universal de la familia a exponente \(L+11\). Esta elección precede al grafo anfitrión. El teorema de descomposición de anfitriones típicos se aplica entonces a cada una de las diecisiete familias sin dejar como hipótesis un productor latino no demostrado.

Esto demuestra las existencias para esos representantes. La exhaustividad requiere además el censo finito de la Sección 9; no se deduce de la construcción del absorbedor.

\textbf{Convención de escala y pérdida de once.} Un certificado con exponente \(L\) proporciona el absorbedor universal cuando se cumplen las hipótesis de disponibilidad protegida y \[\begin{gathered}
0<\omega\le1,\quad 0<\delta,\quad 2^{63}\delta\le\omega,\\
2^{24}\rho\le\left(\frac{\omega^2\delta}{2^L}\right)^2,
\quad 1\le\rho n.
\end{gathered}\] Así, \(L\) interviene en la cota superior de densidad de la reserva; no significa que \(\delta=2^{-L}\). El certificado utiliza disponibilidad protegida de orden \(2^{21}\) con parámetro \(\omega n\) y un anfitrión admisible. Su conclusión es un absorbedor universal de grado máximo a lo sumo \(2\delta n\) para una reserva de entrada de grado máximo a lo sumo \(2\rho n\).

El cambio de escala al dividir etiquetas es el siguiente. Sean \(1\le t'\le8\), \(n\le n'\le t'n\) y \(x=n/n'\). Utilizamos \[\omega'=\omega x,\qquad \delta'=\delta x/t',\qquad \rho'=\rho x.\] Se conservan las condiciones de grado y orden. Escribamos \(P=(\omega^2\delta/2^L)^2\). La condición original del certificado con exponente \(L+11\) da \(2^{24}\rho'\le Px/2^{22}\), mientras que la condición buscada con exponente \(L\) tiene lado derecho \(Px^6/(t')^2\). Basta que \[(t')^2\le2^{22}x^5.\] En efecto, \(x\ge1/8\) y \((t')^2\le64\), por lo que el lado derecho es al menos \(2^{22}/8^5=128\). Además, \(\rho'n'=\rho n\) y \(\delta'n'=\delta n/t'\). La última igualdad asigna a cada componente su parte del presupuesto de grado antes de proyectar. Esto demuestra la pérdida de exponente indicada; el documento complementario identifica el lema de cambio de escala y su uso en el ensamblaje de familias conservadas.

\subsection{9. Todas las familias admisibles de cuatro etiquetas}\label{todas-las-familias-admisibles-de-cuatro-etiquetas}

Llamamos \emph{admisible} a una familia sobre cuatro etiquetas si cada elemento es un multiconjunto de tamaño tres, toda etiqueta aparece en algún elemento y cada etiqueta tiene un compañero prohibido como en (3.2). Es una condición sobre los tipos: no supone una descomposición ni remite a una lista de construcciones satisfactorias.

\subsubsection{Teorema D (familias de cuatro etiquetas; clasificación finita asistida por computador)}\label{teorema-d-familias-de-cuatro-etiquetas-clasificaciuxf3n-finita-asistida-por-computador}

Toda familia admisible de cuatro etiquetas tiene la propiedad de descomposición para grafos típicos (2.7). Existen exactamente 765 de estas familias salvo permutación de las cuatro etiquetas.

La clasificación utiliza tablas explícitas de certificados finitos comprobadas en Lean; las construcciones universales de grafos son los argumentos separados descritos a continuación.

La primera afirmación conserva todas las hipótesis sobre el grafo de (2.7), incluidos los pesos regulares, la extensión y las dos condiciones de retículo entero. No afirma que todo grafo casi equilibrado admita pesos positivos por tipo. El Teorema B se aplica a cada una de estas familias cuando se cumple su hipótesis adicional sobre los pesos.

\emph{Reducción finita y transporte.} Enumeramos las veinte ternas no decrecientes sobre cuatro etiquetas. Asociar a una familia el conjunto de índices de sus ternas da una codificación invertible. Una permutación de las cuatro etiquetas actúa en esa codificación mediante la correspondiente permutación de ternas. Codificar y decodificar, o reetiquetar y aplicar la permutación inversa, son operaciones inversas.

Para una familia admisible, elegimos un compañero prohibido para cada etiqueta. Al excluir los pares correspondientes, su codificación queda contenida en una de veintiséis cajas finitas de tipos permitidos. Cada caja contiene como máximo doce tipos. La comprobación finita de cobertura considera todos los subconjuntos de cada caja que utilizan las cuatro etiquetas. Para cada uno proporciona una de las veinticuatro permutaciones reales de etiquetas que lo lleva a un representante registrado. El certificado aporta, por tanto, una permutación, no sólo igualdad de invariantes numéricos.

La lista contiene 765 máscaras distintas. Se comprueba que cada una es admisible y mínima en su órbita de permutaciones. Si dos máscaras de la lista representan la misma órbita, la minimalidad da ambas desigualdades entre ellas y, por tanto, su igualdad. La cobertura y esta unicidad dan el número de órbitas. Las afirmaciones finitas, las identidades de codificación y su ensamblaje están demostrados en Lean; la búsqueda Python que produjo los certificados candidatos no es una premisa del resultado.

Para cada representante, la biblioteca de construcciones proporciona la propiedad de descomposición para grafos típicos. Esto incluye los diecisiete casos con leyes de conservación de la sección 8. Para pasar del representante a la familia original, reetiquetamos las partes del mismo grafo mediante la permutación inversa. Los pares admisibles, tipos de triángulos, vectores de aristas por tipo y perfiles de esquinas se transportan biyectivamente. Los dos retículos enteros pasan a los retículos correspondientes; los pesos regulares siguen asignados a los mismos triángulos físicos, y los conjuntos de extensión sólo cambian de etiqueta. Aplicamos el teorema del representante y deshacemos el reetiquetado. Así obtenemos la propiedad para la familia original, sin una hipótesis de pertenencia al catálogo.

Los Apéndices A--F proporcionan las construcciones. La ruta PAR cubre 271 representantes; BAL cubre 357, con 301 casos de equilibrio sin carga adicional, 52 con carga y cuatro con torsión. Junto con los 120 casos anteriores y los diecisiete casos con leyes de conservación, son los 765 representantes del certificado finito. Los conteos no prueban por sí solos la cobertura: son esenciales los datos explícitos de componentes y los certificados de permutaciones. Las construcciones universales explican cómo esos datos implican descomposición; la verificación finita establece que cada representante tiene los datos requeridos.

\subsection{Formalización y reproducibilidad}\label{formalizaciuxf3n-y-reproducibilidad}

El desarrollo se ha comprobado en Lean 4 con Mathlib {[}6,7{]}. La campaña alineada compiló la clausura completa de 1.510 módulos de tres entradas explícitas de Paper II y registró quince consultas finales de axiomas que sólo usan extensionalidad proposicional, elección clásica y solidez de los cocientes. Su alcance incluye los certificados de existencia de las diecisiete familias con leyes de conservación, el censo y el transporte a toda familia admisible de cuatro etiquetas. Utiliza el mismo puente de cuatro colores que Paper I. El archivo seleccionado corrige las cabeceras de autoría después de la compilación conservando cada byte del cuerpo de las pruebas; su mapa reversible de fuentes vincula los recibos históricos con los archivos corregidos. El documento complementario {[}8{]}, preprint 1.0, registra las huellas del archivo, la correspondencia de declaraciones, las ubicaciones de certificados finitos, la evidencia de la campaña y los límites de reproducción. Esta integración editorial no ejecuta una compilación Lean ni afirma una reproducción independiente. El documento complementario distingue las consultas finales de las entradas y la evidencia de axiomas a nivel de módulo, incluido el corolario casi completo.

\subsection{Agradecimientos}\label{agradecimientos}

El autor agradece profundamente a su esposa María Paz y a sus hijos Lucas, Juan Cristóbal, Francisca, Raimundo y Benjamín su amor, paciencia y apoyo. Agradece también a Delcourt y Postle el trabajo matemático que motivó este desarrollo, y a las comunidades de Lean y Mathlib la infraestructura utilizada en la formalización.

\subsection{Uso de inteligencia artificial y herramientas computacionales}\label{uso-de-inteligencia-artificial-y-herramientas-computacionales}

Se utilizaron Claude, de Anthropic, y ChatGPT/Codex, de OpenAI, para explorar y revisar argumentos y preparar el manuscrito. Aristotle, de Harmonic, contribuyó a explorar enfoques, buscar contraejemplos, desarrollar demostraciones y formalizarlas y revisarlas en Lean. Su participación no se limitó a traducir pruebas terminadas.

Certo {[}9{]} y Jacobian {[}10{]} se utilizaron para explorar casos finitos, comprobar identidades y evaluar construcciones candidatas. Las comprobaciones finitas no se trataron como demostraciones de enunciados universales. La búsqueda del censo produjo certificados finitos. Lean verificó después la cobertura del universo admisible, la unicidad salvo permutaciones reales y la conexión con los teoremas de construcción; esto es más que compilar una lista de ejemplos.

El autor es responsable de los argumentos, las citas, el código y la presentación. Ningún sistema de inteligencia artificial figura como autor.


\clearpage
\subsection{Referencias}\label{referencias}

\begingroup\small\raggedright
\setlength{\parskip}{3pt}

{[}1{]} J. P. Traverso Gianini, \emph{Triangle decompositions I: Formalizing the Nash--Williams theorem in Lean}, preprint 1.0 (2026). https://doi.org/10.5281/zenodo.23428173. \href{https://github.com/jtraverso/nash-williams-triangle-decompositions/blob/main/preprints/PAPER_I/01_manuscript/PAPER_I_preprint_v1.0_es.md}{Documento}.

{[}2{]} M. Delcourt y L. Postle, \emph{A Proof of Nash-Williams' Conjecture}, arXiv:2606.11178v1 (2026). https://arxiv.org/abs/2606.11178v1.

{[}3{]} M. Delcourt y L. Postle, \emph{Refined absorption for partitioned hypergraphs and beyond}, manuscrito, 2026. Descripción bibliográfica según {[}2, referencia 17{]}; no se ha localizado una versión pública durante la preparación de esta serie.

{[}4{]} M. Delcourt y L. Postle, \emph{Refined Absorption: A New Proof of the Existence Conjecture}, arXiv:2402.17855 (2024). https://arxiv.org/abs/2402.17855.

{[}5{]} P. Keevash, \emph{The existence of designs II}, arXiv:1802.05900 (2018). https://arxiv.org/abs/1802.05900.

{[}6{]} L. de Moura y S. Ullrich, \emph{The Lean 4 Theorem Prover and Programming Language}, CADE 28 (2021), 625--635. https://doi.org/10.1007/978-3-030-79876-5\_37.

{[}7{]} The mathlib Community, \emph{The Lean Mathematical Library}, CPP 2020, 367--381. https://doi.org/10.1145/3372885.3373824.

{[}8{]} J. P. Traverso Gianini, \emph{Formalization and reproducibility companion to Triangle decompositions II}, preprint 1.0 (2026), correspondiente al freeze \texttt{PAPER\_\allowbreak{}II\_\allowbreak{}LEA\allowbreak{}N\_\allowbreak{}20261020\_\allowbreak{}A\allowbreak{}LIGNED\_\allowbreak{}v4}. El complemento identifica el archivo inmutable de las fuentes y su huella SHA-256. \href{https://github.com/jtraverso/nash-williams-triangle-decompositions/blob/main/preprints/PAPER_II/05_formalization/FORMALIZATION_COMPANION_preprint_v1.0_es.md}{Documento}.

{[}9{]} J. P. Traverso Gianini, \emph{Certo}, repositorio de software. https://github.com/jtraverso/certo-math.

{[}10{]} \emph{Jacobian}, repositorio de software. https://github.com/morluto/jacobian.

{[}11{]} B. Barber, D. Kühn, A. Lo, D. Osthus y A. Taylor, \emph{Clique decompositions of multipartite graphs and completion of Latin squares}, Journal of Combinatorial Theory, Series A 151 (2017), 146--201. https://doi.org/10.1016/j.jcta.2017.04.005.

{[}12{]} R. Montgomery, \emph{Fractional clique decompositions of dense partite graphs}, Combinatorics, Probability and Computing 26 (2017), 911--943. https://doi.org/10.1017/S0963548317000165. https://arxiv.org/abs/1603.01039.

{[}13{]} B. Barber, D. Kühn, A. Lo y D. Osthus, \emph{Edge-decompositions of graphs with high minimum degree}, Advances in Mathematics 288 (2016), 337--385. https://doi.org/10.1016/j.aim.2015.09.032.

{[}14{]} S. Glock, D. Kühn, A. Lo y D. Osthus, \emph{The existence of designs via iterative absorption: hypergraph F-designs for arbitrary F}, arXiv:1611.06827v3 (2020). https://arxiv.org/abs/1611.06827v3.

{[}15{]} P. Keevash, \emph{The existence of designs}, arXiv:1401.3665v4 (2024; primera versión 2014). https://arxiv.org/abs/1401.3665v4.

{[}16{]} V. Rödl, \emph{On a packing and covering problem}, European Journal of Combinatorics 6 (1985), 69--78. https://doi.org/10.1016/S0195-6698(85)80023-8.

{[}17{]} P. Frankl y V. Rödl, \emph{Near perfect coverings in graphs and hypergraphs}, European Journal of Combinatorics 6 (1985), 317--326. https://doi.org/10.1016/S0195-6698(85)80045-7.

{[}18{]} N. Pippenger y J. Spencer, \emph{Asymptotic behavior of the chromatic index for hypergraphs}, Journal of Combinatorial Theory, Series A 51 (1989), 24--42. https://doi.org/10.1016/0097-3165(89)90074-5.

{[}19{]} P. E. Haxell y V. Rödl, \emph{Integer and fractional packings in dense graphs}, Combinatorica 21 (2001), 13--38. https://doi.org/10.1007/s004930170003.

{[}20{]} S. Ehard, S. Glock y F. Joos, \emph{Pseudorandom hypergraph matchings}, Combinatorics, Probability and Computing 29 (2020), 868--885. https://doi.org/10.1017/S0963548320000280.

{[}21{]} S. Boucheron, G. Lugosi y P. Massart, \emph{Concentration Inequalities: A Nonasymptotic Theory of Independence}, Oxford University Press (2013), Capítulo 2. https://doi.org/10.1093/acprof:oso/9780199535255.001.0001.

\endgroup
\clearpage

\Needspace{12\baselineskip}\subsection{Apéndice A. Álgebra, certificados enteros y absorbedores finitos}\label{apuxe9ndice-a.-uxe1lgebra-certificados-enteros-y-absorbedores-finitos}

En estos apéndices, «compatible entero» y «divisible integralmente para \(F\)» designan las mismas condiciones de retículo local y global de la Sección 2.

\subsubsection{A.1. Convenciones y compatibilidad}\label{a.1.-convenciones-y-compatibilidad}

Sea \(F\) una familia de multiconjuntos de tamaño tres sobre un conjunto finito de etiquetas \(I\). Los vértices se particionan en clases \(V_i\), y \(n_i=|V_i|\), \(n=\sum_i n_i\). Un par de etiquetas está permitido si aparece en un miembro de \(F\), con multiplicidad cuando las dos etiquetas coinciden. El grafo completo permitido contiene exactamente las aristas sin lazos cuyos pares de etiquetas están permitidos. Un triángulo permitido es un conjunto no ordenado de tres vértices distintos cuyo multiconjunto de etiquetas pertenece a \(F\). Su vector de aristas se denota por \([abc]\). Para una cadena de triángulos \(z\), escribimos \(\partial z\) para su vector de aristas y \(\|z\|_1\) para la suma de los valores absolutos de sus coeficientes.

Un vector de aristas \(r\) es simétrico y tiene diagonal nula. Su perfil local y sus conteos globales \textbf{ordenados} son \[\begin{gathered}
d_r(v)_j=\sum_{w\in V_j}r_{vw},\\
T_{ij}(r)=\sum_{v\in V_i}d_r(v)_j.
\end{gathered}\] Así, una arista interna contribuye dos veces a \(T_{ii}\). Fuera de la diagonal, \(T_{ij}=T_{ji}\), y cada una de esas dos coordenadas cuenta la arista una vez. Todas las columnas globales que siguen utilizan esta convención. En particular, \(\sum_{v\in V_i}d_r(v)=T_{i\bullet}(r)\).

Sea \(g_\tau\) el vector de conteos globales de un triángulo de tipo \(\tau\). Para \(i\in\tau\), sea \[b_{\tau,i}(j)=\operatorname{mult}_\tau(j)-\mathbf1_{j=i}\] su perfil en una esquina de etiqueta \(i\). Decimos que \(r\) es compatible sobre los reales si \(T(r)\) pertenece al espacio real generado por los \(g_\tau\) y cada \(d_r(v)\) pertenece al espacio real generado por los \(b_{\tau,p(v)}\). La compatibilidad entera exige las mismas condiciones para las combinaciones lineales enteras. Son dos requisitos distintos.

Para un vector de aristas, \(\|r\|_1\) cuenta aristas no ordenadas. Por tanto, \[\sum_v\|d_r(v)\|_1\le 2\|r\|_1,
\qquad \|\partial z\|_1\le3\|z\|_1.
\tag{A.1.1}\]

\subsubsection{A.2. Una inversa real que incluye los perfiles locales}\label{a.2.-una-inversa-real-que-incluye-los-perfiles-locales}

\textbf{Proposición A.2.1.} Supongamos que \(|I|=4\), \(0<\gamma\le1\) y \(n_i\ge\gamma n\ge8\) para cada etiqueta. Existe una constante \(C_F>0\), que depende únicamente de \(F\), tal que todo vector \(r\) compatible sobre los reales, soportado en el grafo completo permitido y con \(\|r\|_\infty\le\epsilon\), tiene una preimagen triangular \(z\) que satisface \[\begin{gathered}
\partial z=r,\\
\|z\|_\infty\le \frac{C_F\epsilon}{\gamma^6n}.
\end{gathered}
\tag{A.2.1}\]

\textbf{Demostración.} Una aplicación lineal entre espacios reales de dimensión finita tiene una inversa por la derecha sobre su imagen: elegimos una base de la imagen, una preimagen de cada vector de esa base y extendemos linealmente. Su norma de operador es finita. Aplicamos esta observación a la matriz de columnas \(g_\tau\) y a cada matriz de columnas \(b_{\tau,i}\). Elegimos constantes positivas \(K_1,K_2\) que acoten las normas de coordenadas resultantes. Usamos la \emph{misma} inversa local por la derecha en todos los vértices de una clase. Esta elección preservará las sumas.

\paragraph{Corrección global}\label{correcciuxf3n-global}

Como cada coordenada global de \(r\) tiene valor absoluto a lo sumo \(\epsilon n^2\), elegimos coeficientes \(c_\tau\) tales que \[\begin{gathered}
\sum_{\tau\in F}c_\tau g_\tau=T(r),\\
|c_\tau|\le K_1\epsilon n^2.
\end{gathered}\] Pongamos \(m=\gamma n\). El número \(N_\tau\) de triángulos físicos de un tipo dado es al menos \(m^3/64\). En efecto, elegimos una terna ordenada de vértices con las etiquetas prescritas, excluyendo en cada elección los vértices anteriores, y dividimos por un número a lo sumo seis. Para \(m\ge8\), \(m(m-1)(m-2)/6\ge m^3/64\).

Damos a cada triángulo de tipo \(\tau\) peso \(c_\tau/N_\tau\), obteniendo \(z_1\). Entonces \(T(\partial z_1)=T(r)\) y \[\|z_1\|_\infty\le\frac{64K_1\epsilon}{\gamma^3n}.\] Una arista pertenece a lo sumo a \(n\) triángulos. Por tanto, \(r_1=r-\partial z_1\) tiene conteos globales nulos y \[\|r_1\|_\infty\le E_1,
\qquad E_1=\frac{(1+64K_1)\epsilon}{\gamma^3}.
\tag{A.2.2}\] Restar una cadena de triángulos preserva la compatibilidad sobre los reales.

\paragraph{Corrección de perfiles}\label{correcciuxf3n-de-perfiles}

Aplicamos la inversa local por la derecha elegida a cada \(d_{r_1}(v)\). Sus coeficientes \(c_{v,\tau}\), definidos para los tipos que contienen \(i=p(v)\), satisfacen \[\begin{gathered}
\sum_{\tau\ni i}c_{v,\tau}b_{\tau,i}=d_{r_1}(v),\\
|c_{v,\tau}|\le K_2nE_1,
\qquad \sum_{v\in V_i}c_{v,\tau}=0.
\end{gathered}
\tag{A.2.3}\] La última igualdad se deduce por linealidad de \(T(r_1)=0\); elegir las preimágenes de forma independiente no la justificaría.

Registramos el multiplicador exacto de la siguiente construcción. Fijemos \(\tau\) y números \(\phi_\tau(v)\) cuya suma sea cero en cada clase y que sean cero en las clases ausentes de \(\tau\). Damos a un triángulo \(abc\) de tipo \(\tau\) peso \(\phi_\tau(a)+\phi_\tau(b)+\phi_\tau(c)\). En un vértice \(v\) de etiqueta \(i\in\tau\), su perfil es \(\kappa_{\tau,i}\phi_\tau(v)b_{\tau,i}\). Los multiplicadores son los siguientes; escribimos \(\tau=\{i,x,y\}\).

{\begin{longtable}[]{@{}lr@{}}
\toprule\noalign{}
Etiquetas & \(\kappa_{\tau,i}\) \\
\midrule\noalign{}
\endhead
\bottomrule\noalign{}
\endlastfoot
\(i,x,y\) distintas & \(n_xn_y\) \\
\(i=x\ne y\) & \((n_i-2)n_y\) \\
\(i=y\ne x\) & \(n_x(n_i-2)\) \\
\(x=y\ne i\) & \(n_x(n_x-1)/2\) \\
\(i=x=y\) & \((n_i-2)(n_i-3)/2\) \\
\end{longtable}
}

Damos una verificación directa que incluye las etiquetas repetidas. En una raíz de etiqueta \(i\), el grado hacia la etiqueta \(l\) se obtiene sumando sobre las elecciones ordenadas de la otra esquina de etiqueta \(l\) y de la esquina que completa el triángulo, de etiqueta \(z\). Al sustituir por cero cada suma completa de \(\phi\) en una clase, queda el coeficiente \[(n_l-2\mathbf1_{i=l})
(n_z-\mathbf1_{i=z}-\mathbf1_{l=z})
-\mathbf1_{i=z}(n_l-\mathbf1_{i=l})
+\mathbf1_{l=z}\mathbf1_{i=l}.
\tag{A.2.4}\] Sólo aparecen completaciones con \(\{i,l,z\}=\tau\). La sustitución en (A.2.4) da la tabla, dividiendo por dos cuando la propia columna del perfil tiene entrada dos. Esto explica también por qué tratar las tres etiquetas como distintas daría una corrección incorrecta.

Todo multiplicador de la tabla es al menos \(m^2/8\). Pongamos \(\phi_\tau(v)=c_{v,\tau}/\kappa_{\tau,p(v)}\). Los denominadores son constantes dentro de cada clase, por lo que estas funciones siguen sumando cero en cada clase. Sea \(z_2\) la cadena resultante. La ecuación (A.2.3) y la identidad del multiplicador dan \(d_{\partial z_2}(v)=d_{r_1}(v)\) para todo \(v\). Además, \[\begin{gathered}
\|z_2\|_\infty\le\frac{24K_2E_1}{\gamma^2n},\\
\|r_2\|_\infty\le E_2:=\frac{(1+24K_2)E_1}{\gamma^2},
\end{gathered}
\tag{A.2.5}\] donde \(r_2=r_1-\partial z_2\) tiene todos sus perfiles locales nulos.

\paragraph{Corrección de aristas}\label{correcciuxf3n-de-aristas}

Para un par permitido de etiquetas \(ij\), sea \(q_{ij}\) el número de tipos que lo contienen. Para uno de esos tipos \(\tau\), sea \(k\) la etiqueta restante, determinada de forma única, y pongamos \[e_{\tau,ij}=n_k-2\mathbf1_{i=k}-2\mathbf1_{j=k}\ge m-4\ge m/2.\] Definimos una función simétrica sobre las aristas \[h_\tau(u,v)=\frac{r_2(u,v)}{q_{ij}e_{\tau,ij}}
\quad (u\in V_i,v\in V_j),\] y le damos valor cero en los pares ausentes de \(\tau\) y en la diagonal. Cada suma de fila restringida a una clase es cero, pues su denominador es constante en cada bloque de etiquetas. Damos a un triángulo de tipo \(\tau\) peso \(h_\tau(a,b)+h_\tau(a,c)+h_\tau(b,c)\).

Para una arista fija \(uv\), la suma sobre su tercer vértice cuenta inicialmente \(h_\tau(u,v)\) exactamente \(n_k-\mathbf1_{i=k}-\mathbf1_{j=k}\) veces. Las otras dos sumas son sumas nulas sobre clases a las que se han quitado los extremos prohibidos; restan otras \(\mathbf1_{i=k}+\mathbf1_{j=k}\) copias. Así, este tipo contribuye \(e_{\tau,ij}h_\tau(u,v)=r_2(u,v)/q_{ij}\). Sumando sobre los \(q_{ij}\) tipos se obtiene \(\partial z_3=r_2\), con \[\|z_3\|_\infty\le\frac{6E_2}{\gamma n}.\] Finalmente, \(z=z_1+z_2+z_3\) demuestra (A.2.1), con la elección explícita \[C_F=64K_1+24K_2(1+64K_1)
+6(1+24K_2)(1+64K_1).
\tag{A.2.6}\] La sustitución de potencias menores de \(\gamma^{-1}\) por \(\gamma^{-6}\) utiliza \(\gamma\le1\). \(\square\)

\textbf{Corolario A.2.2 (compatibilidad del defecto efectivo).} Si \(H\) es compatible sobre los enteros y \(\varphi\) es cualquier ponderación real de sus triángulos permitidos, entonces \(r=\mathbf1_{E(H)}-\partial\varphi\) es compatible sobre los reales. Si \(B\subseteq G\) tiene una descomposición en triángulos permitidos y \(G\) es compatible sobre los enteros, la misma conclusión vale para \(H=G-B\).

\textbf{Demostración.} Globalmente, restamos la suma de \(\varphi(T)g_{\tau(T)}\) del vector de conteos de \(H\), compatible sobre los enteros. Localmente, en \(v\), restamos \(\sum_{T\ni v}\varphi(T)b_{\tau(T),p(v)}\) de su perfil compatible sobre los enteros. Ambas diferencias pertenecen a los espacios reales requeridos. Eliminar un grafo descomponible resta sumas enteras de las mismas columnas y, por tanto, preserva la compatibilidad entera. \(\square\)

\subsubsection{A.3. Trasladar la inversa real al grafo anfitrión}\label{a.3.-trasladar-la-inversa-real-al-grafo-anfitriuxf3n}

\textbf{Proposición A.3.1 (transferencia con signos).} Supongamos que el grafo anfitrión proporciona al menos \(s>0\) elecciones en cualquier clase compatible prescrita, adyacentes a cualquier conjunto de a lo sumo cuatro vértices prescritos. Sea \(z\) una cadena real en el grafo completo permitido tal que \(r=\partial z\) está soportado en aristas del anfitrión. Supongamos que \[\begin{gathered}
\|r\|_\infty\le\epsilon,\\
\sum_T|z(T)|\le S,\\
\max_v\sum_{T\ni v}|z(T)|\le U.
\end{gathered}\] Entonces existe una cadena \(w\) sobre triángulos efectivos del anfitrión con \(\partial w=r\) y, para cada triángulo \(Q\) del anfitrión, \[|w(Q)|\le\frac{9\epsilon}{s}+
\frac{81U}{s^2}+\frac{108S}{s^3}.
\tag{A.3.1}\]

\textbf{Demostración.} Promediamos dos sustituciones explícitas. Las elecciones aleatorias que siguen sólo sirven para definir coeficientes reales mediante promedios finitos; no se utiliza ningún teorema probabilístico asintótico.

Para cada triángulo virtual \(abc\), elegimos \(A\) con la etiqueta de \(a\) y adyacente a \(b,c\); después, \(B\) con la etiqueta de \(b\) y adyacente a \(a,c,A\); finalmente, \(C\) con la etiqueta de \(c\) y adyacente a \(a,b,A,B\). En cada paso elegimos uniformemente entre los vértices disponibles. Hay al menos \(s\) elecciones. La identidad de vectores de aristas \[[abc]=[abC]+[aBc]+[Abc]-[aBC]-[AbC]-[ABc]+[ABC]
\tag{A.3.2}\] se verifica cancelando aristas. Los últimos cuatro triángulos pertenecen al anfitrión. Los tres primeros constan de una arista antigua y una cuña física de dos aristas. No se necesita que todos los vértices de \(a,b,c,A,B,C\) sean distintos: cada triángulo mostrado tiene tres vértices distintos, como exigen las condiciones de adyacencia.

Fijemos una arista antigua \(ab\). Las apariciones de cuñas sobre ella tienen masas con signo \(u_j\), masa total \(\sum_j u_j=r_{ab}\) y masa absoluta total \(m_{ab}=\sum_{T\supset ab}|z(T)|\). Cada distribución del ápice tiene probabilidades puntuales a lo sumo \(1/s\), también para las elecciones posteriores del experimento anterior. Emparejamos masas positivas y negativas, dividiéndolas si es necesario. La masa no emparejada tiene tamaño \(|r_{ab}|\) y un único signo. Para un ápice positivo \(c\) y otro negativo \(d\) emparejados, elegimos independientemente según sus dos distribuciones; después elegimos \(A\) con la etiqueta de \(a\) y adyacente a \(b,c,d\), y \(B\) con la etiqueta de \(b\) y adyacente a \(a,c,d,A\). La identidad \[[abc]-[abd]=
[aBc]+[Abc]+[ABd]-[aBd]-[Abd]-[ABc]
\tag{A.3.3}\] sustituye el par por seis triángulos del anfitrión. En particular, no exige que \(ab\) sea una arista del anfitrión. Las apariciones no emparejadas ya son triángulos del anfitrión: su masa se anula cuando \(ab\notin E(G)\), porque \(r\) está soportado en \(G\). Las ecuaciones (A.3.2) y (A.3.3) demuestran la afirmación sobre la frontera.

Para completar el argumento, acotamos los coeficientes de este promedio. Fijemos un triángulo \(Q\) del anfitrión y pongamos \(O_Q=\sum_T|z(T)|\,|T\cap Q|\le3U\). Una sucesión adaptativa de \(k\) elecciones, cada una con probabilidades puntuales condicionales a lo sumo \(1/s\), tiene probabilidad a lo sumo \((3/s)^k\) de situar todas sus elecciones en \(Q\). Por tanto, las cuatro caras del anfitrión en (A.3.2) contribuyen en valor absoluto a lo sumo \[S(3/s)^3+O_Q(3/s)^2.\] Para una cuña emparejada sobre \(ab\), las seis caras de (A.3.3) contribuyen a lo sumo \[2\bigl((\mathbf1_{a\in Q}+\mathbf1_{b\in Q})(3/s)^2+(3/s)^3\bigr)\] veces su masa emparejada. El doble de esa masa es a lo sumo \(m_{ab}\). Además, \[\begin{gathered}
\sum_{ab}m_{ab}=3\sum_T|z(T)|\le3S,\\
\sum_{ab}m_{ab}(\mathbf1_{a\in Q}+\mathbf1_{b\in Q})=2O_Q.
\end{gathered}\] Las cuñas no emparejadas sólo pueden contribuir sobre las tres aristas de \(Q\); su contribución es a lo sumo \(3\epsilon(3/s)\). El total es, por tanto, \[3\epsilon(3/s)+3O_Q(3/s)^2+4S(3/s)^3,\] que es (A.3.1). \(\square\)

\textbf{Corolario A.3.2 (una reparación cuantitativa).} En la Proposición A.2.1 tomemos \(\gamma=\omega\), donde \(0<\omega\le1\), y supongamos que el anfitrión proporciona \(s=\omega n/2\) elecciones como en la Proposición A.3.1. Entonces un defecto compatible de norma a lo sumo \(\epsilon\) admite una corrección soportada en el anfitrión y acotada por \[\|w\|_\infty\le
\frac{(18+1188C_F)\epsilon}{\omega^9n}.
\tag{A.3.4}\] En efecto, (A.2.1) da \(S\le C_F\epsilon n^2/\omega^6\) y \(U\le C_F\epsilon n/\omega^6\). La sustitución en (A.3.1), seguida de \(\omega\le1\), da (A.3.4). Así, los pesos inicialmente en \([\omega^2/n,27/(\omega n)]\) siguen siendo positivos y son a lo sumo \(54/(\omega n)\) tras la corrección siempre que \[(18+1188C_F)\epsilon\le\omega^{11}/2.
\tag{A.3.5}\] Sus cargas sobre las aristas son entonces exactamente uno.

Las estimaciones de borrado utilizadas con este corolario son explícitas. Si se elimina un grafo descomponible \(B\) de grado máximo a lo sumo \(\kappa n\), la disponibilidad para cuatro vecinos prescritos pierde a lo sumo \(4\kappa n\) elecciones. Restringir los pesos triangulares anteriores elimina a lo sumo \(2\kappa n\) triángulos en cada arista retenida. Bajo la cota superior de pesos recién enunciada, un error de carga inicial \(\epsilon\) pasa a ser a lo sumo \(\epsilon+54\kappa/\omega\). El Corolario A.2.2 proporciona ambas condiciones de compatibilidad para este nuevo defecto. En particular, la inversa algebraica no se aplica basándose únicamente en los conteos globales.

\subsubsection{A.4. Certificados enteros para piezas acotadas}\label{a.4.-certificados-enteros-para-piezas-acotadas}

\textbf{Proposición A.4.1.} Fijemos \(F\) y un entero \(K\ge0\). Existe un entero \(C=C(F,K)\ge0\) con la siguiente propiedad. En un grafo completo permitido con cinco representantes designados distintos en cada clase, todo conjunto de aristas \(P\) compatible sobre los enteros, con \(|P|\le K\), tiene un certificado triangular entero de coste a lo sumo \(C\) que utiliza sólo los vértices de \(P\) y los representantes.

\textbf{Demostración.} Para una matriz entera fija, sólo hay un número finito de vectores de su imagen entera en una caja acotada de coordenadas. Elegimos una preimagen entera de cada uno y tomamos el máximo de sus costes. Este argumento no afirma que todo vector entero del espacio real generado tenga una preimagen entera.

Lo aplicamos primero a las columnas globales, en la caja \([-2K,2K]\), y llamamos \(C_1\) a la cota resultante. Realizamos la combinación de tipos elegida sobre los representantes. Su cadena \(z_0\) tiene coste a lo sumo \(C_1\); el residuo \(u=\mathbf1_P-\partial z_0\) tiene conteos globales nulos, sigue siendo compatible sobre los enteros y satisface \[\|u\|_1\le L:=K+3C_1.\] Elegimos en cada clase un representante distinguido \(o_i\). Aplicamos el argumento de la caja finita a cada matriz de perfiles locales en la caja \([-2L,2L]\) y tomamos \(C_2\) como cota de todos los costes de las preimágenes elegidas.

Para \(v\in V_i\setminus\{o_i\}\), representamos \(d_u(v)\) mediante sus columnas locales. Una columna correspondiente a \(\tau=\{i,j,k\}\) se transfiere de \(o_i\) a \(v\) mediante \[[vab]-[o_iab],
\tag{A.4.1}\] donde \(a,b\) son representantes de etiquetas \(j,k\), distintos entre sí y de \(v,o_i\). Cinco representantes bastan incluso cuando coinciden etiquetas. En \(a,b\), los perfiles por clases se cancelan porque \(v,o_i\) tienen la misma etiqueta. Los conteos globales también se cancelan.

Usamos la corrección nula si el perfil deseado es cero. En otro caso, su norma entera \(\ell^1\) es al menos uno. Por tanto, la corrección en \(v\) cuesta a lo sumo \(2C_2\|d_u(v)\|_1\). Por (A.1.1), su coste total es a lo sumo \(C_g=4C_2L\). Tras restarlas, los perfiles de todos los vértices no distinguidos son cero. El perfil en \(o_i\) también es cero, pues la suma de los perfiles en \(V_i\) era la fila global nula. El residuo \(u'\) tiene todos sus perfiles nulos y \[\|u'\|_1\le L+3C_g.
\tag{A.4.2}\]

Damos ahora el cálculo del núcleo que incluye los bloques internos completos. Para una arista \(xy\) de etiquetas \(i,j\), introducimos el vector formal \[q(xy)=[xy]-[x o_j]-[o_i y]+[o_i o_j].
\tag{A.4.3}\] Los lazos formales sólo se permiten en este cálculo. Al sumar (A.4.3) con coeficientes \(u'_{xy}\) se obtiene \(u'\): los términos que contienen un vértice distinguido se anulan por los perfiles nulos y los términos finales por los conteos globales nulos. En un bloque diagonal, esta última afirmación utiliza la mitad del conteo diagonal ordenado, que también es cero.

Para \(i\ne j\), (A.4.3) es un cuadrilátero, o cero si uno de los extremos es distinguido. Tiene un certificado triangular de coste a lo sumo cuatro, obtenido eligiendo un ápice que lo complete y formando conos sobre sus cuatro aristas con signos alternados. Un quinto representante de la etiqueta del ápice permite evitar las cuatro esquinas, incluso si esa etiqueta coincide con \(i\) o \(j\).

Para un bloque interno, fijemos representantes \(p_i,r_i\) distintos de \(o_i\) y entre sí. Restamos \(q(p_ir_i)\) de \(q(xy)\) siempre que \(x,y\ne o_i\). Esto no cambia la suma total: la suma de los coeficientes de las aristas internas no incidentes con \(o_i\) es la suma interna total menos la fila nula en \(o_i\) y, por tanto, es cero. El lazo formal \([o_io_i]\) se cancela ahora en cada diferencia.

Escribimos \(C(a,b,c,d)=[ab]-[bc]+[cd]-[da]\). Cuando \(x,y\) evitan \(p_i,r_i\), \[q(xy)-q(p_ir_i)
=C(x,y,o_i,p_i)+C(x,p_i,r_i,o_i).
\tag{A.4.4}\] Las cuatro esquinas de cada cuadrilátero son distintas. Si un extremo coincide con \(p_i\), la expresión es \(C(p_i,y,o_i,r_i)\); si coincide con \(r_i\), intercambiamos \(p_i,r_i\). Invertimos la arista para cubrir las coincidencias en el otro extremo. Si sus extremos son exactamente \(p_i,r_i\), la expresión es cero. Esto agota los casos. Formar conos sobre estos dos cuadriláteros como máximo da coste a lo sumo ocho por coeficiente de arista original. Ningún lazo formal se utiliza como arista efectiva.

Hemos demostrado que \(u'\) tiene un certificado de coste a lo sumo \(8\|u'\|_1\). Al añadir las dos correcciones anteriores, obtenemos la proposición con \[C=8(L+3C_g)+C_1+C_g,
\qquad L=K+3C_1,\quad C_g=4C_2L.
\tag{A.4.5}\] Todas las elecciones utilizaron únicamente los vértices originales y los cinco representantes por clase. \(\square\)

\subsubsection{A.5. De un certificado firmado a un absorbedor físico}\label{a.5.-de-un-certificado-firmado-a-un-absorbedor-fuxedsico}

\textbf{Proposición A.5.1.} Sea \(P\) un grafo con un certificado entero de \(k\) apariciones triangulares con signo sobre un conjunto etiquetado de vértices \(W\). Existe un grafo abstracto etiquetado \(A\), sin aristas entre dos vértices de \(W\), tal que tanto \(A\) como \(A\cup P\) tienen descomposiciones en triángulos permitidos. Tiene a lo sumo \(9k\) vértices nuevos y \(30k\) aristas. Los vértices nuevos admiten un orden en el que cada uno tiene a lo sumo cuatro vecinos entre \(W\) y los vértices nuevos anteriores.

\textbf{Demostración.} Expandimos cada coeficiente entero en apariciones unitarias, conservando su signo. Para una aparición \(uvw\), introducimos clones privados \(u',v',w'\) con las mismas etiquetas respectivas. Tomamos el grafo de nueve aristas formado por el triángulo de clones y las seis aristas que unen un vértice antiguo con los clones de las otras dos esquinas. Tiene una descomposición \[(u v'w'),\quad(v u'w'),\quad(w u'v').
\tag{A.5.1}\] Alternativamente, apartamos su triángulo de clones. Las aristas restantes forman tres cuñas de dos aristas, una sobre cada lado del triángulo original. La cuña sobre \(uv\) tiene ápice \(w'\).

Para cada arista antigua \(uv\), emparejamos inyectivamente sus apariciones como lado negativo con sus apariciones como lado positivo. La ecuación del certificado dice que queda exactamente un lado positivo sin emparejar si \(uv\in P\), y ninguno en otro caso. Para un par emparejado, sean \(c,d\) los ápices de las cuñas negativa y positiva. Introducimos dos vértices privados \(a,b\), con las etiquetas de \(u,v\), y el grafo de siete aristas \[H=\{ub,av,ab,ac,bc,ad,bd\}.\] Entonces \(H\) junto con la cuña de ápice \(c\) se descompone mediante \((ubc),(avc),(abd)\), y junto con la cuña de ápice \(d\) mediante \((ubd),(avd),(abc)\). Todos esos triángulos están permitidos: cada uno tiene el tipo de una de las dos apariciones originales.

Sea \(A\) la unión de los grafos de nueve aristas de las apariciones y de estas bisagras de siete aristas. Para descomponer \(A\), usamos (A.5.1) en las apariciones positivas, el triángulo de clones en las negativas y cada bisagra con su cuña negativa. Para descomponer \(A\cup P\), usamos (A.5.1) en las apariciones negativas, el triángulo de clones en las positivas y cada bisagra con su cuña positiva emparejada. Cada cuña positiva no emparejada se completa con su arista correspondiente de \(P\). Estas listas particionan las aristas: los clones y los vértices de las bisagras son privados, y el emparejamiento de lados es inyectivo.

Si hay \(k_-\) apariciones negativas, hay \(3k_-\) bisagras. Así, \(|E(A)|=9k+21k_-\le30k\) y el número de vértices nuevos es \(3k+6k_-\le9k\). Ordenamos primero todos los clones de las apariciones, después el primer vértice de cada bisagra y luego el segundo. Un clon tiene a lo sumo dos vecinos antiguos y dos clones anteriores como vecinos. El primer vértice de una bisagra tiene vecinos \(v,c,d\); el segundo tiene \(u,c,d,a\). Esto demuestra la afirmación sobre el orden. \(\square\)

\textbf{Corolario A.5.2 (absorción terminal en el anfitrión).} Sea \(C=C(F,K)\) como en la Proposición A.4.1. Supongamos que un anfitrión con \(t\) etiquetas proporciona al menos \(a\) elecciones compatibles para cada conjunto de a lo sumo cuatro vecinos. Sea \(D\) un registro de aristas prohibidas de grado máximo a lo sumo \(d\), y \(Z\) un conjunto de vértices auxiliares prohibidos. Si \[|Z|+2K+5t+9C+4d<a,
\tag{A.5.2}\] entonces todo \(P\subseteq G\) compatible sobre los enteros, con \(|P|\le K\), tiene un absorbedor \(A\subseteq G\) con a lo sumo \(30C\) aristas, disjunto de \(D\cup P\) y que no utiliza vértices de \(Z\) fuera de \(V(P)\).

\textbf{Demostración.} Elegimos cinco representantes nuevos en cada clase fuera de \(Z\cup V(P)\). Esto es posible por (A.5.2), usando la disponibilidad sin restricciones de vecindad. Aplicamos la Proposición A.4.1 en el grafo completo permitido virtual sobre estos representantes y \(V(P)\), de tamaño a lo sumo \(2K+5t\), y después la Proposición A.5.1. Insertamos sus vértices nuevos en el orden que ésta proporciona. En cada paso, las adyacencias requeridas en el anfitrión dejan al menos \(a\) candidatos. Excluir \(Z\), los vértices fijos, los nuevos anteriores y los vecinos prohibidos por el registro de los a lo sumo cuatro vecinos anteriores elimina menos de \(a\) candidatos por (A.5.2). La inserción es, por tanto, inyectiva y evita el registro. La construcción no tiene aristas entre vértices fijos, de modo que también evita \(P\). Las dos descomposiciones abstractas se convierten en descomposiciones físicas del anfitrión. \(\square\)

\subsubsection{A.6. Absorber un catálogo fijo sin concentrar la carga auxiliar}\label{a.6.-absorber-un-catuxe1logo-fijo-sin-concentrar-la-carga-auxiliar}

El siguiente enunciado de ensamblaje separa la construcción de un catálogo de su utilización. No afirma la existencia del catálogo latino.

\textbf{Proposición A.6.1.} Sean \(X,R\subseteq G\) disjuntos en aristas y sea \((B_j)_{j=1}^M\) un catálogo fijo de subgrafos de \(X\cup R\) compatibles sobre los enteros, cada uno con a lo sumo \(K\) aristas. Supongamos que, para todo \(L\subseteq X\) admisible, incluido \(L=\varnothing\), una subcolección del catálogo particiona \(R\cup L\) en piezas disjuntas en aristas. Supongamos que \[\begin{gathered}
\Delta(X\cup R)\le D_0,\\
\sum_j d_{B_j}(v)\le\Lambda\quad\hbox{for every }v.
\end{gathered}
\tag{A.6.1}\] Ponemos \(C=C(F,K)\) y \(c=30C\). Supongamos que la disponibilidad para cuatro vecinos es al menos \(a\) y elijamos \(T>0,D\) tales que \[\begin{split}
4(|X\cup R|+Mc)&\le Ta,\\
\max\{D_0,T+c\}+c\Lambda&\le D,\\
2K+5t+9C+4D&<a/2.
\end{split}
\tag{A.6.2}\] Entonces existe un grafo fijo \(A\subseteq G-X\) con \(\Delta(A)\le D\) tal que \(A\cup L\) tiene una descomposición en triángulos permitidos para todo \(L\) admisible.

\textbf{Demostración.} Construimos sucesivamente absorbedores privados \(A_j\). Antes del paso \(j\), el registro es \(X\cup R\cup\bigcup_{i<j}A_i\); marcamos como cargados los vértices cuyo grado en el registro supera \(T\). El número total de aristas del registro es a lo sumo \(|X\cup R|+Mc\), por lo que (A.6.2) implica que hay a lo sumo \(a/2\) vértices cargados. Aplicamos el Corolario A.5.2 a \(B_j\), tomando el registro como aristas prohibidas y los vértices cargados como auxiliares prohibidos.

Damos la inducción sobre los grados que justifica cada aplicación sin suponer la conclusión. Un paso añade a lo sumo \(c\) en cualquier vértice. Si se usa un vértice como auxiliar, su grado inmediatamente antes de ese paso era a lo sumo \(T\). Después de su último uso auxiliar, sólo puede recibir contribuciones en pasos en los que es raíz de \(B_j\). Cada uno contribuye a lo sumo \(c\), y el número de esos pasos es a lo sumo \(\sum_j d_{B_j}(v)\le\Lambda\). Un vértice que nunca se utiliza como auxiliar empieza con grado a lo sumo \(D_0\). Por tanto, todo prefijo del registro tiene grado máximo a lo sumo \(\max\{D_0,T+c\}+c\Lambda\le D\). Junto con la última desigualdad de (A.6.2), esto valida (A.5.2) en el paso siguiente y completa la inducción.

Pongamos \(A=R\cup\bigcup_j A_j\). Para un \(L\) admisible, elegimos índices \(J\) cuyas piezas particionen \(R\cup L\). Descomponemos \(A_j\cup B_j\) para \(j\in J\) y \(A_j\) para \(j\notin J\). Estas descomposiciones son disjuntas en aristas y cubren exactamente \(A\cup L\), pues todos los absorbedores privados evitan \(X\cup R\) y se evitan entre sí. La cota de grado se deduce de la cota del registro. El caso \(L=\varnothing\) también descompone \(A\). Todas las elecciones para construir \(A\) preceden a la elección de \(L\). \(\square\)

\bigskip

\Needspace{12\baselineskip}\subsection{Apéndice B. Enlaces de tres pasos y ensamblaje cuantitativo}\label{apuxe9ndice-b.-enlaces-de-tres-pasos-y-ensamblaje-cuantitativo}

\subsubsection{B.1. La afirmación y el orden de construcción}\label{b.1.-la-afirmaciuxf3n-y-el-orden-de-construcciuxf3n}

Escribimos \(F\) para una familia fija de ternas de un conjunto finito de etiquetas y decimos que \(ij\) es activo si \(ijk\in F\) para algún \(k\). Un anfitrión está etiquetado por \(p\). Una disponibilidad protegida para \(s\) vértices de tamaño \(a\) significa que todo conjunto de a lo sumo \(s\) vértices cuyas etiquetas sean activas con \(i\) tiene al menos \(a\) vecinos comunes de etiqueta \(i\). Utilizamos las condiciones locales y globales de divisibilidad entera de este artículo, incluidas las multiplicidades de las etiquetas repetidas.

Demostramos la siguiente afirmación intermedia. Si \(F\) tiene enlaces de tres pasos, existe un entero \(S=S(F)\ge4\) tal que, para todo \(0<\omega\le1\) y \(\delta>0\), existen \(\rho>0,N\) con esta propiedad: en todo anfitrión permitido de orden \(n\ge N\) con disponibilidad protegida \((S,\omega n)\), cada \(X\subseteq E(G)\) con \(\Delta(X)\le2\rho n\) tiene un conjunto fijo \(A\subseteq E(G)\setminus X\) que satisface

\[\begin{gathered}
\Delta(A)\le2\delta n,\\
A\text{ y }A\cup L\text{ admiten }F\text{-descomposiciones}
\end{gathered}
\tag{B.1.1}\]

para todo \(L\subseteq X\) integralmente divisible. El orden de elección es \(F,S;\omega,\delta,\rho,N;G;X;A;L\). En particular, ni la biblioteca ni su realización física se eligen después de \(L\).

Combinada con la inversa real del Apéndice A y los lemas de empaquetamiento y cobertura descendente, (B.1.1) demuestra el Teorema A de Paper II. La hipótesis de que cada etiqueta tiene una pareja prohibida sigue formando parte de las condiciones sobre la familia de ese teorema; no sustituye ninguna de las condiciones de retículo local que siguen.

\subsubsection{B.2. Qué proporciona un enlace de tres pasos}\label{b.2.-quuxe9-proporciona-un-enlace-de-tres-pasos}

Fijemos una etiqueta raíz \(r\). Sea \(L_r\) su enlace, permitiendo lazos. Sus columnas enteras de incidencia son \(e_a+e_b\) para \(rab\in F\); un lazo da \(2e_a\). A lo largo de un recorrido \(a,v,w,b\) de longitud tres,

\[(e_a+e_v)-(e_v+e_w)+(e_w+e_b)=e_a+e_b.
\tag{B.2.1}\]

Así, los enlaces de tres pasos sitúan cada columna de par activo en el retículo local. Recíprocamente, toda columna del enlace tiene suma de coordenadas par. Fijemos una etiqueta activa \(a\). Si un vector entero soportado \(d\) tiene suma \(2z\), entonces

\[d=\sum_jd_j(e_j+e_a)-z(2e_a).
\tag{B.2.2}\]

Las ecuaciones (B.2.1)--(B.2.2) demuestran que la divisibilidad local equivale a tener grado total par en el soporte activo. Esto \textbf{no} elimina la condición del retículo global. Sea \(\mathcal Q_F\) el grupo de vectores enteros de tipos de aristas módulo el subgrupo generado por las columnas triangulares. Un grafo activo sin lazos es integralmente divisible exactamente cuando sus grados son pares y su clase en \(\mathcal Q_F\) es cero. Conservamos este grupo abeliano, incluida su torsión.

\subsubsection{B.3. Una biblioteca par fija y su posterior realización mediante caminos}\label{b.3.-una-biblioteca-par-fija-y-su-posterior-realizaciuxf3n-mediante-caminos}

En B.3, \(n\) es el orden del grafo ambiente, mientras que H.6 utiliza primero el orden local \(N=|W|\). El reescalado al final de H.6 da la formulación con parámetros ambientales que sigue; en particular, \(M=2^h\) en (B.3.2).

Recordamos primero el refinador ordinario en la forma que utilizamos. El dato \(X\) es un grafo sin lazos sobre \(W\) con \(\Delta(X)\le2\rho n\). Su alcance preciso y su fuente se indican en el Apéndice G.4. Sobre un conjunto de vértices \(W\), su grafo auxiliar puede ser completo. Para parámetros positivos \(\lambda,\rho\) con

\[64\lambda n\le |W|,\quad 2^{24}\rho\le\lambda^2,
\quad \rho n\ge1,\quad\lambda n\ge256,
\tag{B.3.1}\]

refina todos los subgrafos pares de \(X[W]\) cuyo número de aristas es divisible por tres en una biblioteca fija indexada de piezas pares de a lo sumo 42 aristas. Su resto fijo tiene grado máximo a lo sumo \(4\lambda n\), el número de índices es a lo sumo \(1224\rho n^2\) y la suma de los grados de \textbf{todas} las piezas de la biblioteca en cualquier vértice es a lo sumo \(4096\lambda n\).

El mecanismo finito de este teorema auxiliar resulta útil aquí. Un árbol binario envía cada vértice a un centro de un conjunto de tamaño \(2^h\) y después envía cada centro a su padre. Elegimos

\[64\rho n\le\lambda2^h\le128\rho n.
\tag{B.3.2}\]

En cada nivel, los bits de las aristas seleccionadas se transportan por cadenas de extremos y cadenas de pares de centros. Una cadena de bits \(x_j\) usa los acarreos \(c_{j+1}=c_j+x_j\pmod2\). Sus dos caminos alternativos de dos aristas se asignan según el acarreo entrante. Cada conexión se asigna una vez y su frontera cancela los dos acarreos adyacentes. En el último nivel del árbol no queda ningún par. Una segunda cadena transporta el conteo de aristas módulo tres, utilizando tres ciclos de cuatro vértices alternativos por conexión. Hasta dos ciclos de ajuste utilizados siempre fijan el residuo del resto. Conservar de antemano todas las alternativas da la biblioteca universal. Las cotas anteriores cuentan todas las alternativas y todos los caminos solicitados, no sólo los caminos de una selección final. Éste es el refinador universal ordinario que utilizamos; en particular, la cota 42 corresponde a sus piezas finales módulo tres.

Lo aplicamos a los conjuntos de vértices de las componentes conexas del grafo de etiquetas activas. La restricción de un \(L\) integralmente divisible a una componente de este tipo es par y tiene un número de aristas divisible por tres: cada columna triangular está íntegramente en una componente. Los conjuntos de vértices de componentes distintas son disjuntos. Ampliando las cotas, ponemos

\[K_v=4096t,\qquad \ell=t+1,\]

donde \(t\) es el número de etiquetas. Obtenemos un resto virtual \(R_v\) y una biblioteca \(P_i^v\) de piezas de tamaño a lo sumo 42, con número de piezas a lo sumo \(K_v\rho n^2\) y carga total por vértice a lo sumo \(K_v\lambda n\). Las aristas del resto virtual no tienen por qué ser aristas permitidas del anfitrión.

Sustituimos una vez cada arista del resto virtual por un camino físico. Las etiquetas de sus extremos pertenecen a una misma componente activa. Un camino mínimo de etiquetas, con un desvío de dos aristas si es necesario, tiene longitud entre dos y \(t+1\). Todos sus pares consecutivos son activos. Damos a una arista virtual peso

\[w(e)=|\{i:e\in P_i^v\}|.
\tag{B.3.3}\]

La selección del residuo vacío cubre \(R_v\), de modo que allí \(w(e)\ge1\). Además,

\[\sum_{e\ni v}w(e)\le K_v\lambda n,
\qquad \sum_{e\in R_v}w(e)\le42K_v\rho n^2.
\tag{B.3.4}\]

Realizamos los caminos sucesivamente, evitando \(X\), las aristas de caminos anteriores, los vértices repetidos del camino actual y los vértices cuya carga auxiliar ponderada acumulada sea al menos \(4K_v\lambda n\). La carga auxiliar total es a lo sumo \(2\ell\sum_e w(e)\); por tanto, el número de vértices cargados es a lo sumo \(21\ell\rho n/\lambda\). El último vértice interno puede requerir dos adyacencias especificadas, razón por la que se utiliza disponibilidad de orden dos. Los márgenes suficientes son

\[64K_v\lambda\le\omega,\quad
128\ell\rho\le\omega\lambda,\quad
\rho\le\lambda,\quad4(\ell+1)\le\omega n.
\tag{B.3.5}\]

Dejan un candidato en cada paso. Un vértice interno recibe a lo sumo un exceso de tamaño \(2w(e)\le2K_v\lambda n\), por lo que su carga auxiliar ponderada es a lo sumo \(6K_v\lambda n\). Incluyendo los extremos originales, la carga total de la biblioteca es a lo sumo \(7K_v\lambda n\). Se utiliza el mismo camino físico en todas las alternativas que contienen su arista virtual. En consecuencia, toda selección virtual disjunta se convierte en una selección física disjunta; sus aristas son exactamente el resto fijo realizado junto con \(L\). La subdivisión preserva la paridad. El tamaño de las piezas es a lo sumo \(42(t+1)\).

Como referencia, la construcción congelada utiliza

\[c_0=\min\left\{2^{-24},\frac1{1024(4096t)(t+1)}\right\},\quad
m_0=42(t+1),\quad K_0=7\cdot4096t.
\tag{B.3.6}\]

Tomar \(\lambda\le c_0\omega\), \(\rho\le c_0\lambda^2\), \(\rho n\ge1\) y \(c_0\lambda n\ge1\) implica (B.3.1) y (B.3.5). Hemos obtenido piezas pares, todavía no piezas globalmente divisibles.

\subsubsection{B.4. Eliminación de todas las cargas globales}\label{b.4.-eliminaciuxf3n-de-todas-las-cargas-globales}

Hay un número finito de grafos pares activos etiquetados sin lazos con a lo sumo \(m_0\) aristas y \(2m_0\) vértices. Enumeramos sus clases en \(\mathcal Q_F\) como \(g_1,\ldots,g_r\). Cada pieza de la biblioteca tiene una de esas clases. Para cada generador, elegimos su patrón como ficha positiva. Para construir una ficha negativa, sustituimos cada arista \(ab\) de ese patrón por \(az_e,z_eb\), eligiendo la etiqueta de \(z_e\) de modo que \(abz_e\) sea un triángulo permitido. Módulo las columnas triangulares, esto cambia la clase de la arista por su opuesta. Ambas fichas son pares y sus tamaños son a lo sumo \(2m_0\). Son grafos finitos efectivos, no aristas negativas formales.

Utilizamos la cadena descendente

\[H_0=\langle g_1,\ldots,g_r\rangle\supseteq H_1
=\langle g_2,\ldots,g_r\rangle\supseteq\cdots\supseteq H_r=0.
\tag{B.4.1}\]

Supongamos que las piezas actuales tienen valores en \(H_{j-1}\). Elegimos enteros \(c_i\) con \([P_i]-c_ig_j\in H_j\). Si el orden de \(g_j+H_j\) es \(d>0\), elegimos \(0\le c_i<d\). Si su orden es infinito, ponemos \(d=0\) y utilizamos representantes enteros cualesquiera. El tamaño acotado de las piezas actuales da una cota fija \(a_j\ge\max_i|c_i|,d\).

Para cada familia seleccionada \(T\) que refina \(R\cup L\), la divisibilidad de \(L\) da \[\sum_{i\in T}[P_i]=[R].
\tag{B.4.2}\]

Sea \(T_0\) una selección para el residuo vacío. Con \(K=(\sum_{i\in T_0}c_i)\bmod d\), donde la reducción módulo cero es la identidad, (B.4.2) implica

\[d\mid \sum_{i\in T}c_i-K.
\tag{B.4.3}\]

Para \(d=0\), todas las demandas suman exactamente \(K\). Para \(d>0\), añadimos \(a_j|\mathcal P|\) \textbf{posiciones vacías fijas} de demanda \(-d\) y seleccionamos exactamente

\[s_T=\frac{\sum_{i\in T}c_i-K}{d}\]

de ellas. La convención de representantes demuestra \(0\le s_T\le a_j|\mathcal P|\). Todas las selecciones ampliadas tienen ahora la misma demanda total. Una posición vacía recibirá \(d\) fichas positivas, cuya clase conjunta pertenece a \(H_j\). Éste es el paso que conserva la torsión; sustituir (B.4.3) por una igualdad sin estas posiciones sería incorrecto.

El registro fijo de fichas es el siguiente. Una pieza de demanda \(c>0\) toma \(c\) fichas negativas. Una pieza de demanda \(c<0\) toma \(-c\) fichas positivas privadas. Cada ficha positiva privada tiene una pareja negativa; si no se selecciona su pieza madre, el par se selecciona como pieza de liberación. Las posiciones variables negativas restantes tienen cardinalidad \(\sum_i\max(c_i,0)-K\), independiente de \(T\). El lema de emparejamiento robusto acotado proporciona un grafo bipartito fijo que empareja exactamente ese número de posiciones variables, junto con sus posiciones de reserva fijas, con reservas positivas, para cada elección permitida de las posiciones variables. Una arista del grafo especifica una pieza de cancelación de dos fichas. Por tanto, un emparejamiento perfecto consume cada ficha restante exactamente una vez. Las piezas antiguas seleccionadas, los pares de liberación y los pares del emparejamiento particionan la selección antigua y \textbf{todas} las fichas nuevas.

En la versión cuantitativa utilizada en la versión congelada, un registro de demanda absoluta total \(z\) requiere a lo sumo \(9z\) fichas; cada ficha aparece en a lo sumo 44 piezas de la biblioteca, y el número de índices nuevos es a lo sumo el número antiguo más \(97z\). Incluidas las posiciones vacías, obtenemos

\[\begin{gathered}
\#\text{fichas}\le9a_j(a_j+1)|\mathcal P|,\\
|\mathcal P'|\le[1+98a_j(a_j+1)]|\mathcal P|.
\end{gathered}
\tag{B.4.4}\]

El lema de emparejamiento robusto es un teorema auxiliar finito, no una solicitud de recursos nuevos del anfitrión dependientes de \(L\): su grafo y cada ficha se fijan antes de \(T\). La inserción en el anfitrión se realiza a continuación.

Insertamos las fichas de forma voraz, disjuntas en aristas y fuera de \(R\cup X\), utilizando disponibilidad protegida hasta el orden máximo de las fichas. Limitamos su carga acumulada por vértice a \(\Delta=\lceil\lambda n\rceil+2m_0+1\le2\lambda n\). Marcamos los vértices de carga superior a \(\Delta-2m_0\); la masa total de aristas acota su número. Evitarlos impide más de un exceso. Las constantes de esta iteración quedan fijadas por una recursión finita, no por el anfitrión. Pongamos \(N_0=2m_0\) y enumeremos, permitiendo repeticiones, las cargas de todos los patrones activos, sin lazos y pares sobre \(N_0\) posiciones etiquetadas, con a lo sumo \(m_0\) aristas. Denotemos esta lista por \(g_1,\ldots,g_r\). En particular, \[r\le t^{N_0}2^{\binom{N_0}{2}}.\] Sea \(H_j=\langle g_{j+1},\ldots,g_r\rangle\). El cociente \(H_{j-1}/H_j\) es cíclico; escribamos \(d_j=0\) si su orden es infinito y, en otro caso, sea \(d_j\) su orden. Elijamos un coeficiente entero \(a_j(x)\) que represente cada \(x\in H_{j-1}\) módulo \(H_j\), con \(0\le a_j(x)<d_j\) cuando \(d_j>0\). Para una cota \(m_j\) del tamaño de las piezas, las cargas posibles pertenecen al conjunto finito \(\mathcal V(m_j)\) de imágenes de vectores de aristas por tipos ordenados cuyas coordenadas están en \(\{0,\ldots,2m_j\}\). Partiendo de \(m_1=m_0\), definamos \[\begin{aligned}
A_j&=d_j+\max\bigl(\{0\}\cup\{|a_j(x)|:x\in\mathcal V(m_j)\cap H_{j-1}\}\bigr),\\
m_{j+1}&=m_j+2m_0A_j+4m_0,\\
L_j&=1+98A_j(A_j+1).
\end{aligned}\] Fijemos \(\mathcal L_{r+1}=1\), \(\mathcal T_{r+1}=0\) y, recorriendo la lista hacia atrás, \[\begin{gathered}
\mathcal L_j=45L_j\mathcal L_{j+1},\\
\mathcal T_j=9A_j(A_j+1)+L_j\mathcal T_{j+1}.
\end{gathered}\] Utilizamos \(m=m_{r+1}\), \(L_c=\mathcal L_1\) y \(T_c=\mathcal T_1\); para la lista vacía son \(m_0,1,0\). La construcción de una capa utiliza a lo sumo \(9A_j(A_j+1)\) fichas por índice entrante, multiplica el número de índices por a lo sumo \(L_j\) y acota la nueva carga con multiplicidades mediante el factor \(45L_j\). Aplicando sucesivamente estas cotas se obtiene exactamente la recursión anterior. Los coeficientes y máximos se eligen en el grupo abeliano finitamente generado fijo, antes del anfitrión o del residuo. Así quedan definidas constantes finitas, sin afirmar una cota cerrada optimizada.

Una elección explícita suficiente es la siguiente. Pongamos \(N_0=2m_0\), \(U=N_0+N_0(N_0+1)/2\) y

\[\begin{gathered}
B=(U+1)(K_0+r+3)(2m_0+3)(8m_0T_c+1)(K_0+1),\\
c=\min\{c_0,1,(8B)^{-1}\}.
\end{gathered}
\tag{B.4.5}\]

Las hipótesis \(\lambda\le c\omega\), \(\rho\le c\lambda^2\), \(\rho n\ge1\) y \(c\lambda n\ge1\) dan ambas desigualdades

\[U+U(K_0\lambda n+2\rho n+r\Delta+\Delta)<\omega n/2,\] \[8m_0T_cK_0\rho n^2\le(\Delta-2m_0)\omega n.
\tag{B.4.6}\]

La primera desigualdad permite evitar los vecinos ya utilizados y mantener vértices distintos; la segunda permite evitar los vértices cargados. Justifican cada inserción de fichas. Tras la última capa, cada pieza tiene clase cero en \(\mathcal Q_F\) y sigue siendo par. La Sección B.2 demuestra su divisibilidad entera completa. Ampliando \(K\) a \(L_c(K_0+2r+2)\), obtenemos una biblioteca divisible fija con

\[|P_i|\le m,\quad |\mathcal P|\le K\rho n^2,\quad
\Delta R\le K\lambda n,\quad
\sum_i d_{P_i}(v)\le K\lambda n.
\tag{B.4.7}\]

\subsubsection{B.5. De la biblioteca al absorbedor universal}\label{b.5.-de-la-biblioteca-al-absorbedor-universal}

Sea \(C_{\rm cert}=C(F,m)\) el coste del certificado firmado proporcionado por A.4 y pongamos \(C_{\rm abs}=30C_{\rm cert}\). El Corolario A.5.2 acota las aristas del absorbedor por \(C_{\rm abs}\); estas dos constantes cumplen funciones distintas. Definamos \[S_t=\left\lceil\max\{4,2m+5t+9C_{\rm cert}\}\right\rceil,
\qquad S=\max(S_t,U,S_{\rm ref},4).\] Aquí \(U\) es el orden de la plantilla en (B.4.5), y \(S_{\rm ref}=\max(2,U)\) basta para el refinador enrutado y sus plantillas de carga. El valor mayor de \(S_t\) es una cota de presupuesto además de un orden de disponibilidad. Reduzcamos el presupuesto de grado deseado a

\[\delta'=\min\{\delta,\omega/[8(S_t+1)]\}.\]

Utilizamos los siguientes parámetros en (B.4.7):

\[\begin{gathered}
\lambda=\frac{\min(c,1)\omega\delta'}{16(K+1)(C_{\rm abs}+1)},\\
\rho=\frac{\min(c,1)\lambda^2}{16(K+1)(C_{\rm abs}+1)}.
\end{gathered}
\tag{B.5.1}\]

Basta tomar

\[N\ge\lceil1/\rho\rceil+\lceil1/(c\lambda)\rceil
+\lceil2C_{\rm abs}/\delta'\rceil+\lceil4(S_t+1)/\omega\rceil.
\tag{B.5.2}\]

Damos a cada índice de la biblioteca un absorbedor terminal privado \(A_i\), disjunto de \(X,R\) y de los \(A_j\) insertados anteriormente. La inserción con control de carga de la Sección A.6 se aplica porque

\[S_t+2S_t\delta'n<\omega n/2,\] \[\max\{(2\rho+K\lambda)n,\delta'n+C_{\rm abs}\}+C_{\rm abs}K\lambda n\le2\delta'n,\] \[4(2\rho+K\lambda+K\rho C_{\rm abs})\le\delta'\omega.
\tag{B.5.3}\]

Estas desigualdades controlan, respectivamente, la evitación de aristas utilizadas, la carga en las raíces y el número de vértices cargados. Se deducen por sustitución en (B.5.1), usando \(2C_{\rm abs}\le\delta'n\). Pongamos \(A=R\cup\bigcup_i A_i\). Para un \(L\subseteq X\) divisible elegido, utilizamos una selección disjunta de la biblioteca \(R\cup L=\dot\bigcup_{i\in T}P_i\). Descomponemos \(A_i\cup P_i\) para \(i\in T\) y \(A_i\) para los índices restantes. Tomar \(L=\varnothing\) también descompone \(A\). Esto demuestra (B.1.1), con sus cuantificadores universales.

La implicación hacia A.6 se puede comprobar ahora sin una constante implícita. Tomemos \(K=m\), \(a=\omega n\), \(T=\delta'n\), \(D=2\delta'n\), \(c=C_{\rm abs}\) y \(\Lambda=K_{\rm lib}\lambda n\), donde \(K_{\rm lib}\) denota la constante de biblioteca escrita \(K\) en (B.4.7). La segunda desigualdad de (B.5.3) da directamente la desigualdad de carga de A.6; multiplicar la tercera por \(n^2\) da su desigualdad de masa. La condición restante se deduce de \[2m+5t+9C_{\rm cert}+8\delta'n
\le S_t+2S_t\delta'n<\omega n/2.\] En efecto, (B.5.2) da \(S_t<\omega n/4\), mientras que \(\delta'\le\omega/[8(S_t+1)]\) da \(2S_t\delta'n<\omega n/4\). Esto explica también por qué el coste del certificado no puede sustituir a la cota de aristas del absorbedor en las desigualdades de carga o masa.

\subsubsection{B.6. Deducción cuantitativa del Teorema A}\label{b.6.-deducciuxf3n-cuantitativa-del-teorema-a}

En (2.7), elegimos \(h=S\), \(\omega_0=1\) y tomamos como tolerancia de regularidad \(c\) el valor \(\varepsilon\) de (B.6.2). Todas las constantes que definen \(S\) dependen solo de la familia.

Sea \(C_F\) la constante de la inversa completa de la Sección A.2 y sea \(H_F>0\) la constante correspondiente para la inversa en el anfitrión físico de la Sección A.3. Pongamos

\[\begin{gathered}
\kappa=\min\{\omega/8,\omega^{12}/(216H_F)\},\\
\varepsilon_{\rm inv}=\omega^{11}/(4H_F).
\end{gathered}
\tag{B.6.1}\]

Aplicamos (B.1.1) con \(\delta=\kappa/8\), obteniendo \(\rho_0,N_0\), y ponemos \(\rho=\min(\rho_0,\kappa/8)\). Si \(\tau(\eta)\) denota la tolerancia de carga del lema de empaquetamiento tipado, tomamos

\[\varepsilon=\tfrac12\min\{\varepsilon_{\rm inv},
\tau(\min(\omega/8,1,\omega^2))\}.
\tag{B.6.2}\]

Además de \(N_0\), el umbral de orden es el máximo de

\[\lceil8/\omega\rceil,\quad
\lceil2^{27}\rho^{-9}\omega^{-6}\rceil,\quad
N_{\rm pack}(27/\omega,\min(\omega/8,1,\omega^2)),\quad
N_{\rm cover}(54/\omega,\rho^2\omega/12).
\tag{B.6.3}\]

Son umbrales de los lemas de empaquetamiento y cobertura descendente, fijados antes del anfitrión. Los pesos regulares originales pertenecen a \([\omega/n,1/(\omega n)]\subseteq[\omega^2/n,27/(\omega n)]\). La hipótesis de vecindades comunes proporciona la disponibilidad protegida utilizada antes.

Para comprobar la hipótesis restante del empaquetamiento inicial, fijemos un vértice \(v\) y elijamos mediante (3.2) una clase prohibida \(j\) para su etiqueta. La disponibilidad protegida sin vecinos prescritos da \(|V_j|\ge\omega n\). Como todas las aristas tienen tipos admisibles, \(N_G(v)\cap V_j=\varnothing\), incluso si \(j=p(v)\). En consecuencia, \[d_G(v)+|V_j|\le n,\qquad \Delta(G)\le(1-\omega)n.\] La disponibilidad protegida para dos raíces proporciona además la condición de anti-aristas del Apéndice G.1. Estas son exactamente las hipótesis adicionales de los enunciados de empaquetamiento con cargas próximas a uno y selección de reserva del Apéndice G.2; las condiciones de retículo local se utilizan por separado.

El lema de empaquetamiento proporciona primero un empaquetamiento triangular grande del que se elige una reserva \(X\) descomponible en triángulos, de grado máximo a lo sumo \(2\rho n\), que conserva disponibilidad para antiaristas \(\rho^2\omega n/12\). Construimos \(A\) mediante (B.1.1). Tanto \(A\) como \(X\) son descomponibles. Por tanto, eliminar \(A\cup X\) preserva la compatibilidad local y global requerida por la inversa real del Apéndice A. El grado máximo eliminado es a lo sumo \((2\delta+2\rho)n\le\kappa n/2\). Las ecuaciones (B.6.1)--(B.6.2) cubren el defecto inicial de pesos y la pérdida de triángulos debida a las aristas eliminadas. La reparación real da pesos exactos no negativos en \(G-(A\cup X)\), con cota superior \(54/(\omega n)\).

El lema de cobertura descendente, aplicado con los pesos reparados y la disponibilidad reservada para antiaristas, empaqueta \(G-A\) dejando un \(L\subseteq X\) integralmente divisible. Descomponemos ahora \(A\cup L\). Estas dos descomposiciones disjuntas en aristas cubren \(G\), lo que demuestra el Teorema A. Es esencial recalcular los pesos después del borrado: restringir los pesos originales no justificaría por sí solo este paso.

\bigskip

\Needspace{12\baselineskip}\subsection{Apéndice C. El teorema semidenso de triángulos mixtos}\label{apuxe9ndice-c.-el-teorema-semidenso-de-triuxe1ngulos-mixtos}

\subsubsection{C.1. Enunciado y convenciones}\label{c.1.-enunciado-y-convenciones}

Sea \(F_0=\{001,011\}\). Para todo \(\alpha>0\) existen \(\varepsilon>0,N\) tales que lo siguiente vale para \(n\ge N\). Supongamos que \(V(G)=V_0\dot\cup V_1\) y

\[\begin{gathered}
|V_i|\ge(1-\varepsilon)n/2,\\
(1/4-\varepsilon)n\le d(v,V_i)\le(1/4+\varepsilon)n
\quad(v\in V_i),
\end{gathered}
\tag{C.1.1}\] \[d(v,V_{1-i})\ge(1/2-\varepsilon)n\quad(v\in V_i).
\tag{C.1.2}\]

Para al menos una parte \(D=V_i\), el grafo \(G[D]\) tiene distancia de edición de aristas mayor que \(\alpha|D|^2\) tanto de la unión de dos cliques como del grafo bipartito completo asociados a \textbf{cada} corte \(X,D\setminus X\) con \(|X|=\lfloor|D|/2\rfloor\). Las distancias cuentan aristas no ordenadas. Si todos los grados de \(G\) son pares y

\[e(V_0,V_1)=2[e(V_0)+e(V_1)],
\tag{C.1.3}\]

entonces \(G\) tiene una descomposición en triángulos mixtos.

Las condiciones de (C.1.3) y la paridad de los grados equivalen exactamente a la divisibilidad entera respecto de \(F_0\). Globalmente, cada triángulo mixto tiene una arista interna y dos cruzadas. Localmente, las columnas en cualquiera de las etiquetas generan todos los vectores enteros de dos coordenadas y suma par: las columnas son \((1,1)\) y \((0,2)\), o las obtenidas intercambiando las coordenadas. Esto demuestra la equivalencia, no sólo la necesidad.

\subsubsection{C.2. La no extremalidad proporciona recorridos cortos en la parte buena}\label{c.2.-la-no-extremalidad-proporciona-recorridos-cortos-en-la-parte-buena}

Damos la estimación elemental de estabilidad que se utiliza repetidamente. Sea \(J\) un grafo sobre un conjunto \(D\) de \(m\) vértices con grados en \([(1/2-\theta)m,(1/2+\theta)m]\), donde \(0<\theta\le1/100\) y \(\theta m\ge1\). Para subconjuntos \(U,W\), sea \(a(U,W)\) el número de pares ordenados adyacentes. Supongamos que \(|U|=(1/2\pm\theta)m\), \(|W|\ge(1/2-\theta)m\) y \(a(U,W)\le\theta m^2\). Escribiendo \(w=|U\cap W|\), tenemos

\[\begin{gathered}
|D\setminus(U\cup W)|\le2\theta m+w,
\\
w[(1/2-3\theta)m-w]\le2\theta m^2.
\end{gathered}
\tag{C.2.1}\]

Para la segunda desigualdad, sumamos la cota inferior de grado sobre \(U\cap W\). Los vecinos fuera de \(U\cup W\) contribuyen a lo sumo \(w(2\theta m+w)\); los vecinos en \(U\cup W\) contribuyen a lo sumo \(2a(U,W)\). Si \(w\le m/4\), el corchete de (C.2.1) es al menos \(m/8\), de donde \(w\le16\theta m\). En otro caso, (C.2.1) da \(w\ge(1/2-11\theta)m\).

Movemos a lo sumo \(\theta m+1\) vértices para convertir \(U\) en un conjunto equilibrado \(X\). En el primer caso, \(W\) difiere poco de \(D\setminus X\) y

\[a(X,D\setminus X)\le20\theta m^2+m.\]

En el segundo caso, \(U\setminus W\) tiene tamaño a lo sumo \(12\theta m\), por lo que \(a(U,U)\le13\theta m^2\) y

\[a(X,X)\le15\theta m^2+2m.\]

Las cotas de grado y estas estimaciones dan distancia de edición a lo sumo \(80\theta m^2\) respecto de los modelos equilibrados de unión de cliques y bipartito, respectivamente. Para ver por qué la conclusión utiliza distancia de edición, empleamos las identidades

\[|E(J)\triangle E(K_X\cup K_{D\setminus X})|
=e(K_X\cup K_{D\setminus X})-e(J)+2e_J(X,D\setminus X)\]

y la identidad correspondiente para el modelo bipartito completo. Los términos lineales en \(m\) son a lo sumo sus coeficientes multiplicados por \(\theta m^2\).

Por tanto, cuando la distancia respecto de ambos modelos supera \(100\theta m^2\), cada par ordenado \(u,v\in D\), incluido \(u=v\), tiene más de \(\theta m^2\) recorridos de longitud tres: aplicamos la contraposición con \(U=N_J(u)\), \(W=N_J(v)\). Excluir un conjunto de \(z\) posibles vértices internos descarta a lo sumo \(2zm\) de esos recorridos. Una arista precedente da una disponibilidad similar de recorridos de longitud cuatro. Estas estimaciones se mantienen cuando se exige que el recorrido esté en la vecindad cruzada común de una colección fija acotada de vértices, pues cada uno carece de a lo sumo \(3\varepsilon n/2\) vecinos cruzados.

En todos los usos posteriores elegimos

\[\begin{gathered}
\theta=\min(a,1)/100,\qquad e=\theta/2^{30},\\
\theta n\ge4096.
\end{gathered}
\tag{C.2.2}\]

Aquí \(a\) es el parámetro de no extremalidad del anfitrión considerado. Los tamaños de las partes en (C.1.1) convierten las cotas de grado en las cotas próximas a la mitad del orden utilizadas arriba, mediante \(100e\le\theta\). En particular, el grafo interno bueno tiene muchos triángulos enraizados, tomando \(u=v\). No se requiere una no extremalidad análoga en la otra parte.

\subsubsection{C.3. Un certificado entero explícito sobre triángulos físicos}\label{c.3.-un-certificado-entero-expluxedcito-sobre-triuxe1ngulos-fuxedsicos}

Sea \(P\subseteq E(G)\) divisible y \(|P|\le m_0\), donde finalmente \(m_0=214\). Restamos un triángulo mixto efectivo en cada arista interna de \(P\). Existe por (C.1.2). Denotamos su suma con signos por \(L\). La función entera residual de aristas \(R=1_P-\partial L\) está soportada en aristas cruzadas, tiene grado par en cada vértice y coeficiente total cero por (C.1.3). Además,

\[\|L\|_1\le |P|,\qquad\|R\|_1\le4|P|.
\tag{C.3.1}\]

Vemos \(R\) como una matriz desde la parte buena \(D\) hacia la otra parte \(B\). Escribimos sus sumas de filas como \(2r_u\). Su total es cero, por lo que \(\sum_u r_u=0\). Elegimos una única arista interna \(bb'\in E(G[B])\) cuyos extremos sean adyacentes a todas las raíces de filas con \(r_u\ne0\). El número de raíces residuales es a lo sumo \(8m_0\); la cota de la unión para vecinos cruzados ausentes y el grado mínimo interno garantizan que esta arista existe. En detalle, basta que \(e(10m_0+10)\le1/100\), lo cual vale por (C.2.2) y \(m_0=214\).

Restamos \(\sum_u r_u[ubb']\). La arista interna \(bb'\) se cancela porque \(\sum r_u=0\); la matriz cruzada restante \(R'\) tiene sumas de filas nulas y sumas de columnas pares. Su norma de coeficientes satisface

\[\|R'\|_1\le\|R\|_1+2\sum_u|r_u|\le2\|R\|_1.
\tag{C.3.2}\]

Elegimos un ancla \(b_0\in B\) adyacente a todas las raíces de filas. Para cada columna no nula \(c\), trabajamos en el grafo interno sobre \(D\cap N(c)\cap N(b_0)\), que contiene todas las raíces de esa columna y un ancla común. La Sección C.2 proporciona un recorrido de longitud cuatro desde cada raíz hasta el ancla. Si sus vértices son \(v_0,\ldots,v_4\), la cadena alternada de aristas tiene incidencia sin signos

\[[v_0v_1]-[v_1v_2]+[v_2v_3]-[v_3v_4]
\longmapsto e_{v_0}-e_{v_4}.
\tag{C.3.3}\]

Su coste es cuatro. Un vecino \(w\) del ancla \(v_*\), menos la cadena correspondiente a \(e_w-e_{v_*}\), da \(2e_{v_*}\) con coste a lo sumo cinco. Así, un vector entero \(x\) de total par \(2z\) tiene una representación mediante aristas internas físicas

\[\begin{gathered}
\sum_u x_u T_u+zD_*,\\
\text{coste}\le4\|x\|_1+5|z|\le\tfrac{13}{2}\|x\|_1.
\end{gathered}
\tag{C.3.4}\]

Para cada arista interna \(uv\) de esta representación, tomamos \([uvc]-[uvb_0]\). Las aristas internas se cancelan y el resultado tiene frontera cruzada \(\sum_u x_u([uc]-[ub_0])\). Ambos triángulos son físicos por la restricción de vecinos comunes. El coste es a lo sumo \(13\|x\|_1\). Al sumar sobre las columnas se cancelan todas las aristas de \(b_0\), pues las filas de \(R'\) suman cero. Combinando (C.3.1)--(C.3.4), la corrección de filas y todas las columnas cuestan a lo sumo \(27\|R\|_1\); por tanto, el certificado completo cuesta a lo sumo \(109|P|\le126|P|\).

Hemos demostrado un certificado triangular con signos en el anfitrión efectivo, con una cota independiente de \(n\). Un certificado del retículo de pares de clases no habría proporcionado por sí solo las identidades físicas utilizadas en (C.3.3)--(C.3.4).

\subsubsection{C.4. Bisagras físicas y absorbedores terminales acotados}\label{c.4.-bisagras-fuxedsicas-y-absorbedores-terminales-acotados}

Dos triángulos mixtos distintos que comparten una arista tienen raíces \(012\) y \(013\). Los siguientes intercambios finitos especifican por completo la bisagra. En cada fila, el primer y segundo modo cubren las mismas aristas de bisagra más, respectivamente, las dos aristas no compartidas de \(012\) y \(013\).

{\begin{longtable}[]{@{}
  >{\raggedright\arraybackslash}p{(\linewidth - 6\tabcolsep) * \real{0.2500}}
  >{\raggedright\arraybackslash}p{(\linewidth - 6\tabcolsep) * \real{0.2500}}
  >{\raggedright\arraybackslash}p{(\linewidth - 6\tabcolsep) * \real{0.2500}}
  >{\raggedright\arraybackslash}p{(\linewidth - 6\tabcolsep) * \real{0.2500}}@{}}
\toprule\noalign{}
\begin{minipage}[b]{\linewidth}\raggedright
Caso
\end{minipage} & \begin{minipage}[b]{\linewidth}\raggedright
Vértices nuevos
\end{minipage} & \begin{minipage}[b]{\linewidth}\raggedright
Primer modo
\end{minipage} & \begin{minipage}[b]{\linewidth}\raggedright
Segundo modo
\end{minipage} \\
\midrule\noalign{}
\endhead
\bottomrule\noalign{}
\endlastfoot
Arista interna compartida en la parte buena & 4,5 en esa parte & 042,453,512 & 043,452,513 \\
Arista interna compartida en la otra parte & 4 en esa parte; 5,6,7 en la parte buena & 042,125,456,167,473 & 043,425,156,467,173 \\
Arista cruzada compartida & 4,5 con etiquetas adecuadas & 052,142,453 & 053,143,452 \\
\end{longtable}
}

En la primera fila, las aristas de la bisagra son (04,45,51,42,43,52,53). En la segunda son (04,42,43,25,56,67,73,45,46,47,15,16,17). En la última son (05,14,42,52,43,53,45). Enumerar las aristas de cada modo verifica ambas identidades y la disjunción de aristas dos a dos. Todos los triángulos enumerados son mixtos.

Para una arista interna buena, elegimos el camino (0,4,5,1) entre los recorridos disponibles, exigiendo además que sus vértices internos sean adyacentes a 2 y 3 a través de la partición. Para una arista interna de la otra parte, elegimos primero 4 entre los vecinos internos de 0 y los vecinos cruzados comunes de 2,3. Hay al menos \(n/8\) elecciones. Después elegimos el camino (2,5,6,7,3) en la parte buena con sus vértices internos adyacentes a 1 y 4. Hay al menos \(\theta |D|^3/16\) elecciones. Para una arista cruzada compartida, los cuatro patrones posibles de etiquetas de 2,3 permiten dos elecciones sucesivas, utilizando en cada paso una vecindad interna y un número acotado de vecindades cruzadas casi completas; hay al menos \(n^2/64\) elecciones ordenadas.

Para precisar, las etiquetas de 4,5 en el caso cruzado, cuando las de 0,1 son (i,1-i), son respectivamente ((1-i,1-i)), ((i,i)), ((i,1-i)), ((i,1-i)) para los patrones ((i,i)), ((1-i,1-i)), ((i,1-i)), ((1-i,i)) de 2,3; en el último patrón intercambiamos los dos triángulos raíz antes de utilizar el tercer patrón. Las elecciones de adyacencia siguen el mismo intercambio. Así evitamos exigir dos intersecciones grandes de vecindades internas en la parte no controlada.

Eliminamos las elecciones con vértices repetidos. Las cotas de recorridos que evitan vértices de la Sección C.2 ya cubren esta pérdida. Pasar de elecciones etiquetadas a conjuntos físicos distintos de aristas divide por un factor a lo sumo \(i^{i-4}\) para un soporte de orden \(i\in\{6,8\}\). Con \(|D|\ge n/3\), las cotas anteriores implican

\[\#\{\text{bisagras físicas de orden }i\}
\ge\beta n^{i-4},\qquad \beta=\theta/2^{30}.
\tag{C.4.1}\]

El exponente es \(i-4\) porque las cuatro raíces están fijas. Son descomposiciones positivas de conjuntos simples de aristas, no intercambios con signos contados con multiplicidad.

El lema que convierte certificados firmados en absorbedores se aplica ahora a la Sección C.3. Cada aparición del certificado se separa utilizando vértices nuevos; el emparejamiento de apariciones positivas y negativas de aristas se realiza mediante las bisagras anteriores. Sus dos modos dan descomposiciones de \(A_P\) y \(A_P\cup P\), con \(A_P\cap P=\varnothing\). Para \(|P|\le214\), su tamaño es a lo sumo

\[C=310\cdot126\cdot214.
\tag{C.4.2}\]

Pongamos \(K=41(126\cdot214)+1\). Basta que \(K<n/8\) y \(2^{36}K<\beta n\); estas desigualdades permiten evitar todos los vértices y aristas utilizados previamente en la construcción finita. Por tanto, un umbral de orden terminal es

\[N_t=\lceil4096/\theta\rceil+\lceil8K+1\rceil
+\lceil(2^{36}K+1)/\beta\rceil+1.
\tag{C.4.3}\]

\subsubsection{C.5. Pesos exactos acotados: completar el grafo cruzado y después restringir}\label{c.5.-pesos-exactos-acotados-completar-el-grafo-cruzado-y-despuuxe9s-restringir}

Construimos ahora los pesos necesarios para el teorema de empaquetamiento, independientemente de los certificados enteros anteriores. Sea \(A=D\) la parte buena y \(B\) la otra parte. Completamos temporalmente todas las aristas cruzadas, conservando ambos grafos internos. Sean \(I,J\) sus matrices de incidencia vértice-arista sin signos. Una ponderación triangular se describe mediante dos matrices \(X\in\mathbb R^{E(A)\times B}\), \(Y\in\mathbb R^{A\times E(B)}\). Sus cargas internas son \(X\mathbf1,Y^{\mathsf T}\mathbf1\) y su carga cruzada es

\[IX+YJ^{\mathsf T}.
\tag{C.5.1}\]

Para cada \(u\in A\), promediamos la identidad

\[\tfrac12([uv]+[uw]-[vw])\longmapsto e_u
\tag{C.5.2}\]

sobre todos los triángulos ordenados que pasan por \(u\) en la parte buena. Define una inversa lineal por la derecha \(Q\) de \(I\). La cota inferior de triángulos enraizados de la Sección C.2 da una cota de norma de orden \(1/(\theta n)\).

Damos una solución explícita del problema con todas las aristas cruzadas presentes. Para cargas internas deseadas \(f,g\) y cargas cruzadas \(Z\), el equilibrio significa \(\sum Z=2(\sum f+\sum g)\). Empezamos con \(X_0=f\mathbf1^{\mathsf T}/|B|\) y \(Y_0=\mathbf1g^{\mathsf T}/|A|\). Pongamos

\[R=Z-IX_0-Y_0J^{\mathsf T},\quad r=R\mathbf1,\quad
Y_1=\frac{r\mathbf1^{\mathsf T}}{2|E(B)|},\quad
X_1=Q(R-Y_1J^{\mathsf T}).
\tag{C.5.3}\]

El total de \(R\), y por tanto de \(r\), es cero. Como cada columna de \(J\) suma dos, \(Y_1J^{\mathsf T}\mathbf1=r\). Así, \(X_1\mathbf1=0\) y \(Y_1^{\mathsf T}\mathbf1=0\). La ecuación (C.5.1), con \(X=X_0+X_1,Y=Y_0+Y_1\), tiene exactamente las cargas requeridas.

Pongamos \(\tau=\theta/100\). Las cotas geométricas dan \(|A|,|B|\ge n/3\), \(|E(B)|\ge n^2/40\) y al menos \(2\tau n^2\) triángulos ordenados en cada vértice de \(A\). Al sustituir en (C.5.2)--(C.5.3) obtenemos la siguiente estimación de norma. Si cada coordenada de demanda tiene valor absoluto a lo sumo \(M\), \(X_0,Y_0\) son a lo sumo \(3M/n\); por tanto, \(R\) es a lo sumo \(7M\), \(r\) a lo sumo \(7Mn\), \(Y_1\) a lo sumo \(140M/n\) y \(R-Y_1J^{\mathsf T}\) a lo sumo \(147M\). La norma de coeficientes de \(Q\) es a lo sumo \(3/(2\tau n)\), contando las apariciones de una arista fija en la identidad promediada (C.5.2). La cota uniforme resultante es \((143+147\cdot3/(2\tau))M/n\le512M/(\tau n)\), pues \(0<\tau\le1\). Así,

\[\|Sx\|_\infty\le\frac{512}{\tau n}\|x\|_\infty
\tag{C.5.4}\]

para la inversa con todas las aristas cruzadas presentes. La constante, deliberadamente holgada, absorbe las cotas de grado por fila de ambas matrices de incidencia interna.

Volvemos a \(G\). Una arista superviviente pierde a lo sumo \(3en\) triángulos mixtos al eliminar las aristas cruzadas ausentes: cada tercer vértice descartado debe ser un vecino cruzado ausente de uno de sus extremos. Extendemos una demanda equilibrada de aristas por cero al grafo completado, aplicamos \(S\) y retenemos únicamente los triángulos físicos. El operador de frontera resultante es \(I-E\), con

\[\|E\|_\infty\le3e(512/\tau)\le1/2,
\tag{C.5.5}\]

porque \(307200e\le\theta\). Todas las demandas residuales están equilibradas, pues los triángulos físicos tienen balance cero. Por tanto, la serie geométrica convergente \(S_{\rm phys}\sum_{j\ge0}E^j\) es una inversa exacta sobre ese espacio, acotada por \(1024/(\tau n)\).

Damos a cada triángulo mixto el peso inicial

\[\phi_0(T)=\frac2{(1+4e)n}.
\tag{C.5.6}\]

El conteo directo del tercer vértice mediante (C.1.1)--(C.1.2) da cargas de aristas en \([1-14e,1]\) y pesos iniciales en \([1/n,2/n]\). Por (C.1.3), el defecto \(1-\partial\phi_0\) está equilibrado. Lo reparamos con la inversa física. La corrección tiene valor absoluto a lo sumo \(14e(1024/\tau)/n\le1/n\), pues \(e=\theta/2^{30}\). Los pesos reparados tienen carga exacta uno en cada arista y pertenecen a \([0,4/n]\). La condición necesaria aquí es la no negatividad, no una cota inferior estrictamente positiva. Esto completa la construcción de pesos.

\subsubsection{C.5a. La biblioteca semidensa compatible}\label{c.5a.-la-biblioteca-semidensa-compatible}

El anfitrión semidenso no proporciona consultas internas arbitrarias de dos vecinos. Por eso utilizamos el siguiente resultado de biblioteca compatible, demostrado en H.1--H.4. Sus declaraciones formales se localizan en {[}8, Sección 8{]}. Es distinto del refinador enrutado del Apéndice B.

Sea \(n\ge3\) y fijemos una reserva sin lazos \(X\subseteq E(G)\) con \(\Delta(X)\le\Delta_X\), un entero \(a\ge0\), \(M=2^a\) y una cota entera de grado \(\Delta_X\). Existe una biblioteca abstracta fija con resto \(R\), compatible con las dos etiquetas de las raíces, que refina \(R\cup L\) para todo \(L\subseteq X\) enteramente \(F_0\)-divisible. Sus piezas tienen a lo sumo 214 aristas. Su número de índices es a lo sumo \[I_X=57635(|X|+2M^2).\] Las cargas totales con multiplicidades en vértices antiguos, centros de solicitudes, vértices de fichas y concentradores están acotadas, respectivamente, por \[\begin{gathered}
99d_X(v),\qquad66,\qquad88,\\
330(\lfloor n/M\rfloor+1)\Delta_X+1452M.
\end{gathered}\] Si las cuatro cotas son a lo sumo \(L_{\max}\le q\), donde \(q\) es un entero positivo, un anfitrión con disponibilidad de anti-aristas \(\alpha\) realiza toda esta biblioteca mediante una única aplicación, evitando un grafo de aristas prohibidas de grado máximo \(d\), siempre que \[2\cdot214L_{\max}+\frac{2\cdot214 I_X}{q}
+2(d+\Delta_X+2q)<\alpha n.\] La aplicación fija los vértices antiguos, es inyectiva en cada pieza y en la unión de sus identidades de aristas y deja una carga física con multiplicidades a lo sumo \(2q\). Por tanto, la misma biblioteca insertada sirve para toda elección posterior de \(L\).

La construcción comprime primero el montículo de paridad en \(M\) concentradores. Sus piezas físicas de paridad tienen a lo sumo 22 aristas. La etiqueta intermedia de cada solicitud de dos aristas se elige para completar un triángulo de \(F_0\) con las etiquetas de los extremos. La corrección de carga se construye después, con estas etiquetas ya fijadas; reetiquetar una biblioteca ya corregida no justificaría este paso. Su cota de tamaño es \(9|P|+16\), lo que da \[9\cdot22+16=214.\] Los conteos de la capa de carga y del montículo dan las cuatro cotas de multiplicidad e \(I_X\). Son cotas de toda la biblioteca, no de una descomposición seleccionada. La inserción compartida utiliza únicamente las solicitudes permitidas de dos aristas resultantes. Sus exclusiones comprenden colisiones de piezas, vértices cuya carga con multiplicidades alcanzó su cuota y vecinos prohibidos o ya utilizados. Sus cotas son los tres términos de la desigualdad estricta. Por eso basta la disponibilidad de anti-aristas; nunca se solicitan vecinos comunes para una terna no permitida. Las Secciones H.1--H.4 deducen las cotas de la jerarquía finita, la capa de carga y la inserción compartida; H.5 demuestra el resultado de posiciones robustas utilizado en la capa de carga.

Para la aplicación en C.6 tomemos \(\alpha=1/8\), \(d=0\) y \[\Delta_X=\lceil\kappa^2n\rceil,\quad
\kappa n\le M\le2\kappa n,\quad
q=L_{\max}=\lceil8192\kappa n\rceil.\] Existe una potencia de dos \(M\) en ese intervalo porque \(\kappa n\ge1\). Las condiciones (C.6.1) implican \[\begin{gathered}
I_X\le524288(\kappa n)^2,\\
\Delta_X+2q\le32768\kappa n.
\end{gathered}\] También acotan cada una de las cuatro cargas abstractas por \(q\). En efecto, \(\Delta_X\le2\kappa^2n\), \(q\ge8192\kappa n\) y la expresión de los concentradores es a lo sumo \(4224\kappa n\). Como \(q\le8193\kappa n\), las exclusiones de inserción son a lo sumo \[(428\cdot8193+32768+65536)\kappa n
=3604908\kappa n<n/8,\] donde la última desigualdad se deduce de \(2^{27}\kappa\le1\). Aquí el término central se acotó mediante \(2\cdot214I_X/q\le32768\kappa n\).

Por último, sea \(C\) la cota de \textbf{aristas} del absorbedor privado acotado de C.6. Las condiciones de masa y grado para colocar uno de estos absorbedores por índice de biblioteca son \[\begin{aligned}
\max\{\Delta_X+2q,d_*n/4+C\}+2Cq&\le d_*n,\\
4\{n(\Delta_X+2q)+CI_X\}&\le(d_*n/4)(d_*n).
\end{aligned}\] En este despliegue \(d_*\) denota el parámetro llamado \(d\) en (C.6.1), no el grado prohibido anterior. Las cotas \(2^{19}C\kappa\le d_*\), \(2^{21}C\kappa\le1\), \(2^{22}\kappa\le d_*^2\) y \(C\le d_*n/4\) implican estas desigualdades: el máximo es a lo sumo \(d_*n/2\), \(2Cq\le16386C\kappa n<d_*n/2\) y \(CI_X\le\kappa n^2/4\). Esto verifica la interfaz de colocación privada y da (C.6.3) con el mismo orden de cuantificadores.

\subsubsection{C.6. Robustez y reserva final}\label{c.6.-robustez-y-reserva-final}

La construcción terminal debe seguir siendo utilizable después de evitar vértices y aristas ya utilizados. Partimos de no extremalidad \(2a\) y error \(0\le d\le1/3\). Eliminamos a lo sumo \(\eta n/2\) vértices y un conjunto de aristas de grado máximo a lo sumo \(\delta n\). Si \(\eta,\delta\le a/100\), el anfitrión retenido tiene no extremalidad \(a\); su error de grado y tamaño está acotado por \(2(d+\eta+\delta)\). Para la afirmación sobre la distancia al modelo, extendemos un corte equilibrado de la parte retenida a un corte equilibrado de la parte original sin cambiar el lado de los vértices retenidos. Escribiendo \(D'\) para la parte retenida y \(B_{\rm del}\) para las aristas eliminadas, el coste adicional de edición es a lo sumo \[|D\setminus D'|\,|D|+|B_{\rm del}|
\le(\eta/2+\delta)n^2\le an^2/9\le a|D|^2.\] Aquí \(|D|\ge n/3\), y basta la cota de aristas deliberadamente holgada \(|B_{\rm del}|\le\delta n^2\). La distancia original supera \(2a|D|^2\), por lo que la distancia retenida supera \(a|D|^2\ge a|D'|^2\). Por esta razón se duplica el parámetro original antes de aplicar las construcciones locales.

Fijemos ahora el \(\alpha>0\) original y pongamos \(a=\alpha/2\). Utilizamos (C.2.2), (C.4.2), y tomamos \(N_0\) mayor o igual que los umbrales terminal y de pesos. Pongamos

\[d=\min\{e/16,a/1000,1/1000\},\] \[\kappa=\min\{2^{-27},d/(2^{19}C),(2^{21}C)^{-1},d^2/2^{22}\},
\quad\rho=\kappa^2,\quad\sigma=\rho/2,\quad\gamma=\sigma^2/96.
\tag{C.6.1}\]

Estas elecciones dan \(\kappa^2\le d\), \(2^{19}C\kappa\le d\), \(2^{21}C\kappa\le1\) y \(2^{22}\kappa\le d^2\). Sea \(N_n\) el umbral de empaquetamiento tipado con cota de pesos \(4/n\) y grado máximo residual \(n/64\), y \(N_c\) el umbral de cobertura descendente tipada para cota de pesos \(4/n\) y parámetro de antiaristas \(\gamma\). Tomamos

\[N=2N_0+N_n+N_c+
\left\lceil2^{27}\sigma^{-9}(1/8)^{-6}\right\rceil
+\lceil1/\rho\rceil+\lceil4C/d\rceil+1.
\tag{C.6.2}\]

Utilizamos los pesos exactos de la Sección C.5 para obtener el primer empaquetamiento. El lema de selección de reserva elige un \(X\) descomponible en triángulos de grado máximo \(2\sigma n=\kappa^2n\), con parámetro de antiaristas \(\gamma\). El propio anfitrión tiene parámetro de antiaristas \(1/8\) por sus cotas de grados cruzados e internos. La biblioteca semidensa compatible de C.5a produce piezas divisibles de a lo sumo 214 aristas con una única inserción compartida. Aplicamos sucesivamente la construcción terminal robusta anterior a toda esa biblioteca. Las ecuaciones (C.6.1)--(C.6.2) son precisamente los presupuestos que acotan la carga utilizada en las raíces, permiten evitar los vértices cargados y cubren el coste constante \(C\). El absorbedor resultante \(A\subseteq G-X\) satisface

\[\begin{gathered}
\Delta A\le dn,\qquad A\text{ descomponible},\\
A\cup L\text{ descomponible para todo subgrafo divisible }L\subseteq X.
\end{gathered}
\tag{C.6.3}\]

Su construcción precede al segundo empaquetamiento. Como \(A,X\) son disjuntos y descomponibles, \(G-(A\cup X)\) sigue siendo integralmente divisible. El grado máximo eliminado es a lo sumo \(2dn\). La estimación de robustez sitúa este grafo dentro de las hipótesis de la Sección C.5, por lo que construimos allí \textbf{nuevos pesos exactos}. El teorema de cobertura descendente con esos pesos deja un \(L\subseteq X\) divisible en \(G-A\). Descomponemos \(A\cup L\) mediante (C.6.3). Esto demuestra el teorema con \(\varepsilon=d\) y \(N\) de (C.6.2).

\bigskip

\Needspace{12\baselineskip}\subsection{Apéndice D. División de vértices según las componentes de sus enlaces}\label{apuxe9ndice-d.-divisiuxf3n-de-vuxe9rtices-seguxfan-las-componentes-de-sus-enlaces}

\subsubsection{D.1. División de las etiquetas}\label{d.1.-divisiuxf3n-de-las-etiquetas}

Sea \(F\) una familia de tipos triangulares sobre un conjunto finito \(I\), donde cada etiqueta aparece en al menos un tipo. Para \(r\in I\), su enlace \(L_r\) tiene una arista \(ab\) exactamente cuando \(\{r,a,b\}\in F\). Se permiten lazos en este \emph{grafo de etiquetas}. Representan etiquetas repetidas, no lazos en el anfitrión. Retenemos sólo los vértices activos de \(L_r\) y denotamos por \(\mathcal C_r\) sus componentes conexas. Escribimos \(C_r(s)\) para la componente que contiene una etiqueta activa \(s\). El conjunto de etiquetas divididas es \[I'=\{(r,C):r\in I,\ C\in\mathcal C_r\},\qquad t'=|I'|,\] con proyección \(\pi(r,C)=r\).

Si \(\{a,b,c\}\in F\), la arista \(bc\) de \(L_a\) da \(C_a(b)=C_a(c)\). Por tanto, \[\sigma\{a,b,c\}
=\{(a,C_a(b)),(b,C_b(a)),(c,C_c(a))\}
\tag{D.1.1}\] está bien definido como multiconjunto, independientemente de la enumeración de sus esquinas. Esto sigue siendo válido cuando se repiten etiquetas. Pongamos \(F'=\sigma(F)\). La proyección recupera el tipo original, de modo que \(\sigma\) es inyectiva.

\textbf{Lema D.1.1.} El enlace activo en \((r,C)\) dentro de \(F'\) es isomorfo a \(L_r[C]\). El isomorfismo, incluidos los lazos, es \[s\longmapsto (s,C_s(r)),\qquad s\in C.
\tag{D.1.2}\]

\textbf{Demostración.} La aplicación es inyectiva porque su primera coordenada recupera \(s\). Una arista \(sq\) de \(L_r[C]\) corresponde al tipo \(\{r,s,q\}\). Su tipo dividido contiene \((r,C)\) y aporta la arista entre las dos imágenes de (D.1.2). Recíprocamente, todo tipo que contiene \((r,C)\) es la división de uno de esos tipos: la proyección da \(r,s,q\) y su esquina en \(r\) especifica que \(s,q\) pertenecen a \(C\). Estas afirmaciones también valen para \(s=q\). \(\square\)

En particular, todos los enlaces activos divididos son conexos. La división no elimina la bipartición: el enlace dividido es bipartito exactamente cuando lo es la componente original. Por tanto, esta operación por sí sola no resuelve una familia cuya dificultad restante sea un enlace bipartito.

\subsubsection{D.2. División del anfitrión}\label{d.2.-divisiuxf3n-del-anfitriuxf3n}

Sea \(G\) un grafo simple con aplicación de etiquetas \(p:V\to I\), cuyas aristas tienen todas pares de etiquetas permitidos. Sustituimos \(v\) por una copia \((v,C)\) para cada \(C\in\mathcal C_{p(v)}\), con etiqueta \((p(v),C)\). Enviamos una arista \(uv\) a \[\Phi(uv)=
(u,C_{p(u)}(p(v)))\,(v,C_{p(v)}(p(u))).
\tag{D.2.1}\] Sea \(G'\) el grafo que tiene exactamente estas aristas. No hay aristas entre dos copias del mismo vértice original. La proyección sobre los vértices originales es una inversa por la izquierda de \(\Phi\), por lo que \(\Phi\) es inyectiva y preserva cardinalidades, uniones y disjunción de conjuntos de aristas.

\textbf{Lema D.2.1.} Para cada \(E\subseteq E(G)\), \(\Phi\) da una biyección entre las descomposiciones de \(E\) en triángulos permitidos y las descomposiciones de \(\Phi(E)\) en triángulos de \(F'\).

\textbf{Demostración.} Para un triángulo permitido \(uvw\), sus dos aristas incidentes en \(u\) eligen la misma copia de \(u\), por la identidad de componentes anterior a (D.1.1). Lo mismo ocurre en \(v,w\). Por tanto, su imagen es un triángulo de tipo \(\sigma\{p(u),p(v),p(w)\}\).

Recíprocamente, un triángulo de tipo \(F'\) de \(G'\) se proyecta en tres vértices originales distintos, pues cada par de sus esquinas está unido y las copias de un mismo vértice nunca están unidas. Su tipo proyectado pertenece a \(F\). Cada una de sus aristas es la única imagen bajo \(\Phi\) de su proyección, de modo que levantar el triángulo proyectado recupera el triángulo original. La biyección de aristas transporta entonces las descomposiciones en ambas direcciones. \(\square\)

\textbf{Lema D.2.2.} La compatibilidad entera de \(E\) equivale a la compatibilidad entera de \(\Phi(E)\).

\textbf{Demostración.} Para los conteos globales utilizamos la convención ordenada del Apéndice A. En los pares ordenados permitidos definimos la inclusión de coordenadas \[e_{rs}\longmapsto e_{(r,C_r(s)),(s,C_s(r))}.
\tag{D.2.2}\] Es inyectiva, pues la proyección recupera \(r,s\). El vector de conteos globales de \(\Phi(E)\) es la imagen del de \(E\), y la imagen de la columna de \(\tau\) es la columna de \(\sigma(\tau)\). Para etiquetas repetidas, esta afirmación cuenta los mismos pares de esquinas con sus multiplicidades; en particular, una entrada diagonal sigue contándose dos veces. Como \(F'=\sigma(F)\) y (D.2.2) es inyectiva, la pertenencia a los retículos globales es equivalente.

En un vértice de etiqueta \(r\), el retículo entero local es el retículo de incidencia de su enlace: \[\Lambda_r=\left\langle e_s+e_q:sq\in E(L_r)\right\rangle_{\mathbb Z}
=\bigoplus_{C\in\mathcal C_r}
\left\langle e_s+e_q:sq\in E(L_r[C])\right\rangle_{\mathbb Z}.
\tag{D.2.3}\] La suma directa es literal: las componentes tienen soportes de coordenadas disjuntos. Para un lazo, el generador es \(2e_s\). El perfil en la copia \((v,C)\) es la restricción del perfil en \(v\) a \(C\), transportada mediante (D.1.2). El Lema D.1.1 identifica su retículo con el sumando correspondiente de (D.2.3). La pertenencia del perfil original equivale, por tanto, a la pertenencia en todas las copias. Junto con la afirmación global, esto demuestra el lema. \(\square\)

\subsubsection{D.3. Cotas de disponibilidad y grado}\label{d.3.-cotas-de-disponibilidad-y-grado}

\textbf{Lema D.3.1.} La división de vértices tiene las siguientes propiedades cuantitativas.

\begin{enumerate}
\def\labelenumi{\arabic{enumi}.}
\tightlist
\item
  Una clase dividida \((r,C)\) tiene exactamente \(|V_r|\) vértices.
\item
  Si \(n'=|V(G')|\), entonces \(n\le n'\le t'n\).
\item
  Para todo conjunto de aristas \(X\), \(\Delta(\Phi(X))\le\Delta(X)\).
\item
  Para \(A'\subseteq E(G')\), su conjunto de aristas proyectado satisface \(\Delta(\Phi^{-1}(A'))\le t'\Delta(A')\).
\item
  Una disponibilidad de vecinos comunes compatibles para a lo sumo \(h\) vértices prescritos se transfiere a \(G'\) con el mismo \(h\) y el mismo número \emph{absoluto} de candidatos.
\end{enumerate}

\textbf{Demostración.} Las dos primeras afirmaciones se deducen directamente de la definición de las copias; toda etiqueta original tiene un enlace activo. Las aristas incidentes con una copia concreta forman un subconjunto de las incidentes con su vértice original, lo que demuestra la tercera. Las aristas incidentes con un vértice original se particionan entre sus copias. Tiene a lo sumo \(t'\) copias, lo que demuestra la cuarta.

Para la última afirmación, prescribimos raíces divididas y una etiqueta de destino compatible \((r,C)\). Proyectamos las raíces a \(V\); esto no puede aumentar su número. Un vecino común \(w\in V_r\) de las raíces proyectadas da el candidato \((w,C)\). La compatibilidad de la etiqueta dividida de destino con cada raíz dividida fuerza exactamente las dos clases de copias de (D.2.1), por lo que este candidato es adyacente a todas las raíces divididas. Distintos \(w\) dan candidatos distintos. Por tanto, la disponibilidad de vecinos comunes en \(G\) es una cota inferior para la de \(G'\). \(\square\)

La distinción entre cotas absolutas y relativas es importante. Una clase de tamaño al menos \(\gamma n\) tiene garantizado tamaño al menos \((\gamma/t')n'\), no necesariamente \(\gamma n'\). Una disponibilidad de \(an\) se convierte en al menos \((a/t')n'\). Las cotas de vecinos ausentes en cada clase permitida se transfieren sin cambio en términos absolutos: la clase candidata pertinente en el grafo dividido es una copia de la original y el criterio de adyacencia es el mismo. Ninguna de estas observaciones afirma la transferencia automática de todas las condiciones de regularidad ponderada de patrones utilizadas en otras partes del artículo.

\subsubsection{D.4. Transferencia del absorbedor universal al grafo original}\label{d.4.-transferencia-del-absorbedor-universal-al-grafo-original}

Para precisar, decimos que \(G\) tiene la propiedad de absorbedor universal en \((\rho,\delta)\) si todo \(X\subseteq G\) con \(\Delta(X)\le2\rho|V(G)|\) tiene un grafo \(A\subseteq G-X\), fijado antes de elegir el residuo, con \(\Delta(A)\le2\delta|V(G)|\), tal que \(A\) y \(A\cup L\) se descomponen para todo \(L\subseteq X\) compatible sobre los enteros.

\textbf{Proposición D.4.1.} Si \(G'\) tiene esta propiedad para \(F'\) en \((\rho,\delta/t'^2)\), entonces \(G\) la tiene para \(F\) en \((\rho,\delta)\).

\textbf{Demostración.} Dado \(X\), pongamos \(X'=\Phi(X)\). Por el Lema D.3.1, \(\Delta(X')\le2\rho n\le2\rho n'\). Obtenemos su absorbedor fijo \(A'\subseteq G'-X'\) y lo proyectamos a \(A\). Sus aristas son aristas efectivas de \(G\), y es disjunto de \(X\) por la inyectividad de \(\Phi\). Su grado máximo satisface \[\Delta(A)\le t'\Delta(A')
\le t'\frac{2\delta n'}{t'^2}
\le2\delta n.
\tag{D.4.1}\] Para un \(L\subseteq X\) compatible sobre los enteros, el Lema D.2.2 hace que \(\Phi(L)\) también lo sea. Así, \(A'\cup\Phi(L)\) se descompone. Por el Lema D.2.1, su proyección \(A\cup L\) también se descompone. El mismo argumento proyecta la descomposición de \(A'\). Ninguna elección de \(A\) depende de \(L\). \(\square\)

El cuadrado \(t'^2\) de esta cota segura cubre dos efectos distintos: el posible aumento del número de vértices y la fusión de varias copias en una raíz. No es legítimo acotar el grado de un absorbedor proyectado por el grado de una sola copia.

Finalmente, si el grafo de pares permitidos de la familia dividida tiene varias componentes, cada tipo triangular está contenido en una de ellas. Tanto sus columnas globales como las locales tienen soporte en esa componente. Restringir un conjunto de aristas compatible sobre los enteros a una componente preserva, por tanto, la compatibilidad entera. Los absorbedores fijos construidos en estas componentes pueden unirse: las componentes son disjuntas en vértices, por lo que sus cotas de grado máximo no se suman. Esto permite construirlos por separado en las familias divididas conexas y aplicar después la Proposición D.4.1.

\subsubsection{D.5. Alcance en el argumento de clasificación}\label{d.5.-alcance-en-el-argumento-de-clasificaciuxf3n}

Los casos PAR utilizan esta reducción para pasar de enlaces desconectados a sus componentes conexas y aplicar allí la construcción de absorbedores ya demostrada que corresponda. Los datos finitos deben verificar todavía que cada componente resultante satisface las hipótesis de esa construcción. La prueba anterior no deduce un absorbedor sólo de la conexidad de los enlaces ni afirma que los enlaces bipartitos dejen de serlo tras la división.

\bigskip

\Needspace{12\baselineskip}\subsection{Apéndice E. Paridad, equilibrio, carga y torsión}\label{apuxe9ndice-e.-paridad-equilibrio-carga-y-torsiuxf3n}

\subsubsection{E.1. Datos finitos y lo que debe demostrarse a partir de ellos}\label{e.1.-datos-finitos-y-lo-que-debe-demostrarse-a-partir-de-ellos}

Sea \(K\) una familia componente sobre un conjunto finito de etiquetas. Sea \(c\) la aplicación que identifica las componentes de su grafo de pares activos. Todas las ternas de \(K\) están en una componente. La división del Apéndice D hace que los enlaces raíz pertinentes sean conexos en sus soportes activos. Hay entonces dos posibilidades para el retículo local.

\begin{itemize}
\tightlist
\item
  En un enlace conexo no bipartito, el retículo de incidencia sin signos consta de los vectores enteros de total par.
\item
  En un enlace conexo bipartito de lados \(L_r,R_r\), consta de los vectores enteros cuyas sumas en ambos lados son iguales.
\end{itemize}

Para la segunda afirmación, multiplicamos las filas de un lado por \(-1\). La matriz de incidencia sin signos del enlace se convierte en una matriz de incidencia orientada. Eliminar las hojas de un árbol generador representa todo vector entero de total cero, lo que demuestra la suficiencia. Para la primera, alternar a lo largo de un recorrido da sumas o diferencias de extremos según su paridad; un recorrido cerrado impar da \(2e_a\). Transportarlo a un ancla fija y utilizar el total par demuestra la suficiencia. Son argumentos enteros; la conexidad por sí sola no distinguiría los dos retículos.

Sea \(B_c\) el conjunto de etiquetas con enlaces bipartitos y sean \(\sigma_r(j)\in\{-1,0,1\}\) sus signos de lado, nulos fuera del soporte activo. Para un vértice físico \(v\) definimos su desequilibrio

\[b_v(P)=\sum_j\sigma_{p(v)}(j)d_P(v,V_j).
\tag{E.1.1}\]

El certificado local finito afirma que los perfiles pares soportados, con \(b_v=0\) cuando \(p(v)\in B_c\), pertenecen al retículo local. Puede darse mediante la prueba de recorridos y árboles anterior o mediante fórmulas explícitas de coeficientes enteros. Cada triángulo tiene desequilibrio cero en sus vértices de enlace bipartito.

Para los conteos globales usamos \(u_{ij}=u_{ji}\), contando una arista interna una vez en la diagonal. Su conteo de incidencia ordenada es \((1+\mathbf1_{i=j})u_{ij}\). El certificado global proporciona coeficientes triangulares enteros para todo vector no negativo soportado de conteos que tenga sumas pares de grados por clase, total de aristas divisible por tres, los balances sumados por clase y las condiciones adicionales indicadas. Los tres regímenes BAL son

\[\begin{gathered}
\text{sólo equilibrios};\\
\text{equilibrios y }\chi(P)=0;\\
\text{equilibrios y }g_0q\mid\psi(P).
\end{gathered}
\tag{E.1.2}\]

Aquí \(\chi,\psi\) son pesos enteros especificados sobre pares ordenados de etiquetas. Los certificados escalares correspondientes son fórmulas de soluciones enteras, no comprobaciones de rango real. Sumar (E.1.1) sobre una clase de etiquetas da exactamente

\[\sum_{v:p(v)=r}b_v(P)
=\sum_j\sigma_r(j)(1+\mathbf1_{r=j})u_{rj}.
\tag{E.1.3}\]

Por tanto, los certificados locales y escalares implican conjuntamente la divisibilidad entera completa de una pieza física que satisface (E.1.2). Esta implicación es el último paso de cada construcción, no una hipótesis sobre su resultado.

\subsubsection{E.2. La biblioteca inicial común dentro de una componente}\label{e.2.-la-biblioteca-inicial-comuxfan-dentro-de-una-componente}

Aplicamos el refinador ordinario de árbol binario de la Sección B.3 dentro del conjunto de vértices de la componente. Realizamos cada arista virtual mediante un recorrido permitido de etiquetas de longitud exactamente cuatro, proporcionado por el certificado finito de la componente. Cada arista virtual se pondera por uno más su multiplicidad en toda la biblioteca. La realización voraz evita las aristas utilizadas y los vértices que superan un umbral de carga ponderada \(Q\). Produce piezas pares con número de aristas divisible por tres y tamaño a lo sumo \(S_0=168\), un resto físico \(R\) y las cotas

\[\begin{gathered}
|\mathcal P_0|\le1224\rho n^2,\\
\sum_{P\in\mathcal P_0}d_P(v)\le\Lambda_0,
\end{gathered}\] \[\Lambda_0=4096\lambda n+2(Q+4096\lambda n+1),
\quad
\Delta(R\cup X)\le d_0:=2\rho n+4\lambda n+2(Q+4096\lambda n+1).
\tag{E.2.1}\]

Sustituir una arista por cuatro añade tres aristas, de modo que la construcción por recorridos conserva el invariante módulo tres de este refinador de componentes. Sus alternativas finitas, datos de realización y resto se fijan antes del residuo. Para cada \(L\subseteq X\) divisible, una subfamilia disjunta cubre exactamente \(R\cup L\). La realización ponderada es necesaria: acotar el grado de \(R\) sin acotar todas las apariciones de la biblioteca no justificaría los absorbedores privados posteriores.

Para las componentes PAR cuyos retículos locales sólo imponen paridad, hay tres certificados globales finitos en la biblioteca existente: basta la divisibilidad ordinaria; queda una carga entera; o queda una congruencia cíclica. En el primer caso, (E.2.1) ya es una biblioteca divisible. Para una carga, usamos el registro fijo de la Sección B.4 con fichas físicas positivas y negativas; para una congruencia, la cadena cíclica de la Sección E.5. Las plantillas finitas tienen a lo sumo seis aristas. En la normalización de una carga utilizada por las componentes PAR, sus cargas son \(\pm12\), los pesos de pares tienen valor absoluto a lo sumo cuatro y las piezas resultantes tienen a lo sumo 840 aristas. El certificado finito de retículo y la construcción terminal de las Secciones A.4--A.5 completan entonces cada componente. Éste es el contenido matemático reutilizado por las tablas de división PAR; la tabla no supone una conclusión de descomposición.

\subsubsection{E.3. Fichas de equilibrio enraizadas y una componente concreta}\label{e.3.-fichas-de-equilibrio-enraizadas-y-una-componente-concreta}

Para cada etiqueta raíz \(r\) de enlace bipartito, los datos BAL especifican dos grafos finitos enraizados \(T_r^+,T_r^-\), con a lo sumo \(g\) aristas y \(m+1\) vértices. Son sin lazos, activos, pares y están contenidos en la componente. Su raíz tiene etiqueta \(r\) y

\[\begin{gathered}
b_{\rm root}(T_r^\pm)=\pm2,\\
b_y(T_r^\pm)=0\quad(y\ne\text{raíz}).
\end{gathered}
\tag{E.3.1}\]

Cada vértice de (E.3.1) es físico; las dos fichas son grafos con aristas positivas. También satisfacen \(M\mid e(T_r^\pm)\), donde \(M\in\{1,3\}\). Si \(M=1\), un potencial local separado demuestra que las piezas pares equilibradas tienen tamaño divisible por tres. Esta alternativa es necesaria: no toda ficha raíz válida tiene tamaño divisible por tres.

Por ejemplo, tomemos \(K=\{001\}\). En una raíz de etiqueta 0, \(b=d(\cdot,V_0)-d(\cdot,V_1)\); la etiqueta 1 no tiene condición de equilibrio. La ficha positiva es el ciclo (0,1,2,3,0), con etiquetas sucesivas (0,0,1,0). Su raíz tiene desequilibrio 2 y los demás vértices de etiqueta 0 tienen cada uno un vecino de cada etiqueta. La ficha negativa es el ciclo (0,1,2,3,4,0), etiquetado (0,1,0,0,1). Su raíz tiene desequilibrio \(-2\) y los demás vértices de etiqueta 0 están equilibrados. Ambas son activas y pares. Sus tamaños son cuatro y cinco, por lo que usamos \(M=1\). Una vez equilibrados todos los vértices de etiqueta 0,

\[e_{01}=2e_{00},\qquad e(P)=3e_{00}.
\tag{E.3.2}\]

Esto demuestra la divisibilidad del tamaño requerida y proporciona el coeficiente triangular global \(e_{00}\). En la etiqueta 0, el coeficiente local es \(d_0=d_1\); en la etiqueta 1 es \(d_0/2\). Ésta es la componente completa que subyace a cada una de las dos componentes del ejemplo \(F_{\rm BAL000}=\{001,223\}\).

En general, el mecanismo de potencial finito tiene la misma forma. Si un peso de par es congruente módulo \(q\) con la suma de dos potenciales de extremos y cada potencial de vértice se anula módulo \(q\) en perfiles pares equilibrados, sumar sobre los vértices demuestra la divisibilidad global módulo \(q\). Para la divisibilidad de la carga por seis se utiliza un potencial para \(\chi\) módulo dos y otro para \(\chi-1\) módulo tres, junto con

\[(\chi-1)(P)=\chi(P)-2e(P).
\tag{E.3.3}\]

Las ecuaciones (E.3.2)--(E.3.3) explican por qué las tablas de componentes incluyen tanto fórmulas de coeficientes locales como potenciales; un núcleo racional no demostraría las congruencias requeridas.

\subsubsection{E.4. Equilibrar todos los vértices del anfitrión sin pérdida exponencial}\label{e.4.-equilibrar-todos-los-vuxe9rtices-del-anfitriuxf3n-sin-puxe9rdida-exponencial}

Procesamos los vértices cuyas etiquetas pertenecen a \(B_c\). En un vértice actual \(x\), cada pieza de la biblioteca tiene desequilibrio par: su grado es par y cada signo activo es impar. Su demanda entera es \(c_i=b_x(P_i)/2\). Para un residuo divisible \(L\), \(b_x(L)=0\); por tanto, toda familia seleccionada tiene la misma demanda total, a saber, \(b_x(R)/2\). Utilizamos el registro fijo de la Sección B.4, realizando sus fichas positivas y negativas mediante copias enraizadas de \(T_{p(x)}^\pm\) en \(x\). Cada pieza nueva tiene desequilibrio cero en \(x\). Los vértices tratados anteriormente siguen equilibrados por (E.3.1), y los desequilibrios pendientes no cambian en el descendiente de una pieza antigua; las piezas formadas sólo por fichas no tienen desequilibrio pendiente.

El invariante esencial de tamaño es

\[\mathcal V(P)=e(P)+g\sum_{y\ \mathrm{pending}}|b_y(P)/2|.
\tag{E.4.1}\]

Una pieza antigua necesita a lo sumo \(|b_x(P)/2|\) fichas, que añaden a lo sumo esa cantidad multiplicada por \(g\) aristas, mientras que exactamente ese término desaparece de la suma. Por tanto, el potencial de su descendiente no aumenta. Una pieza de liberación o emparejamiento consta de dos fichas y tiene potencial a lo sumo \(2g\). Inicialmente, \(\sum_y|b_y(P)/2|\le e(P)\le S_0\), de modo que todas las piezas tienen en todo momento tamaño a lo sumo \(S_0(1+g)\), que también domina \(2g\). Esta cota no depende del número de vértices procesados.

El mismo argumento de herencia da una cota uniforme de las demandas futuras. Sea \(N\ge9\Lambda_0\). En cualquier vértice pendiente \(y\),

\[9\sum_i|b_y(P_i)/2|\le N.
\tag{E.4.2}\]

Inicialmente esto se deduce de \(|b_y(P_i)|\le d_{P_i}(y)\); después, los únicos valores no nulos son los heredados por los descendientes únicos de las piezas antiguas. Por tanto, una capa utiliza a lo sumo \(N\) fichas nuevas. Si \(U\) es la unión de las fichas colocadas hasta ese momento, las cotas inductivas son

\[e(U)\le gN\cdot\#\{\text{raíces tratadas}\},\quad
d_U(v)\le\theta+g+gN\mathbf1_{\{v\text{ tratado}\}},
\tag{E.4.3}\] \[\begin{gathered}
\sum_{P\in\mathcal P}d_P(v)\le\Lambda_0+44d_U(v),\\
|\mathcal P|\le|\mathcal P_0|+11N\cdot\#\{\text{raíces tratadas}\}.
\end{gathered}
\tag{E.4.4}\]

El coeficiente 44 se suma sobre los conjuntos de fichas físicamente disjuntos; no se multiplica en cada capa. La cota 11 cubre los a lo sumo \(97N/9\) índices nuevos del registro.

Para realizar una capa, prohibimos los vértices que no son raíz y cuya carga actual de fichas supera \(\theta\); por (E.4.3), hay a lo sumo \(2gNn/\theta\) de ellos. También prohibimos los vértices utilizados en el lote actual y los vecinos por aristas ya utilizadas. Una desigualdad suficiente de disponibilidad protegida es

\[(m+1)+\frac{2gNn}{\theta}+(m+1)N
+(m+1)(d_0+\theta+g+2gN)<\omega n.
\tag{E.4.5}\]

Los cuatro términos cubren, respectivamente, la distinción de vértices dentro de una plantilla, los vértices cargados, la disjunción de vértices del lote actual fuera de su raíz fija y la evitación de aristas. En cada paso, un candidato restante permite insertar el siguiente vértice de la plantilla. Un vértice que no es raíz tiene a lo sumo un exceso de carga de \(g\), mientras que la raíz fija recibe a lo sumo \(gN\); esto demuestra (E.4.3) para la capa siguiente. Demuestra la existencia de la disponibilidad requerida, en lugar de suponer que todas las fichas raíz existen simultáneamente.

Tras todas las capas, cada pieza es equilibrada y par. Conserva la divisibilidad del tamaño por \(M\); la implicación del potencial finito da divisibilidad por tres. En el régimen de sólo equilibrio, el certificado escalar de la Sección E.1 demuestra ahora la divisibilidad entera completa. El grado máximo, la carga total de la biblioteca y el número de piezas son los de (E.4.3)--(E.4.4).

\subsubsection{E.5. Carga global adicional y torsión}\label{e.5.-carga-global-adicional-y-torsiuxf3n}

\paragraph{Una carga global}\label{una-carga-global}

Si las fichas raíz son neutras para \(\chi\), su adición no cambia la carga de un descendiente de una pieza antigua; las piezas formadas sólo por fichas tienen carga cero. Si, en cambio, sus cargas son \(\pm c_r\), una pieza formada sólo por fichas sigue teniendo carga cero, pues utiliza el mismo número de fichas de ambos signos. Una pieza de carga cero y sin desequilibrio pendiente no recibe fichas. Por tanto, el número de piezas que llevan carga o un desequilibrio pendiente nunca aumenta. Está acotado por el tamaño inicial de la biblioteca, no por el número mucho mayor de (E.4.4).

Después de las capas de vértices, sea \(u>0\) la unidad de carga especificada por el certificado de potencial finito. La carga de cada pieza es divisible por \(u\). Los datos de la componente incluyen dos plantillas pares equilibradas de a lo sumo seis aristas y cargas \(\pm u\). Utilizamos el mismo registro fijo con demanda \(\chi(P)/u\). Las fichas nuevas no cambian los balances de los vértices. La carga se anula en cada pieza nueva. En la normalización \(|\chi_{ij}|\le4\),

\[\begin{gathered}
|\chi(P)|\le8e(P),\\
K_\chi=8S_0(1+g)|\mathcal P_0|
\end{gathered}
\tag{E.5.1}\]

acota la demanda absoluta total, también para la variante con raíces cargadas. El registro adicional utiliza a lo sumo \(9K_\chi\) fichas, \(54K_\chi\) aristas y \(97K_\chi\) índices. El tamaño de pieza es a lo sumo \(49S_0(1+g)\): seis aristas por ficha multiplicadas por \(8e(P)\), más la pieza antigua. Aplicamos ahora el certificado escalar con carga de (E.1.2).

\paragraph{Una congruencia cíclica}\label{una-congruencia-cuxedclica}

Supongamos que cada pieza preliminar par y de tamaño divisible por tres tiene carga divisible por \(g_0>0\) y que la obstrucción es \(\psi(P)/g_0\pmod q\). Los datos finitos incluyen una ficha par de seis aristas, tamaño divisible por tres y valor uno módulo \(q\). Sus copias físicas no tienen por qué estar equilibradas inicialmente. Aplicamos esta capa \textbf{antes} de las capas de equilibrio de vértices.

Enumeramos la biblioteca antigua en \(s\) filas. Entre filas consecutivas colocamos \(q\) fichas privadas de valor uno. Una alternativa en la fila \(i\) tiene un bit que especifica si utiliza la pieza antigua, un acarreo entrante \(b<q\) y un acarreo saliente \(a=(b+x_i)\bmod q\), donde \(x_i\) es su valor antiguo seleccionado. Salvo en la última fila, recibe \(b\) fichas de la conexión anterior y \(q-a\) de la siguiente. El valor total es \(x_i+b+q-a=0\pmod q\). La siguiente fila recibe las \(a\) fichas restantes, de modo que cada conexión se particiona exactamente. El primer acarreo es cero. El último es el valor fijo \(r\) del resto antiguo. En la última fila no hay conexión siguiente: su contribución saliente se sustituye por \((q-r)\bmod q\) fichas finales fijas. Cancelan \(x_i+b\equiv r\pmod q\). Los residuos divisibles tienen valor cero, por lo que todos tienen ese mismo último acarreo. La selección vacía fija \(r\) antes de cualquier residuo posterior.

Hay a lo sumo \(2qs\) alternativas y \(qs\) fichas en el resto fijo. Una pieza recibe a lo sumo \(3q\) fichas, por lo que su tamaño es a lo sumo \(S_0+18q\). Su carga de biblioteca es a lo sumo \(q\Lambda_0\) más \(4q\) veces la carga de fichas. Equilibramos ahora los vértices con fichas cuya carga \(\psi\) es exactamente cero. La congruencia cero se conserva en todo momento; el argumento de potencial de la Sección E.4 se aplica con tamaño inicial \(S_0+18q\). Finalmente, el certificado escalar de torsión da divisibilidad entera. La neutralidad exacta de estas fichas posteriores es necesaria; saber sólo que sus desequilibrios en las raíces son correctos no preservaría la congruencia.

\subsubsection{E.6. Elecciones cuantitativas y ensamblaje terminal}\label{e.6.-elecciones-cuantitativas-y-ensamblaje-terminal}

Sea \(C_t\) la constante del mayor certificado terminal para las piezas acotadas anteriores, de modo que un absorbedor privado utiliza a lo sumo \(30C_t\) aristas, como en el Apéndice A. Estas constantes quedan fijadas por los datos finitos de la componente antes de elegir el anfitrión. Elegimos un exponente \(L\) suficientemente grande para que

\[2^{40}(30C_t+1)(g+1)(m+2)(t+1)\le2^L.
\tag{E.6.1}\]

Para las componentes de torsión, donde las tablas congeladas tienen \(q\le12\), utilizamos \(2^{46}\) en lugar de \(2^{40}\). Imponemos

\[0<\omega\le1,\quad 0<\delta,\quad2^{63}\delta\le\omega,\quad
\lambda=\omega^2\delta/2^L,\quad
0<2^{24}\rho\le\lambda^2,\quad\rho n\ge1.
\tag{E.6.2}\]

Para el caso de sólo equilibrio elegimos

\[Q=2^{15}\lambda n/\omega,\quad
\Lambda_0=4096\lambda n+2(Q+4096\lambda n+1),\quad
N=\lceil9\Lambda_0\rceil,\quad
\theta=2^6(g+1)N/\omega.
\tag{E.6.3}\]

Para verificar que estas elecciones bastan, medimos los costes en \(b=\lambda n/\omega^2\). Entonces

\[\delta n=2^Lb,\quad \omega n\ge2^{63}\delta n,\quad
Q\le2^{15}b,\quad \Lambda_0\le2^{17}b,\quad
N/\omega\le2^{21}b,\quad\theta\le2^{27}(g+1)b,
\tag{E.6.4}\]

La estimación intermedia más fuerte \(\Lambda_0/\omega\le2^{17}b\) demuestra la cota para \(N/\omega\), incluido su techo. Esto es necesario al dividir por \(\omega\) en la definición de \(\theta\). Además, \(2gNn/\theta\le\omega n/32\). También \(\lambda n\ge2^{12}\), por lo que los techos y las plantillas de tamaño constante quedan absorbidos en estas cotas. Las ecuaciones (E.6.1)--(E.6.4) implican (E.4.5) y las tres desigualdades de absorbedores privados de la Sección A.6: el presupuesto de vecinos prohibidos es menor que la mitad de la disponibilidad, la carga de raíces más la auxiliar es a lo sumo \(2\delta n\) y la masa total de aristas permite evitar los vértices cargados. Explícitamente, con \(D_U=\theta+g+gN\), esas tres desigualdades son \[2S_0(1+g)+5t+9C_t+8\delta n<\omega n/2,\] \[\max\{d_0+D_U,\delta n+30C_t\}
+30C_t(\Lambda_0+44D_U)\le2\delta n,\] \[4\{2\rho n^2+16\lambda n^2+gNn
+30C_t(1224\rho n^2+11Nn)\}\le\delta\omega n^2.
\tag{E.6.5}\] En particular, la cota de grado cuenta todas las apariciones de la biblioteca en una raíz; usar sólo el grado de su unión no la demostraría.

La variante de raíces cargadas utiliza la cota (E.5.1), una densidad separada de fichas \(\rho_T=2^{11}(g+1)\rho\) y \(\lambda_T=2^{26}(g+1)\delta/2^L\); la variante de carga neutra es su caso más sencillo. Tomamos \(K_\chi\) como el techo de \(8S_0(1+g)1224\rho n^2\), que acota (E.5.1). Estas escalas adicionales se eligen a partir de los datos de la componente ya fijados. En el caso de torsión tomamos \(\lambda_T=\Lambda_0/n\), \(\rho_T=2^{11}\rho\) y \(N=\lceil45q\Lambda_0\rceil\). Valen las mismas estimaciones con \(b'=64b\) y \(2^{L'}=2^L/64\). Esto explica el factor adicional \(2^6\) en (E.6.1). El documento complementario {[}8, Sección 7, mapa del Apéndice E{]} identifica los módulos fuente que contienen las desigualdades adicionales explícitas y sus sustituciones.

Insertamos un absorbedor privado para cada pieza final de la biblioteca, evitando \(X\), el resto fijo y todos los absorbedores privados anteriores. Existen por las Secciones A.4--A.6 y los presupuestos recién comprobados. Para cualquier residuo divisible posterior \(L\), elegimos la subfamilia refinadora, utilizamos el modo de absorción en sus índices y el modo vacío en los demás. El resultado es un absorbedor universal de la componente de grado máximo a lo sumo \(2\delta n\). Esta construcción también trata una componente vacía tomando un absorbedor vacío.

\subsubsection{E.7. Transporte, certificados finitos y conclusión para las familias}\label{e.7.-transporte-certificados-finitos-y-conclusiuxf3n-para-las-familias}

PAR no afirma que una componente sea resoluble por llevar el nombre PAR. Su certificado consta de aplicaciones explícitas de división, una lista finita de ternas reducidas, una aplicación de componentes e isomorfismos de componentes con las construcciones demostradas en las Secciones E.2--E.6. Por ejemplo,

\[F_{\rm PAR000}=\{000,001,002,111,113\}\]

se divide en las dos componentes disjuntas

\[\{000,001,003\}\quad\text{y}\quad\{222,224\}.
\tag{E.7.1}\]

La primera es un reetiquetado de la componente monofamiliar registrada FM01; la segunda es la componente de dos etiquetas KE2. Un vértice de etiqueta 1 del grafo original tiene dos copias: una para su componente del enlace que contiene la etiqueta 0 y otra para la que contiene las etiquetas 1,3. La aplicación explícita sobre pares determina qué copia utiliza cada arista. Así, cada arista tiene exactamente un levantamiento y cada triángulo permitido tiene una terna levantada coherente. El Apéndice D demuestra la preservación de la divisibilidad entera, la disponibilidad protegida y la proyección de descomposiciones.

Los absorbedores de componentes reducidas distintas son disjuntos en vértices en el anfitrión dividido. Sus grados máximos no se suman allí. Al proyectar, el Apéndice D contempla la pérdida \(t'^2\) indicada, donde \(t'\) es el número de etiquetas divididas. Esta pérdida se incorpora reduciendo el presupuesto por componente antes de construir los absorbedores. No debe omitirse confundiendo un vértice dividido con el original que contiene todas sus copias.

Las tablas BAL especifican, para cada componente reducida, los signos efectivos, las fórmulas de coeficientes, los potenciales, las plantillas raíz y las plantillas opcionales de carga o torsión que satisfacen las Secciones E.1, E.3 y E.5. Contienen más información que los nombres de las familias. Sus identidades finitas forman parte de las tablas de certificados adjuntas {[}8{]}, y los argumentos universales de componentes se han demostrado arriba. La siguiente tabla registra cómo interviene esa evidencia finita en el teorema.

{\begin{longtable}[]{@{}
  >{\raggedright\arraybackslash}p{(\linewidth - 4\tabcolsep) * \real{0.3333}}
  >{\raggedright\arraybackslash}p{(\linewidth - 4\tabcolsep) * \real{0.3333}}
  >{\raggedright\arraybackslash}p{(\linewidth - 4\tabcolsep) * \real{0.3333}}@{}}
\toprule\noalign{}
\begin{minipage}[b]{\linewidth}\raggedright
Ruta
\end{minipage} & \begin{minipage}[b]{\linewidth}\raggedright
Evidencia finita
\end{minipage} & \begin{minipage}[b]{\linewidth}\raggedright
Argumento universal
\end{minipage} \\
\midrule\noalign{}
\endhead
\bottomrule\noalign{}
\endlastfoot
PAR & Aplicaciones de división, isomorfismos de componentes y certificados de componentes base {[}8{]} & Apéndice D y componentes de paridad, carga o torsión de la Sección E.2 \\
BAL sin carga global adicional & Soluciones de coeficientes locales, soluciones escalares, fichas raíz equilibradas y potenciales {[}8{]} & Secciones E.1, E.3, E.4, E.6 \\
BAL con una carga & Los datos anteriores más la unidad de carga, la identidad de anulación, las cargas de fichas raíz y las fichas globales equilibradas & Secciones E.1, E.3--E.6, incluida la herencia de piezas activas \\
BAL con torsión & Módulo entero, ficha de valor uno, fichas de equilibrio neutras y solución escalar bajo la congruencia & Secciones E.1, E.3--E.6, con la capa cíclica primero \\
\end{longtable}
}

El conjunto congelado tiene 271 representantes PAR y 357 representantes BAL en estas rutas (301 sin carga adicional, 52 con una carga y cuatro con torsión), además de 120 representantes cubiertos previamente. Los diecisiete representantes restantes con leyes de conservación utilizan la construcción del Apéndice F. La enumeración finita completa, el transporte por permutaciones y el teorema de cobertura son objetos demostrados por separado en la versión congelada. La cobertura exhaustiva utiliza el certificado finito separado y el argumento de transporte de la Sección 9; no se deduce simplemente sumando estos conteos.

Para cada componente, las constantes terminales y de la inversa real son finitas. Tomar máximos sobre sus componentes, que son finitas, y después sobre las tablas finitas de familias fija todos los parámetros antes de elegir un anfitrión. Aplicamos el ensamblaje de inversa real, reserva y cobertura descendente de la Sección B.6 (con las escalas flexibles (E.6.2) en el caso BAL). Convierte estos absorbedores universales y pesos regulares en una descomposición exacta. Reetiquetar de vuelta da el resultado afirmado para anfitriones típicos de cada familia PAR/BAL. Junto con el Apéndice F y el teorema del censo, esto proporciona la parte constructiva del Teorema D.La inversa real de la Proposición A.2.1 se aplica sobre el anfitrión original de cuatro etiquetas; dividir etiquetas para el absorbedor de componentes no requiere aplicar esa proposición a más de cuatro etiquetas.

\bigskip

\subsubsection{E.8. Los casos conexos y la interfaz con posiciones adicionales}\label{e.8.-los-casos-conexos-y-la-interfaz-con-posiciones-adicionales}

Los 120 representantes denominados casos anteriores tienen enlaces activos de raíz conexos y no bipartitos. De ellos, 93 tienen enlaces de tres pasos y están cubiertos directamente por el Teorema A. Los otros 27 requieren sus certificados finitos de componentes de paridad; la conexión y la no bipartición de los enlaces no establecen por sí solas la condición de tres pasos. La misma ruta de certificados cubre los 271 representantes PAR después de la separación de componentes del Apéndice D. Por tanto, los 27 casos no se deducen del Teorema A ni del total numérico del censo.

Describimos la interfaz intermedia para explicitar esta distinción. Una plantilla con posiciones adicionales tiene \(m+3\) posiciones etiquetadas en cada una de las \(t\) clases, de modo que su orden es \(q=t(m+3)\). Para un tipo triangular se distinguen tres de esas posiciones; una inserción debe preservar todas las etiquetas y ser inyectiva, pero sólo se exige que los tres pares distinguidos formen ese triángulo. La regularidad de estas plantillas asigna a las inserciones pesos en \([\omega n^{2-q},\omega^{-1}n^{2-q}]\), con cargas de aristas dentro de una tolerancia \(c\) de uno. La extensión de plantillas exige al menos \(\omega n^{k-|S|}\) extensiones de toda aplicación inyectiva etiquetada de raíces prescrita sobre \(S\), en cada plantilla de orden \(k\le K\), imponiendo los pares permitidos con un extremo nuevo. No exige que aparezcan aristas en pares prohibidos.

La propiedad de descomposición con posiciones adicionales utiliza \[\begin{gathered}
h=2^{50q^3},\qquad d_q=2^{-1000q^5},\\
c=\omega^{h^{20}},\qquad K=qh.
\end{gathered}\] Para cada \(m\) elige una cota superior positiva de densidad y un umbral de orden antes del anfitrión. Da una descomposición para anfitriones enteramente divisibles con aristas permitidas, clases de tamaño al menos \(n/h\), regularidad y extensión de plantillas como las anteriores, siempre que \(n^{-d_q}<\omega\) y se cumpla la cota superior de densidad. Es una interfaz que debe demostrarse para cada familia, no una hipótesis de que todas las familias la satisfacen.

Los certificados finitos de componentes proporcionan un absorbedor universal y un inverso en el anfitrión completo. La reducción con posiciones adicionales suma los pesos de todas las inserciones de una misma copia triangular física, obteniendo pesos entre \(\omega^2/n\) y \(27/(\omega n)\) con la misma tolerancia de carga por arista. La extensión de plantillas proporciona disponibilidad protegida tomando una plantilla con un solo vértice nuevo. La condición de compañero prohibido da el margen de grado. Los argumentos de reserva, corrección y cobertura hacia la reserva demuestran entonces la propiedad con posiciones adicionales. Son resultados y reducciones separados en el mapa de fuentes de {[}8, Sección 8{]}; ninguno supone la conclusión de descomposición bajo otro nombre.

Recíprocamente, esta propiedad implica (2.7). Dada la cota inferior \(\gamma n\) para las clases, tomamos \(m=\lceil1/\gamma\rceil\). Distribuimos el peso de cada triángulo físico entre sus extensiones inyectivas con posiciones adicionales. Para \(n\) suficientemente grande hay suficientes vértices no utilizados en cada clase; los conteos de extensiones acotan cada peso resultante y conservan su carga total en las aristas. Aplicar sucesivamente la condición de vecinos comunes da la extensión de plantillas. Reduciendo la densidad y la tolerancia y aumentando el umbral de orden, todo antes del anfitrión, se obtienen las hipótesis requeridas. La descomposición resultante está formada por los triángulos físicos originales. Este es el puente utilizado para los 120 casos conexos y las familias PAR, cuyos certificados finitos proporcionan los resultados específicos de cada familia. Las tablas finitas completas siguen siendo evidencia necesaria; un ejemplo representativo no las sustituye.

\Needspace{12\baselineskip}\subsection{Apéndice F. Un catálogo fijo para las familias con leyes de conservación}\label{apuxe9ndice-f.-un-catuxe1logo-fijo-para-las-familias-con-leyes-de-conservaciuxf3n}

En todo este apéndice, un digrafo cíclico tiene tres clases de vértices y cada arco va de la clase i a la clase i+1 módulo 3. Los digrafos no tienen arcos repetidos. El grado d(v) significa \textbf{grado de entrada más grado de salida}. Una circulación con signos es un vector entero sobre los arcos cuyas sumas entrante y saliente son iguales en cada vértice. Una pieza positiva es un conjunto no vacío de arcos equilibrado en ese mismo sentido. No tiene por qué ser un único ciclo. Su tamaño es su número de arcos.

Para un catálogo finito C de piezas, su carga en v es

\[\operatorname{load}_{\mathcal C}(v)=\sum_{P\in\mathcal C}d_P(v).\]

Esto cuenta el catálogo completo, no sólo las piezas seleccionadas para un grafo residual concreto. Todas las reservas y catálogos que siguen se eligen antes del residuo L. Sólo la partición elegida del catálogo depende de L.

\subsubsection{F.1. Redondeo entero y una inversa bilateral}\label{f.1.-redondeo-entero-y-una-inversa-bilateral}

Orientamos arbitrariamente las aristas de un grafo simple finito y denotamos por B su matriz de incidencia con signos, con una columna \(e_{v}-e_{u}\) para una arista orientada de u a v.

\textbf{Redondeo de flujos.} Si Bf es entero, existe un vector entero z con \(Bz=Bf\) y \(\lfloor f_e\rfloor\le z_e\le\lceil f_e\rceil\) para cada arista e. En particular, \(|z_{e}-f_{e}|<1\).

Para verlo, consideramos las aristas en las que f no es entero. Ningún vértice tiene exactamente una arista incidente de ese tipo: su ecuación de incidencia tendría un lado izquierdo no entero y un lado derecho entero. Si queda alguna de esas aristas, su grafo subyacente contiene, por tanto, un ciclo. Añadimos un múltiplo con signo de su vector de incidencia cíclica. Esto preserva Bf. Aumentamos el múltiplo hasta que una arista alcance por primera vez un extremo entero de su intervalo permitido. Ninguna otra coordenada sale de su intervalo y desaparece al menos una coordenada no entera. Las coordenadas anteriormente enteras no cambian. La iteración termina y demuestra la afirmación.

Supongamos que G satisface la desigualdad de corte

\[|\partial_G S|\ge h\min\{|S|,|V(G)\setminus S|\},\qquad h>0.\tag{F.1}\]

Todo vector real g de suma cero con \(|g_{v}|\le M\) tiene una preimagen real bajo B de norma supremo a lo sumo M/h. En efecto,

\[\left|\sum_{v\in S}g_v\right|\le M\min\{|S|,|V(G)\setminus S|\}\le(M/h)|\partial_G S|.\]

Estas desigualdades de corte implican la afirmación sobre flujos. Para completar la prueba, la imagen de la caja \([-M/h,M/h]^E\) bajo B es un conjunto convexo compacto. Si g estuviera fuera, un potencial separador p satisfaría

\[\langle p,g\rangle>(M/h)\sum_{uv\in E}|p_u-p_v|.\]

Restamos una constante a p, lo que no cambia el lado izquierdo, y expresamos p como suma de indicadores de sus conjuntos de nivel anidados multiplicados por diferencias consecutivas de nivel. Aplicamos la desigualdad de corte a cada término y sumamos. Obtenemos la desigualdad opuesta, una contradicción. Si g es entero, podemos redondear como antes sin cambiar su divergencia.

\textbf{Inversa bilateral.} Sean A y B grafos conexos con constantes de expansión \(h_{A},h_{B}\). Si X es una matriz entera con sumas de filas y columnas nulas y \(|X_{ab}|\le M\), entonces

\[B_A Z B_B^{\mathsf T}=X,\qquad |Z_{ef}|<\frac{2M}{h_Ah_B}+2.\tag{F.2}\]

para una matriz entera Z.

Levantamos primero cada columna de X a través de \(B_{A}\), obteniendo columnas \(u_{b}\) acotadas por \(M/h_{A}\). Pongamos \(s_b=u_b-\operatorname{mean}_{b\prime}u_{b\prime}\). Entonces \(B_{A}\) \(s_{b}=X_{b}\), pues las filas de X suman cero; las filas de la matriz s también suman cero y \(|s_{eb}|\le 2M/h_{A}\). Levantamos cada fila a través de \(B_{B}\). Esto da una matriz real C con

\[B_A C B_B^{\mathsf T}=X,\qquad |C_{ef}|\le\frac{2M}{h_Ah_B}.\tag{F.3}\]

También existe una preimagen entera \(C_{0}\), sin afirmar una cota de norma: expresamos X como suma de \(X_{ab}(e_a-e_{a_0})(e_b-e_{b_0})^{\mathsf T}\), para raíces fijas \(a_0,b_0\), y levantamos las dos diferencias mediante caminos en A y B.

Redondeamos C preservando su imagen bilateral. Elegimos un árbol generador de A. Sea P la aplicación que envía un vector de aristas al único vector soportado en ese árbol con la misma incidencia. Así, \(B_AP=B_A\) y P se anula en \(\ker B_A\). Pongamos

\[D=C-C_0,\qquad V=PD,\qquad U=D-V.\]

Entonces \(B_AU=0\). Además, \(B_A D B_B^{\mathsf T}=0\), por lo que \(P(DB_B^{\mathsf T})=0\) y, en consecuencia, \(VB_B^{\mathsf T}=0\). Redondeamos cada columna de U como flujo de divergencia cero en A y cada fila de V como flujo de divergencia cero en B. Las matrices enteras resultantes \(U_{Z},V_{Z}\) conservan esas identidades y cambian cada entrada en menos de uno. Por tanto, \(Z=C_{0}+U_{Z}+V_{Z}\) demuestra (F.2). No se utiliza ninguna cota sobre P. En particular, este argumento no supone que un producto tensorial de matrices de incidencia sea totalmente unimodular.

\subsubsection{F.2. Un expansor de grado acotado para cada tamaño}\label{f.2.-un-expansor-de-grado-acotado-para-cada-tamauxf1o}

Para cada \(n\ge 1\) existe un grafo simple G sobre n vértices con grado máximo a lo sumo ocho y expansión \(h=1/20\). El caso de un vértice es inmediato. Para n mayor, tomamos cuatro permutaciones uniformes independientes y formamos la unión no dirigida de sus aristas de permutación, descartando lazos y repeticiones.

Damos el argumento de conteo con su estimación numérica. Fijemos un conjunto S de tamaño \(1\le s\le n/2\) y pongamos \(r=\lfloor (s-1)/20\rfloor\). Una permutación con a lo sumo r imágenes que salen de S envía S a \(S\cup R\) para algún conjunto R de r elementos fuera de S. Su probabilidad es a lo sumo

\[\binom nr\binom{s+r}s\big/\binom ns.\]

Por tanto, la probabilidad de que algún conjunto de s elementos falle para las cuatro permutaciones es a lo sumo

\[U_s=\binom ns\left[\binom nr\binom{s+r}s\big/\binom ns\right]^4.\]

Verificamos \(U_{s}\le 2^{-s}\). Escribimos \(A=\binom{n}{r}\), \(B=\binom{s+r}{s}\), \(\alpha=s/n\) y \(x=(s+r)/n\). Como \(B\le A\) y \(B/\binom{n}{s}\le x^{s}\) (comparando los factores de los dos coeficientes binomiales),

\[U_s\le A^5x^{3s}.\]

La estimación elemental \(\binom{n}{r}\le (3n/r)^{r}\) para \(r\ge 1\) da \(A^{20}\le (60/\alpha)^{s}\) cuando \(20r\le s\). Una forma de comprobar esta última implicación es utilizar la monotonía de \((60n/t)^{t}\) para \(0<t\le n\), en \(t=20r\) y \(t=s\). El caso \(r=0\) da directamente la misma desigualdad. Además,

\[x^{20}\le(21/20)^{20}\alpha^{20}\le3\alpha^{20}.\]

Elevar a la potencia veinte la cota para \(U_{s}\) da

\[U_s^{20}\le[60^5\cdot27\cdot\alpha^{55}]^s\le(2^{-20})^s,\]

porque \(\alpha\le 1/2\) y \(60^{5}\cdot 27\le 2^{35}\). Sumar \(2^{-s}\) en el intervalo finito \(1\le s\le \lfloor n/2\rfloor\) da un número estrictamente menor que uno. Por tanto, existe una tupla de permutaciones que sirve para todos esos conjuntos. Para cada S, alguna permutación tiene al menos \(r+1\ge s/20\) imágenes salientes. Éstas dan aristas distintas de la frontera simple, incluso tras descartar repeticiones entre permutaciones. Tomar complementos demuestra (F.1) para todos los conjuntos. Cada permutación aporta a lo sumo dos aristas incidentes por vértice, lo que da grado a lo sumo ocho. La desigualdad de corte también implica conexidad.

\subsubsection{F.3. Un relleno con signos acotado en el anfitrión cíclico completo}\label{f.3.-un-relleno-con-signos-acotado-en-el-anfitriuxf3n-cuxedclico-completo}

Sea Q un grafo con tres clases A,B,C, cada una indexada por \(\mathbb Z/n\mathbb Z\), y todos los arcos cíclicos entre clases consecutivas. Fijamos una orientación del grafo G de la Sección F.2 sobre el conjunto de índices. Para índices a,b,c, escribimos

\[t(a,b,c)=[A_aB_b]+[B_bC_c]+[C_cA_a].\]

Utilizamos las siguientes circulaciones fijas con signos como caras:

\begin{enumerate}
\def\labelenumi{\arabic{enumi}.}
\tightlist
\item
  los \(n^{2}\) triángulos latinos t(a,b,a+b);
\item
  para cada arista orientada \(f\to s\) de G y cada a, los dos movimientos
\end{enumerate}

\[\begin{aligned}m_B(e,a)&=t(a,s,a+f)-t(a,f,a+f),\\m_C(e,a)&=t(a,f-a,s)-t(a,f-a,f);\end{aligned}\]

\begin{enumerate}
\def\labelenumi{\arabic{enumi}.}
\setcounter{enumi}{2}
\tightlist
\item
  en cada uno de los tres bloques de clases consecutivas, los rectángulos \((e_s-e_f)(e_{s\prime}-e_{f\prime})^{\mathsf T}\) indexados por dos aristas orientadas de G.
\end{enumerate}

Cada cara es no nula, tiene coeficientes en \{-1,0,1\} y a lo sumo cuatro arcos no nulos. Los movimientos tienen cuatro arcos, no seis, porque se cancela un lado triangular común. El único cambio de perfil saliente de \(m_{B}\) es \(e_{s}-e_{f}\) en la clase B; el de \(m_{C}\) es \(e_{s}-e_{f}\) en la clase C. Cada rectángulo tiene sumas de filas y columnas nulas y, por tanto, perfil saliente cero en cada clase.

Hay a lo sumo \(57n^{2}\) caras: G tiene a lo sumo \(4n\) aristas, por lo que los tres conteos son \(n^{2}\), a lo sumo \(8n^{2}\) y a lo sumo \(48n^{2}\). Para un arco fijo, las cotas del número de caras incidentes en los cuatro grupos (triángulos, movimientos B, movimientos C y rectángulos) son, respectivamente,

\[3,\quad32,\quad32,\quad256.\tag{F.4}\]

Para los movimientos, elegir cuál de sus dos extremos variables coincide con un índice fijo reduce el conteo a las aristas incidentes de G, con a lo sumo cuatro elecciones de este tipo; fijada esta arista, el índice restante queda determinado. Para los rectángulos, elegimos un extremo de cada arista orientada, obteniendo la cota segura \((2\cdot 8)^{2}\). La cota tres para las caras latinas es deliberadamente holgada. Estos conteos siguen siendo válidos si distintas caras indexadas tienen el mismo vector de arcos.

Sea L cualquier subgrafo equilibrado de Q. Denotamos sus grados salientes en las tres clases por \(d_{A},d_{B},d_{C}\). Sus sumas coinciden y cada grado es a lo sumo n.~Elegimos

\[z_{ab}=\begin{cases}1,&\text{si el representante de }b-a\pmod n\text{ es }<d_A(a),\\0,&\text{en otro caso.}\end{cases}\]

La suma de cadenas \(z_{ab}\) t(a,b,a+b) tiene el perfil A requerido. Sus errores de perfil B y C tienen suma cero y valor absoluto a lo sumo n.~Por la Sección F.1, cada error tiene un flujo entero en G con coeficientes de valor absoluto menor que \(20n+1\). Dividimos cada coeficiente en n sumandos enteros de valor absoluto a lo sumo 21. Los utilizamos como coeficientes de \(m_{B}(e,a)\) y \(m_{C}(e,a)\), respectivamente. Las dos correcciones sólo afectan a sus perfiles respectivos. Obtenemos una circulación con signos Y con exactamente el perfil saliente de L en las tres clases.

Por (F.4), el uso absoluto de caras sobre cualquier arco en esta construcción es a lo sumo

\[3+2\cdot21\cdot32=1347.\]

En consecuencia, el residuo \(X=1_{L}-Y\) tiene entradas acotadas por 1348. Todos sus perfiles salientes se anulan. Al ser una circulación, sus perfiles entrantes también se anulan. Cada matriz de clases consecutivas tiene, por tanto, filas y columnas de suma cero. Aplicamos (F.2) a cada una de estas tres matrices con ambas constantes de expansión iguales a 1/20. Los coeficientes de rectángulos son enteros de valor absoluto menor que

\[2\cdot1348\cdot20^2+2=1078402,\]

y, por tanto, a lo sumo 1078401. Combinando esto con (F.4), obtenemos un relleno con signos de \(1_{L}\) cuyo uso absoluto total sobre cada arco es a lo sumo el entero fijo

\[T=3+2\cdot21\cdot32+1078401\cdot256.\tag{F.5}\]

Las caras y T se fijaron antes de L. Los coeficientes del relleno dependen de L.

\subsubsection{F.4. Convertir rellenos con signos en un catálogo positivo}\label{f.4.-convertir-rellenos-con-signos-en-un-catuxe1logo-positivo}

Fijemos D subconjunto de Q. Para cada arco \(e=u\to v\) de Q, introducimos T caminos positivos de u a v de longitud cuatro y T caminos negativos de v a u de longitud cinco, todos con vértices internos privados. Los caminos distintos tienen vértices internos y arcos distintos. Sus longitudes tienen los residuos apropiados módulo tres, por lo que todos son cíclicos. Sea R su unión, disjunta de D.

El catálogo consta de tres tipos de piezas positivas equilibradas:

\begin{itemize}
\tightlist
\item
  para cada cara, cada signo y cada elección de un canal sobre cada uno de sus arcos no nulos, sustituimos cada arco con signo por un camino en su dirección indicada;
\item
  para cada arco de Q, la unión de cualquier canal positivo con cualquier canal negativo;
\item
  para cada e en D, e junto con el canal negativo de índice cero de e.
\end{itemize}

Sustituir una circulación entera por caminos dirigidos preserva el equilibrio. Cada pieza enumerada es no vacía y contiene a lo sumo 20 arcos. En particular, no se afirma que el primer tipo sea un único ciclo dirigido; el equilibrio es la propiedad utilizada por la construcción posterior.

Demostramos cobertura exacta, no sólo cobertura de L. Dado L equilibrado subconjunto de D, tomamos el relleno de la Sección F.3 y expandimos sus coeficientes en apariciones unitarias con signos. En un arco e, sean \(p_{e},n_{e}\) sus conteos positivos y negativos. Entonces

\[p_e-n_e=\mathbf1_L(e),\qquad p_e+n_e\le T.\]

Para e en L, seleccionamos primero la pieza del tercer tipo que utiliza el canal negativo cero. Asignamos canales positivos distintos a las \(p_{e}\) apariciones positivas y canales negativos restantes distintos a las \(n_{e}\) apariciones negativas. Hay espacio suficiente, pues \(n_{e}+1_{L}(e)=p_{e}\le T\). Tanto los canales positivos como los negativos no asignados son \(T-p_{e}\), por lo que los emparejamos y seleccionamos las piezas del segundo tipo. Las piezas del primer tipo realizan las apariciones de caras con signo. Cada arco de \(R\cup L\) se utiliza ahora exactamente una vez. Las piezas así seleccionadas son distintas, pues son no vacías y disjuntas en arcos. Forman, por tanto, un subconjunto del catálogo fijo, no un multiconjunto que requiera copias adicionales de un arco físico.

El catálogo tiene tamaño a lo sumo \(C_{\rm card}n^2\) y carga por vértice a lo sumo \(C_{\rm load}n\), donde las constantes congeladas son

\[\begin{aligned}C_{\rm card}&=114T^4+9T^2+9,\\C_{\rm load}&=(646T^4+T^2+1)(120T+22),\\C_d&=C_{\rm card}+C_{\rm load}.\end{aligned}\tag{F.6}\]

Para la primera estimación hay a lo sumo \(2\cdot 57n^{2}\cdot T^{4}\) piezas de caras, \(9n^{2}\cdot T^{2}\) piezas de canales emparejados y \(9n^{2}\) piezas de entrada. Las cotas \(9n^{2}\) cuentan todos los pares ordenados, por lo que sobreestiman de forma segura los arcos de Q.

Para la estimación de carga, un arco físico pertenece a lo sumo a \(2\cdot 323\cdot T^{4}+T^{2}+1\) piezas del catálogo. Un vértice base incide con a lo sumo \(6n\) arcos de referencia. Contar a lo sumo cinco arcos por canal y permitir una ficha de canal adicional para un vértice privado da a lo sumo

\[2[5(2T\cdot6n+1)+6n]\le(120T+22)n.\]

arcos físicos incidentes, para \(n\ge 1\). Contar dos veces los pares (pieza, arco incidente) demuestra (F.6). Estas cotas deliberadamente holgadas cubren tanto los vértices base como los privados. Cualquier digrafo cíclico D sobre N vértices se inserta en Q con N índices por clase; añadir vértices preserva los originales. La fórmula (F.6) da una reserva y un catálogo fijos con número de piezas \(\le C_dN^2\) y carga \(\le C_dN\). Para \(N=0\) usamos la reserva y el catálogo vacíos. Esto demuestra la construcción densa con piezas de tamaño veinte; no es una hipótesis adicional.

\subsubsection{F.5. Reducir una entrada de grado diecisiete a carga de orden raíz cuadrada}\label{f.5.-reducir-una-entrada-de-grado-diecisiete-a-carga-de-orden-rauxedz-cuadrada}

Sea Y cíclico con \(d_{Y}(v)\le 17\) y \(m=|Y|>0\). Pongamos

\[s=\lfloor\sqrt{2m}\rfloor+1,\quad M=307+290s,\quad A_0=9\cdot597^2\cdot8.\tag{F.7}\]

Asignamos el soporte de Y a M grupos de cada color, cada uno con a lo sumo s vértices, de modo que arcos distintos de Y tengan pares ordenados de grupos distintos. Procesamos los vértices uno a uno. Además de la capacidad de los grupos, mantenemos que los vértices a distancia no dirigida a lo sumo dos reciben índices de grupo diferentes. En un vértice nuevo, los índices de grupo prohibidos están acotados por:

\begin{itemize}
\tightlist
\item
  \(17+17^{2}=306\) por vértices a distancia a lo sumo dos;
\item
  \(17^2s=289s\) por arcos existentes cuyo par de grupos coincidiría con el de un arco hacia un vecino procesado anteriormente;
\item
  a lo sumo s-1 grupos llenos, pues el soporte procesado tiene tamaño menor que \(s^{2}\).
\end{itemize}

Para el segundo conteo, un grupo contiene a lo sumo s vértices antiguos, cada uno incidente con a lo sumo diecisiete arcos; hay a lo sumo diecisiete vecinos pertinentes. Los arcos nuevos distintos no pueden coincidir entre sí porque sus extremos antiguos tienen índices de grupo diferentes. Tampoco pueden coincidir las orientaciones opuestas en el vértice nuevo: los colores de sus extremos son distintos. Como M supera la suma de las tres exclusiones, siempre queda un grupo disponible. Esto establece la asignación afirmada, en lugar de limitarse a suponer un cociente simple.

Sea p(v) el grupo de v y sea D el cociente cíclico, con un centro por grupo. La aplicación e-\textgreater(p(tail e),p(head e)) es inyectiva. Hay \(N=3M\) centros y a lo sumo \(17s\) arcos de entrada incidentes con cada centro. Como \(s^{2}\le 8m\) y \(M\le 597s\),

\[N^2\le A_0m,\qquad(\text{incidencia del centro})^2\le A_0m.\tag{F.8}\]

Aplicamos la Sección F.4 a D. Para cada arco original \(e=u\to v\) introducimos cinco vértices privados y los siguientes caminos:

\[\begin{aligned}P_e&:v\to x_0\to u,\\O_e&:u\to x_1\to x_2\to p(u),\\I_e&:p(v)\to x_3\to x_4\to v,\\Q_e&:p(u)\to p(v).\end{aligned}\]

Los colores de p(u),p(v) son los de u,v. Por tanto, los caminos indicados son cíclicos. El ciclo base \(P_{e}+O_{e}+Q_{e}+I_{e}\) tiene longitud nueve. La reserva consta de todos estos ciclos base y de la reserva densa sobre los centros. Su catálogo contiene el catálogo denso y:

\begin{itemize}
\tightlist
\item
  el triángulo \(e+P_{e}\), para cada e en Y;
\item
  el ciclo base de longitud nueve para cada e;
\item
  \(I_{e}+O_{f}\), un ciclo de longitud seis, siempre que \(head(e)=tail(f)\).
\end{itemize}

Dado L equilibrado subconjunto de Y, seleccionamos el ciclo base para cada e fuera de L y el triángulo para cada e en L. En cada vértice antiguo emparejamos sus arcos entrantes de L con sus arcos salientes de L y utilizamos los ciclos de seis correspondientes. Los arcos restantes son exactamente la reserva densa más Q(L). Este último es equilibrado, pues la aplicación de grupos preserva las sumas de divergencias, y es un subgrafo de D por inyectividad. La Sección F.4 particiona este resto. Todas las selecciones son disjuntas en arcos y todas las piezas pertenecen al catálogo fijado antes de L.

Hay a lo sumo \(17m\) pares ordenados compatibles (e,f), por lo que el tamaño total del catálogo es a lo sumo (19+C\_d A\_0)m. El siguiente cálculo de carga registra la cota más fuerte que necesitaremos para los vértices que se identificarán. En un vértice w, sea \(S_{w}\) la suma de sus grados en todos los ciclos base de nueve. Los triángulos contribuyen a lo sumo \(S_{w}+d_{Y}(w)\), los ciclos base a lo sumo \(S_{w}\) y los ciclos de seis a lo sumo \(34S_{w}\): un arco puede ser el primer o segundo miembro de a lo sumo diecisiete pares. El catálogo denso contribuye a lo sumo \(C_{d}\) N y no contiene vértices antiguos. En particular, \(S_{v}=2d_{Y}(v)\) en un vértice antiguo y, por tanto,

\[\operatorname{load}(v)\le73d_Y(v)\quad\text{para cada vértice antiguo}.\tag{F.9}\]

En un vértice general, sea ell su incidencia como centro si es un centro, y cero en otro caso. El número de ciclos base que lo contienen es a lo sumo 17+ell, lo que da

\[\operatorname{load}(w)\le1241+(72+C_d)(\ell+N).\]

Por (F.8), \((\ell+N)^2\le(2A_0)^2m\). Como \(m\ge 1\), la constante congelada

\[C_m=19+C_dA_0+1241+(72+C_d)(2A_0).\tag{F.10}\]

acota, por tanto, el tamaño del catálogo por \(C_{m}\) m y da \(\operatorname{load}(w)^2\le C_m^2m\) en cada vértice. Cada pieza tiene tamaño a lo sumo veinte. Para \(m=0\), el catálogo vacío satisface las mismas afirmaciones. La distinción entre (F.9) y la cota de orden raíz cuadrada es esencial: una cota de orden raíz cuadrada por puerto sería insuficiente tras asignar muchos puertos a una misma raíz física.

\subsubsection{F.6. Una red fija de ocho permutaciones en cada raíz}\label{f.6.-una-red-fija-de-ocho-permutaciones-en-cada-rauxedz}

Para cada conjunto finito P existen ocho permutaciones tales que, para todo S subconjunto de P, el número total de canales dirigidos de permutación que salen de S, contados con sus etiquetas de permutación, es al menos

\[\min\{|S|,|P\setminus S|\}.\tag{F.11}\]

Aquí los canales paralelos se conservan como canales distintos; no es el grafo simple de la Sección F.2. Para \(|P|\le 1\) la afirmación es inmediata. Sea \(n=|P|\) y fijemos S con \(1\le s\le n/2\), \(p=s/n\). Para una permutación uniforme, sea I el número de imágenes de S que permanecen en S. Sus momentos factoriales satisfacen

\[\mathbb E\binom Ik=\frac{\binom sk^2}{\binom nk}\le\binom skp^k.\]

Así, E \(\lambda^{I}\) \(\le (1+p(\lambda-1))^{s}\) para \(\lambda\ge 1\). Si ocho permutaciones independientes tienen menos de s canales salientes, su conteo interno total supera \(7s\). La desigualdad de Markov con \(\lambda=4/p\), seguida de \(\binom{n}{s}\le (3/p)^{s}\), acota la probabilidad de un conjunto de s elementos que falle por

\[\left[3(5-p)^8p^6/4^7\right]^s\le2^{-s}.\]

Para la última desigualdad, \(p^{3}(5-p)^{4}\) es creciente para \(0<p\le 1/2\) y es menor que 52 en ese intervalo, mientras que \(3\cdot 52^{2}<4^{7}/2\). La suma sobre s es estrictamente menor que uno. Por tanto, alguna tupla sirve para todos los conjuntos pequeños. El conteo saliente a través de un corte coincide con el de su complemento para cada permutación, lo que demuestra (F.11).

Descartamos los puntos fijos y damos capacidad unitaria a cada canal restante. Los grados de entrada y salida son ambos a lo sumo ocho. Si los enteros \(b_{p}\) satisfacen \(|b_{p}|\le 1\) y suman \(b_{p}=0\), entonces

\[\left|\sum_{p\in S}b_p\right|\le\min\{|S|,|P\setminus S|\}.\]

Por (F.11), se cumplen todas las condiciones de corte dirigido para encaminar este desequilibrio. Por tanto, un flujo entero de capacidad unitaria selecciona canales con la divergencia requerida. Esta última afirmación también se deduce directamente añadiendo una fuente y un sumidero para las demandas positivas y negativas y aumentando el flujo por caminos residuales: si el aumento se detiene antes de encaminar toda la demanda, el conjunto alcanzable desde la fuente viola precisamente una de estas desigualdades de corte. Las capacidades enteras mantienen cada canal seleccionado con valor cero o uno. La red sirve para todos esos b; no se elige después de conocer b.

\subsubsection{F.7. Puertos de raíces y catálogo completo}\label{f.7.-puertos-de-rauxedces-y-catuxe1logo-completo}

Sea X un digrafo cíclico arbitrario y pongamos \(x=|X|\). Dividimos cada arco de entrada en dos puertos distintos, uno asignado a su raíz de origen y otro a su raíz de destino. La entrada levantada es un emparejamiento sobre \(2x\) puertos. Un puerto conserva el color de su raíz. En cada raíz aplicamos la Sección F.6 a su conjunto de puertos. Sustituimos cada canal de la red elegida por un camino privado de longitud tres cuyos extremos son sus puertos. Los extremos tienen el mismo color, por lo que el camino es cíclico. Sea N la unión de todos los caminos de red e I la entrada levantada.

Cada puerto tiene grado total de red a lo sumo dieciséis, cada vértice interno de red tiene grado dos y el emparejamiento levantado añade a lo sumo uno. Por tanto,

\[\max_v d_{N\cup I}(v)\le17,\qquad |N\cup I|\le49x.\tag{F.12}\]

En efecto, hay a lo sumo ocho canales por puerto, cada uno con tres arcos, de modo que \(|N|\le 24\cdot (2x)\). Las redes quedan fijadas para X. Si L subconjunto de X es equilibrado, su levantamiento tiene divergencia en \{-1,0,1\} en cada puerto y la suma de divergencias en cada fibra de raíz es cero. La Sección F.6 selecciona canales completos cuyos caminos A cancelan esas divergencias. Así, A unión I(L) es equilibrado, con A subconjunto de N.

Aplicamos una vez la Sección F.5 a \(Y=N\cup I\), obteniendo una reserva R disjunta de Y y su catálogo fijo. Ampliamos este catálogo con cada camino individual de canal de N. Dado L, usamos primero la partición mesoscópica para A unión I(L). Después añadimos los caminos individuales de canal no seleccionados en A. Esto particiona

\[R\cup N\cup I(L).\tag{F.13}\]

exactamente, utilizando un subconjunto del catálogo fijo ampliado.

Un camino de canal individual añadido no es equilibrado sobre el conjunto de puertos divididos. Sus dos extremos pertenecen a la misma raíz. Tras identificarlos con esa raíz, sus contribuciones entrante y saliente se cancelan. Todas las demás piezas ya eran equilibradas. La inserción física debe preservar arcos distintos; ese paso separado de la Sección 8 no se sustituye por un cociente arbitrario del grafo abstracto. Bajo una realización inyectiva en arcos, el equilibrio desciende sumando divergencias sobre las fibras, y las particiones de (F.13) también descienden.

Hay a lo sumo \(16x\) piezas adicionales de canales. Su carga en un vértice antiguo es a lo sumo su grado de red, pues los caminos son disjuntos en arcos. Las ecuaciones (F.9) y (F.12) dan, por tanto, la cota uniforme en vértices antiguos

\[\operatorname{load}(v)\le73\cdot17+16=1257.\tag{F.14}\]

Para un vértice mesoscópico realmente nuevo u no hay piezas adicionales de canales, por lo que \(load(u)^{2}\le C_{m}^{2}\cdot 49x\). Las cotas restantes son

\[\begin{aligned}|\mathcal C|&\le(49C_m+16)x,\\\sum_{P\in\mathcal C}|P|&\le(980C_m+48)x.\end{aligned}\tag{F.15}\]

y cada pieza tiene a lo sumo veinte arcos. La segunda cota mantiene separados los tamaños veinte y tres; no es necesario sustituir ambos por veinte.

Hay exactamente \(d_{X}(r)\) puertos asignados a una raíz física r. La precarga total del catálogo transportada por esos puertos es, por tanto, a lo sumo

\[1257d_X(r).\tag{F.16}\]

y no 1257 veces el número total de arcos de entrada. Los vértices internos de red no están fijados a r. Éste es el presupuesto de raíces que utiliza la inserción física ponderada. Tomar L vacío en (F.13) también demuestra que el soporte del catálogo contiene R unión N, condición necesaria al definir sus pesos de arcos.

Las ecuaciones (F.12)--(F.16) completan la construcción del catálogo abstracto. Junto con una realización física inyectiva en arcos y el absorbedor terminal demostrado en el Apéndice A, proporcionan la entrada abstracta necesaria para la construcción de familias con leyes de conservación. Esta afirmación se refiere a esa entrada; las pruebas separadas PAR/BAL y semidensa se desarrollan en los Apéndices E y C, respectivamente.

\subsubsection{F.8. Por qué las piezas satisfacen las condiciones terminales de divisibilidad}\label{f.8.-por-quuxe9-las-piezas-satisfacen-las-condiciones-terminales-de-divisibilidad}

Para el núcleo de tres colores, el único tipo triangular es 012. En un vértice de color i de una pieza cíclica equilibrada P, los números de vecinos de colores i-1 e i+1 coinciden: son sus grados entrante y saliente. Su perfil local es, por tanto, un múltiplo entero del perfil de esquina de 012. Sumar el equilibrio sobre cada clase de color también muestra que los números de arcos en los tres bloques de clases consecutivas coinciden. En consecuencia, el vector global de conteos de aristas es un múltiplo entero de (1,1,1). Así, P satisface ambas condiciones de divisibilidad entera; el argumento no se basa sólo en la divisibilidad del número total de aristas.

Tras la sustitución física de este artículo, vale el mismo razonamiento porque cada pieza es equilibrada y cíclica, incluidos los canales adicionales cuyos extremos se identificaron con su raíz común. Cada pieza tiene a lo sumo ochenta aristas físicas. Por tanto, el resultado de certificados enteros acotados del Apéndice A puede utilizarse con la cota única de tamaño ochenta. Su certificado se construye en el anfitrión completo permitido; no debe afirmarse que todos sus triángulos auxiliares pertenecen ya al anfitrión incompleto. En su lugar, la construcción terminal se inserta mediante extensión. Su grafo no tiene aristas entre dos raíces antiguas y su construcción ordena los vértices nuevos con a lo sumo cuatro vecinos expuestos previamente, como se demuestra en el Apéndice A. Las desigualdades de disponibilidad y evitación de ese apéndice son las condiciones utilizadas para esta inserción.

Finalmente, aplicamos el lema de ensamblaje de catálogo fijo al catálogo físico. Sus entradas son el número de piezas y la carga de todo el catálogo, no la carga de una selección residual concreta. Las construcciones terminales y sus inserciones físicas se fijan para todas las piezas antes de seleccionar L. La partición para L elige entonces una de las dos descomposiciones de cada construcción ya insertada. Esto explica exactamente dónde intervienen (F.15), (F.16) y la estimación de carga física; no introduce un absorbedor nuevo dependiente de L. La elección cuantitativa de las escalas del anfitrión sigue siendo la de la construcción congelada y del lema de realización física de la Sección 8.

\Needspace{12\baselineskip}\subsection{Apéndice G. Enunciados auxiliares y sus aplicaciones}\label{apuxe9ndice-g.-enunciados-auxiliares-y-sus-aplicaciones}

Las pruebas anteriores utilizan los resultados auxiliares siguientes. Enunciamos las condiciones necesarias para las aplicaciones, de modo que puedan comprobarse por separado de la construcción del absorbedor tipado. Estas versiones cuantitativas son resultados del desarrollo complementario {[}8{]}, no citas de un teorema general de descomposición sin especificar. La Sección 8 del documento complementario da las declaraciones y ubicaciones exactas. El Apéndice H demuestra matemáticamente la biblioteca compatible, las posiciones robustas, el refinador ordinario, el empaquetamiento y la selección de reserva, incluidas sus cotas numéricas. Las reducciones de cobertura y caminos ponderados se demuestran a continuación.

\subsubsection{G.1. Anti-aristas}\label{g.1.-anti-aristas}

Una anti-arista para raíces distintas \(u,v\) y etiqueta de completación \(k\) es un camino de dos aristas \(uwv\), donde \(p(w)=k\). No se exige que exista la arista \(uv\). Un grafo tiene \emph{disponibilidad de anti-aristas \(\alpha\) para \(F\)} si, para toda presentación ordenada \(\{i,j,k\}\in F\) y todo par de vértices distintos \(u\in V_i\), \(v\in V_j\), existen al menos \(\alpha n\) vértices \(w\) de ese tipo. Esta es una condición de vecinos comunes tipados, no una afirmación de que falte la arista \(uv\). La disponibilidad protegida de orden al menos dos la implica tomando las dos raíces como vecinos prescritos. En la aplicación semidensa sólo se solicitan las etiquetas de completación de \(F_0\); no se requieren consultas arbitrarias de vecinos comunes internos.

\subsubsection{G.2. Empaquetamiento y selección de una reserva descomponible}\label{g.2.-empaquetamiento-y-selecciuxf3n-de-una-reserva-descomponible}

El método es el nibble clásico {[}16--18{]}; la comparación global entre empaquetamientos fraccionarios y enteros {[}19{]} y el resultado ponderado de emparejamientos {[}20{]} proporcionan contexto relacionado. La cota local del residuo y las constantes explícitas siguientes se demuestran en H.7--H.8, sin inferirlas de esas referencias omitiendo comprobar sus hipótesis.

\textbf{Enunciado auxiliar de empaquetamiento.} Para \(b>0\) y \(0<\eta\le1\) existen una tolerancia positiva \(\tau(\eta)\) y un umbral entero \(N_{\rm pack}(b,\eta)\). Para \(\xi>0\), pongamos \(\zeta=\min\{\eta,1,\xi^2\}\). Si \(n\ge N_{\rm pack}(b,\zeta)\), \(\Delta(J)\le(1-\xi)n\) y los triángulos permitidos de \(J\) tienen pesos en \([0,b/n]\) cuya carga en cada arista difiere de uno en a lo sumo \(\tau(\zeta)\), entonces \(J\) tiene un empaquetamiento de triángulos permitidos que deja grado máximo a lo sumo \(\eta n\). Con cargas unitarias exactas existe también un enunciado sin la hipótesis de margen en el grado, con umbral dependiente sólo de \(b,\eta\).

Estas funciones se eligen uniformemente antes del conjunto de etiquetas, el anfitrión y los pesos. La implementación del enunciado de cargas próximas a uno permite \[N_{\rm pack}(b,\zeta)\le
\left\lceil2^{6500}(b+1)^{40}\zeta^{-4000}\right\rceil
\quad(0<\zeta\le1).\] La tolerancia y el umbral utilizados en B.6 son los de este enunciado. La Sección H.8 especifica la tolerancia y demuestra la positividad y la cota superior mostrada. Los pesos por sí solos no implican ninguna afirmación sobre el retículo entero.

\textbf{Enunciado auxiliar de reserva.} Supongamos que hay a lo sumo cuatro etiquetas, \(0<\rho,\omega\le1\), \(G\) tiene disponibilidad de anti-aristas \(\omega\) y \(Q\subseteq E(G)\) es el conjunto de aristas de un empaquetamiento de triángulos permitidos con \(\Delta(E(G)\setminus Q)\le\omega n/8\). Si \[n\ge2^{27}\rho^{-9}\omega^{-6},\] existe una reserva descomponible \(X\subseteq E(G)\) con \(\Delta(X)\le2\rho n\) y el grafo sobre \(X\) tiene disponibilidad de anti-aristas \(\rho^2\omega/12\). La selección toma triángulos completos de la descomposición suministrada. Aplicar primero el enunciado de empaquetamiento con cargas próximas a uno requiere la cota adicional \(\Delta(G)\le(1-\omega)n\); por eso se realiza la comprobación de B.6.

\subsubsection{G.3. Cobertura hacia la reserva con el absorbedor ya fijado}\label{g.3.-cobertura-hacia-la-reserva-con-el-absorbedor-ya-fijado}

\textbf{Enunciado auxiliar de cobertura hacia la reserva.} Fijemos \(b>0\) y \(0<\alpha\le1\). Existe un umbral que depende sólo de esos parámetros. Sean \(A,X\subseteq E(G)\) disjuntos. Supongamos que \(G\) es enteramente \(F\)-divisible, que \(A\cup L\) es descomponible para todo \(L\subseteq X\) enteramente divisible y que el grafo sobre \(X\) tiene disponibilidad de anti-aristas \(\alpha\). Si \(J=G-(A\cup X)\) tiene cargas triangulares unitarias exactas con pesos en \([0,b/n]\), entonces \(G\) tiene una descomposición en \(F\) por encima de ese umbral.

Para ver la reducción, tomemos un empaquetamiento en \(J\) que deje \(\Delta(L_0)\le\eta n\), donde \(\eta=\alpha^2/256\). Pongamos \(r=\lceil\eta n\rceil\), \(q=\lfloor\alpha n/16\rfloor\) y \(m=|L_0|\le nr/2\). Para \(n\ge256/\alpha^2\), las cotas \(q\ge\alpha n/32\), \(q\le\alpha n/16\) y \(r\le2\eta n\) dan \[\frac m q+2(r+2q)<\alpha n.\] Cubrimos sucesivamente las aristas de \(L_0\) mediante triángulos cuyas otras dos aristas pertenecen a \(X\). Excluimos los vértices ya utilizados \(q\) veces como centro nuevo y los vecinos de reserva ya gastados de los dos extremos. Los tres términos de la desigualdad pagan exactamente esas exclusiones. Cada arista descubierta tiene un tipo de completación, porque su carga unitaria original era positiva. La reserva no utilizada \(L\subseteq X\) es enteramente divisible: restamos el empaquetamiento y los triángulos nuevos de \(E(G)\) y usamos que \(A\) es descomponible (tomando \(L=\varnothing\) en su propiedad de absorción). El absorbedor fijo resuelve ese \(L\). No es necesario elegirlo después del empaquetamiento.

La variante efectiva con cargas próximas a uno de {[}8{]} utiliza \(\tau(\min\{\alpha^2/256,1,\alpha^2\})\) y umbral \[\max\left\{N_{\rm pack}(b,\min\{\alpha^2/256,1,\alpha^2\}),
\left\lceil256/\alpha^2\right\rceil\right\}.\] Las cargas exactas satisfacen automáticamente su tolerancia. Su margen de grado se deduce de la reserva disjunta con anti-aristas, no se impone como una hipótesis nueva en el enunciado de cobertura.

\subsubsection{G.4. Refinadores ordinarios y casillas robustas}\label{g.4.-refinadores-ordinarios-y-casillas-robustas}

El resultado de refinador ordinario de B.3 supone que \(X\) no tiene lazos, que \(\Delta(X)\le2\rho n\) y las cuatro desigualdades escalares (B.3.1). Su conclusión es un único resto y una biblioteca fija, con las cotas allí indicadas: tamaño de pieza 42, grado del resto \(4\lambda n\), a lo sumo \(1224\rho n^2\) índices y carga con multiplicidades \(4096\lambda n\). Todo subgrafo par de \(X\) con número de aristas divisible por tres tiene una selección de esa misma biblioteca. El anfitrión completo virtual proporciona los caminos que necesita la construcción de montículo; B.3 demuestra después que la biblioteca se puede enrutar en el anfitrión real. Este resultado ordinario no proporciona una inserción en el anfitrión real ni divisibilidad tipada global.

El resultado de casillas robustas tiene otra función. Para enteros \(0\le s\le z\), existen \(m\le3z\) casillas fijas y un grafo bipartito entre \(m+s\) vértices izquierdos y \(m\) vértices derechos fijos más \(z\) flexibles. Los grados izquierdos son a lo sumo 24, los derechos a lo sumo 18 y toda selección de exactamente \(s\) casillas flexibles admite un emparejamiento perfecto con todos los vértices izquierdos. La Sección H.5 demuestra este enunciado finito mediante permutaciones y el teorema de Hall. En B.4 cada casilla flexible representa una unidad de carga, no un vértice del anfitrión. Junto con las dos funciones directas de las fichas, da las cotas de multiplicidad y de índices nuevos utilizadas en la construcción por capas. Los factores 9, 44 y 97 corresponden a esos conteos del registro, no a grados del grafo bipartito robusto.

\Needspace{20\baselineskip}\subsubsection{G.5. Selección ponderada de caminos}\label{g.5.-selecciuxf3n-ponderada-de-caminos}

Este es el paso voraz preciso utilizado en la Sección 8. Supongamos que la disponibilidad protegida para dos vecinos en tres etiquetas distintas es al menos \(a\), las aristas prohibidas tienen grado máximo a lo sumo \(d\ge0\) y una función de carga no negativa tiene total a lo sumo \(S\) y cota puntual \(T+2W\), donde \(T>0\). Si \[S/T+4+2d<a,\] entonces, para raíces de dos etiquetas distintas y \(0\le w\le W\), se puede elegir un camino de cuatro aristas con la tercera etiqueta prescrita en su vértice central. Sus tres vértices internos son nuevos, evita las aristas prohibidas y cargar \(2w\) a cada uno conserva la cota puntual y aumenta la carga total en \(6w\).

En efecto, a lo sumo \(S/T\) vértices tienen carga al menos \(T\). Elegimos el primer y el último vértice interno fuera de ese conjunto, las raíces y los vecinos prohibidos; elegimos el vértice central en el vecindario común de los dos vértices elegidos, con la tercera etiqueta. En cada elección hay a lo sumo cuatro exclusiones por distinción y a lo sumo \(2d\) exclusiones por vecinos prohibidos. La desigualdad estricta deja una elección. Cada carga interna nueva es a lo sumo \(T+2W\). Las cargas de los extremos se asignan por separado antes de este paso, como requiere la Sección 8.

\clearpage

\subsection{Apéndice H. Demostraciones de las construcciones auxiliares}\label{apuxe9ndice-h.-demostraciones-de-las-construcciones-auxiliares}

Este apéndice proporciona las construcciones auxiliares utilizadas en C.5a, G.2 y G.4; la reducción de cobertura se presenta por separado en G.3. En las cadenas, «conector» traduce wire; «enlace» designa el grafo asociado a una etiqueta raíz. Una biblioteca es una familia indexada: dos índices se conservan aunque sus conjuntos de aristas coincidan. Su carga es la suma de los grados sobre todos los índices. Cada resto y cada biblioteca se fijan antes de elegir el residuo divisible.

\subsubsection{H.1. Cadenas de paridad}\label{h.1.-cadenas-de-paridad}

Para bits \(x_0,\ldots,x_{N-1}\), ponemos \(c_{-1}=0\) y \(c_j=c_{j-1}+x_j\) en \(\mathbb F_2\). Entre posiciones consecutivas colocamos dos caminos distintos \(A_j,B_j\) con extremos comunes. La posición \(j\) recibe \(B_{j-1}\) cuando \(c_{j-1}=1\); si no es la última, recibe además \(A_j\) y recibe \(B_j\) cuando \(c_j=0\). Cada camino del conector se asigna una vez. En una posición no terminal, el número recibido tiene paridad \[c_{j-1}+1+(1-c_j)=x_j\pmod2.\] En la última posición la paridad es \(x_{N-1}+c_{N-1}\). Hay como máximo tres caminos por posición. Esto incluye una cadena de longitud uno, que no tiene conectores.

Aplicamos esta regla a una familia de conectores \(\Lambda\) y una aplicación \(h\) que no fija ningún extremo de conector. Para cada vértice \(u\), formamos una cadena de sus conectores incidentes con extremos de camino \(u,h(u)\). Para cada par proyectado no degenerado \(\sigma=\{h(u),h(v)\}\), formamos una cadena de sus preimágenes, con extremos de camino \(\sigma\), y dos caminos terminales \(f_\sigma,f'_\sigma\). Los caminos con los mismos extremos siguen siendo peticiones distintas.

Para cada conector incluimos las dieciséis alternativas determinadas por su bit de selección y los tres acarreos entrantes. Una alternativa contiene el conector seleccionado, los caminos recibidos de sus dos cadenas de extremos y los de su cadena de par proyectado. En la última posición de esta cadena añadimos \(f_\sigma\) si el acarreo saliente es uno. Incluimos también la pieza \(f_\sigma\cup f'_\sigma\). Descartamos las alternativas con frontera no nula módulo dos.

Para una entrada seleccionada par, las cadenas de extremos tienen acarreo terminal cero. Sus fronteras cancelan los extremos de entrada y dejan los extremos proyectados; las cadenas de pares cancelan estos últimos. Un par proyectado degenerado contribuye cero. Los pares proyectados con un número impar de preimágenes forman una nueva entrada par, pues su frontera es la imagen de la frontera nula anterior. Si el acarreo final de un par es cero, seleccionamos \(f_\sigma\cup f'_\sigma\); si es uno, \(f_\sigma\) ya se utiliza y \(f'_\sigma\) es el siguiente conector de entrada. Esto demuestra la cobertura disjunta entre niveles consecutivos. Una jerarquía que termina sin conectores refina, por tanto, toda entrada par mediante una biblioteca fija.

Escribimos \(S=\sum_\ell|\Lambda_\ell|\). Hay como máximo \(17S\) índices y \(8S\) peticiones: por nivel, los conectores de extremos cuestan a lo sumo \(4|\Lambda_\ell|\), los de pares \(2|\Lambda_\ell|\) y los terminales \(2|\Lambda_\ell|\). Una pieza tiene como máximo \(1+3+3+4=11\) átomos. Un átomo aparece en no más de \(33\) piezas, pues un conector encuentra como máximo dos posiciones con dieciséis alternativas cada una y un terminal puede aparecer además en su pieza cíclica. Realizar cada petición mediante su propio camino de dos aristas da \[\begin{gathered}
|P_i|\le22,\\
\sum_i|P_i|\le66(8S+|X|).
\end{gathered}\] En una raíz la carga es a lo sumo \(33(d_X+\text{carga de raíces de peticiones})\); en un vértice intermedio nuevo es a lo sumo \(66\). La carga de raíces de peticiones de un nivel en \(u\) es a lo sumo \[2d_\Lambda(u)+10\sum_{v:h(v)=u}d_\Lambda(v).\] El primer término cuenta sus conectores de extremos; el segundo paga los conectores de extremos entrantes, los de pares y los terminales. Estas estimaciones cuentan todas las alternativas.

\subsubsection{H.2. La jerarquía comprimida y las etiquetas compatibles}\label{h.2.-la-jerarquuxeda-comprimida-y-las-etiquetas-compatibles}

Para C.5a introducimos \(2M\) identidades de concentradores, disjuntas de los vértices originales, donde \(M=2^a\). Enviamos el vértice original \(j\) al concentrador \(M+(j\bmod M)\) y los concentradores posteriores \(i\) a \(\lfloor i/2\rfloor\). Los niveles positivos tienen tamaños \(M,M/2,\ldots,1\); el último no tiene enlaces no degenerados. Así, \[S\le|X|+2M^2.\] Un vértice original tiene carga de raíces de peticiones a lo sumo \(2d_X(v)\) y, por tanto, carga de biblioteca a lo sumo \(99d_X(v)\). Un concentrador recibe como máximo \(10(\lfloor n/M\rfloor+1)\Delta_X\) de la primera aplicación. En un nivel positivo de tamaño \(w\), su propio grado es a lo sumo \(w\) y tiene como máximo dos hijos, lo que da carga de raíces a lo sumo \(22w\). Sumar los tamaños da la cota de carga de biblioteca del concentrador \[330(\lfloor n/M\rfloor+1)\Delta_X+1452M.\]

Conservamos las etiquetas originales y damos etiqueta cero a los concentradores. Una petición con extremos de igual etiqueta recibe un vértice intermedio de la etiqueta opuesta; una con extremos de etiquetas distintas recibe uno de etiqueta cero. Cada terna solicitada pertenece así a \(F_0=\{001,011\}\). Estas etiquetas se fijan antes de corregir las cargas algebraicas.

Para un conjunto de aristas etiquetado definimos \[\chi(E)=2e_{01}(E)-4e_{00}(E)-4e_{11}(E).\] Los grados pares implican \(4\mid\chi(E)\), pues el número de aristas cruzadas es par. Un grafo par con carga cero es integralmente \(F_0\)-divisible: sus coeficientes globales son \(e_{00}\) y \(e_{11}\) y sus perfiles locales pertenecen a las redes generadas por \((1,1),(0,2)\) o sus inversiones de coordenadas. Estas son precisamente las redes de suma par.

Utilizamos un ciclo de cuatro vértices con etiquetas \(0011\) como ficha de carga \(-4\), y un ciclo de ocho con etiquetas \(01011001\) como ficha de carga \(4\). Tienen respectivamente dos y seis aristas cruzadas, y cada uno tiene dos aristas internas. Cada ficha tiene grado máximo dos y no más de ocho aristas, sobre identidades nuevas.

\subsubsection{H.3. El registro de cargas y sus recuentos}\label{h.3.-el-registro-de-cargas-y-sus-recuentos}

Para las piezas pares ponemos \(c_i=\chi(P_i)/4\), de modo que \(|c_i|\le|P_i|\). Toda selección para un residuo divisible tiene la misma suma \(K=\chi(R_{\rm par})/4\). Ponemos \[p_i=\max(-c_i,0),\quad m_i=\max(c_i,0),\quad
z=\sum_i(p_i+m_i),\quad s=\sum_i m_i-K.\] Una pieza parental de demanda positiva recibe \(m_i\) fichas negativas variables. Una de demanda negativa recibe \(p_i\) fichas positivas privadas, cada una con una compañera negativa variable. Incluimos la pieza de liberación formada por cada pareja. Para una selección dada, utilizamos sus piezas parentales y las piezas de liberación de las parentales negativas no seleccionadas. Quedan exactamente \(s\) fichas negativas variables: pertenecen a parentales positivas no seleccionadas o son compañeras de parentales negativas seleccionadas. La selección del residuo vacío muestra que \(0\le s\le z\).

Usamos H.5 para fijar un grafo con \(b\le3z\) posiciones negativas fijas, \(b+s\) posiciones positivas y \(z\) posiciones negativas flexibles. Introducimos las fichas auxiliares correspondientes y, por cada arista del grafo, una pieza de cancelación formada por sus dos fichas. Para cualquier residuo, sus \(s\) posiciones variables restantes admiten un emparejamiento perfecto con las posiciones positivas y las negativas fijas. Las piezas parentales, de liberación y de emparejamiento cubren cada ficha una vez y cubren \(R_{\rm par}\cup L\). Cada pieza seleccionada es par y tiene carga cero. Esto demuestra la universalidad.

Hay como máximo \(9z\) fichas y, si la biblioteca anterior tiene \(N\) índices, como máximo \[N+z+24(b+s)\le N+97z\] índices nuevos. Una ficha aparece como máximo \(1+1+24+18=44\) veces, contando sus posibles funciones parental, de liberación y de emparejamiento. Las cargas anteriores se conservan; los vértices de fichas tienen carga a lo sumo \(88\). Las piezas parentales tienen tamaño a lo sumo \(9|P_i|\) y las de cancelación a lo sumo \(16\). Por tanto, la cota uniforme de tamaño es \(9\cdot22+16=214\), y \[17S+97\cdot66(8S+|X|)
=51233S+6402|X|\le57635(|X|+2M^2).\]

El mismo recuento explica B.4. Si una capa tiene cota \(A\) para los coeficientes y el orden finito, sus como máximo \(AN\) índices vacíos de relleno contribuyen demanda a lo sumo \(A^2N\), además de los \(AN\) originales. Los índices vacíos no tienen carga de aristas. Así, la demanda es a lo sumo \(A(A+1)N\), lo que da \(9A(A+1)N\) fichas y \[N+AN+97A(A+1)N\le L_jN,\qquad L_j=1+98A(A+1).\] La desigualdad tiene holgura \(A^2N\). Si la unión de fichas tiene carga a lo sumo \(\Delta\), una capa transforma la carga de biblioteca \(\Lambda\) en como máximo \(\Lambda+44\Delta\). La siguiente cota inductiva incluye otro \(\Delta\), por lo que \[L_{\rm next}(\Lambda+45\Delta)
\le45L_jL_{\rm next}(\Lambda+\Delta).\] El número de fichas restante es \(9A(A+1)N+T_{\rm next}L_jN\). Estos son los factores y las recurrencias hacia atrás de B.4; acotan la biblioteca completa.

\subsubsection{H.4. Una realización física para todos los residuos}\label{h.4.-una-realizaciuxf3n-fuxedsica-para-todos-los-residuos}

Colocamos primero los vértices originales, luego los concentradores, los vértices intermedios de peticiones y los vértices de fichas en orden cíclico. Cada identidad nueva tiene como máximo dos vecinos anteriores. Las consultas de cierre de los dos ciclos de fichas tienen etiquetas \((0,1,1)\) y \((0,0,1)\), respectivamente. Los vecinos anteriores de cada consulta de dos raíces aparecen juntos en una pieza, por lo que la inyectividad por piezas hace distintas sus imágenes. Todas las consultas de dos raíces están permitidas. La abundancia positiva de antiaristas y \(n\ge3\) fuerzan la presencia de ambas etiquetas; añadir una raíz de etiqueta opuesta proporciona también cada consulta necesaria de una raíz. La consulta vacía tiene la misma cota de oferta.

Escribimos \(B=214\), sea \(I\) una cota del número de índices y sea \(\Lambda(v)\le L_{\max}\le q\) la carga de biblioteca abstracta. Al colocar \(v\) cargamos su vértice físico con \(\Lambda(v)\) y utilizamos solo vértices cuya carga acumulada sea menor que \(q\). Después su carga es a lo sumo \(2q\). Hay como máximo \(2BI/q\) vértices saturados, pues la suma total de grados es a lo sumo \(2BI\). Como máximo \(2BL_{\max}\) identidades anteriores aparecen junto a la nueva identidad: cada pieza tiene a lo sumo \(2B\) vértices no aislados y el número de piezas que la contienen es a lo sumo su carga. Excluir sus imágenes conserva la inyectividad en cada pieza.

En cada uno de los como máximo dos vecinos anteriores excluimos los vecinos del grafo prohibido, de \(X\) y de la unión de aristas ya realizadas. Sus cotas de grado son \(d,\Delta_X,2q\). Por tanto, la exclusión total es a lo sumo \[2BL_{\max}+\frac{2BI}{q}+2(d+\Delta_X+2q)<\alpha n.\] Queda un candidato. Esto demuestra la inducción, incluida la inyectividad de aristas de todo el soporte. La aplicación única fija raíces y etiquetas; transportar cualquier selección abstracta conserva la disjunción, la cobertura exacta y la divisibilidad. H.1--H.4 demuestran así C.5a, incluidas sus constantes.

\subsubsection{H.5. El grafo de posiciones robustas}\label{h.5.-el-grafo-de-posiciones-robustas}

Para \(k\ge1\), tomamos \(3k\) vértices izquierdos y dos conjuntos derechos \(Y,Z\) de tamaño \(2k\) cada uno. Damos dos puertos a cada vértice izquierdo y tres a cada derecho. Para cada conjunto derecho elegimos, independientemente, seis permutaciones uniformes de los \(6k\) puertos y unimos sus propietarios. Los grados son a lo sumo \(24\) a la izquierda y \(18\) a la derecha.

Un conjunto izquierdo de tamaño \(a\le k\) tiene todos sus vecinos en \(Y\) dentro de un conjunto dado de tamaño \(a\) con probabilidad a lo sumo \((a/(2k))^{12a}\). Un conjunto derecho de tamaño \(b\le2k\) tiene todos sus vecinos en un conjunto izquierdo dado de tamaño \(b\) con probabilidad a lo sumo \((b/(3k))^{18b}\). Esto utiliza muestreo sin reemplazo, cuya probabilidad está acotada por la potencia correspondiente. Las cotas de unión son a lo sumo \(1/15\) cada una, pues \(\binom Na\le(3N/a)^a\) da \[\begin{gathered}
\binom{3k}{a}\binom{2k}{a}(a/(2k))^{12a}\le16^{-a},\\
\binom{4k}{b}\binom{3k}{b}(b/(3k))^{18b}\le16^{-b}.
\end{gathered}\] En efecto, las bases son \((27/2)x^{10}\le1/16\) para \(x\le1/2\) y \(12y^{16}\le1/16\) para \(y\le2/3\). Existe, por tanto, un grafo en que ambos vecindarios exceden estrictamente el tamaño del conjunto considerado.

Para cualquier subconjunto \(Z'\) de \(Z\) de tamaño \(k\), se cumple la condición de Hall entre el lado izquierdo y \(Y\cup Z'\). Los conjuntos izquierdos pequeños utilizan la expansión hacia \(Y\). Un conjunto mayor \(A\) que incumpliese Hall dejaría un conjunto derecho \(B\) de tamaño mayor que \(3k-|A|\), con vecinos en su complemento. Si \(|B|\le2k\), la expansión derecha da una contradicción; en caso contrario se aplica a un subconjunto de tamaño \(2k\), cuyos vecinos tendrían que caber en menos de \(2k\) vértices izquierdos. Esto demuestra la existencia del emparejamiento perfecto. Para \(k=0\) utilizamos el grafo vacío.

Para \(0\le s\le z\) arbitrarios, tomamos \(k=\max(z-s,s)\), \(p=\max(z-2s,0)\) y \(b=2k+p\). Conservamos \(z\) posiciones flexibles, convertimos las siguientes \(p\) en fijas y descartamos las demás. Las identidades \(s+p=k\), \(b+s=3k\) y \(b\le3z\) demuestran el enunciado requerido al extender cada subconjunto elegido de tamaño \(s\) con las posiciones fijas.

\subsubsection{H.6. El refinador ordinario y sus constantes}\label{h.6.-el-refinador-ordinario-y-sus-constantes}

Trabajamos primero sobre \(N\) vértices con parámetros locales \(\lambda,\rho\). Elegimos una potencia de dos \[64\rho N/\lambda\le M\le128\rho N/\lambda.\] Las hipótesis de B.3 dan \(2\le M\) y \(2M\le N\). Enumeramos los vértices desde cero. La primera aplicación es \(j\mapsto M+((j+1)\bmod M)\); no tiene puntos fijos y sus fibras tienen tamaño a lo sumo \(\lfloor N/M\rfloor+1\). Las aplicaciones posteriores son los padres de la jerarquía binaria. H.1 da \(S\le|X|+2M^2\le2\rho N^2\). Añadimos como máximo dos ciclos disjuntos de cuatro aristas para hacer divisible por tres el tamaño del resto y los incluimos como piezas separadas. La biblioteca par tiene entonces \(N_0\le17S+2\le34\rho N^2+2\) índices.

Entre posiciones consecutivas de esta biblioteca colocamos tres ciclos nuevos de cuatro aristas. Una alternativa \((x,a,b)\in\{0,1\}\times\{0,1,2\}^2\) toma la pieza anterior cuando \(x=1\), los ciclos numerados al menos \(a\) en su conector izquierdo y los numerados por debajo de \(b\) en su conector derecho. Omitimos los conectores de los extremos y conservamos las alternativas de tamaño divisible por tres. Elegimos los acarreos como el negativo del recuento acumulado de aristas seleccionadas módulo tres. Cada pieza tiene tamaño congruente con \(x|P_i|-a+b=0\); los acarreos inicial y final son cero. Cada ciclo del conector se asigna una vez. Se conserva la paridad y el tamaño de pieza es a lo sumo \(22+12+8=42\).

Damos al conector \(j\) raíces \(j\bmod N,(j+1)\bmod N\). Cada uno de sus ciclos utiliza dos peticiones de caminos separadas. Si \(D=\lfloor2\rho N\rfloor\), todas las raíces de peticiones tienen carga a lo sumo \[r_0=2D+10(\lfloor N/M\rfloor+1)D+44M+8
+12(\lfloor(N_0+1)/N\rfloor+1).\] Hay como máximo \(8S+4+6N_0\le512\rho N^2\) peticiones. Las hipótesis escalares dan \[\rho\le\lambda/2^{30},\quad
\lfloor N/M\rfloor D\le\lambda N/32,\quad
44M\le\lambda N/2000,\quad
\lfloor(N_0+1)/N\rfloor\le34\rho N+1.\] Por tanto, \(r_0\le432\rho N+(5/16+1/2000)\lambda N+32\le\lambda N\), usando \(\lambda N\ge256\).

Ponemos \(q=\lfloor\lambda N/64\rfloor\ge\lambda N/128\). Realizamos las peticiones sucesivamente, evitando vértices intermedios ya utilizados \(q\) veces y aristas de \(X\) o de peticiones anteriores. El número de exclusiones es a lo sumo \[\frac{512\rho N^2}{q}+2(2\rho N+r_0+2q)
\le\frac{65536\rho N}{\lambda}+4\rho N+\frac{33}{16}\lambda N
<8\lambda N.\] Así, basta una oferta de \(8\lambda N\) vecinos comunes. El grado del resto es a lo sumo \(r_0+2q\le4\lambda N\). Un átomo anterior aparece como máximo \(18\cdot33\) veces; un conector módulo tres aparece como máximo \(36\) veces. En consecuencia, \[\sum_i d_{P_i}(v)\le594d_X(v)+630(r_0+2q)\le4096\lambda N,
\qquad |I|\le18N_0\le1224\rho N^2.\] Para el subconjunto \(W\) de B.3 usamos \(N=|W|\), \(\lambda'=\lambda n/N\) y \(\rho'=\rho n/N\). Se conservan sus cuatro hipótesis escalares y el anfitrión completo sobre \(W\) proporciona \(N-2\ge8\lambda n\) vecinos comunes. Finalmente, \(1224\rho'N^2\le1224\rho n^2\), mientras que las cotas de grado y carga se reescalan exactamente. Esto demuestra el enunciado virtual utilizado.

\subsubsection{H.7. Progreso, conflictos y cargas supervivientes en una ronda de empaquetamiento}\label{h.7.-progreso-conflictos-y-cargas-supervivientes-en-una-ronda-de-empaquetamiento}

Sean triángulos físicos candidatos distintos de un conjunto activo de aristas \(B\), con pesos \(0\le w_i\le\beta\) y cargas de aristas en \([1-\delta,1+\delta]\). Suponemos \(\delta\le1/4\), \(0<\gamma\le1/128\) y \(\gamma\beta\le1/4\). Retenemos el triángulo \(i\) independientemente con probabilidad \(\gamma w_i\) y conservamos solo los que no comparten aristas con otro retenido. Si \(Q\) es el conjunto de aristas tocadas, ponemos \(B'=B\setminus Q\) y \(W=Q\setminus E(\mathcal P)\). Entonces \[B\setminus E(\mathcal P)=B'\mathbin{\dot\cup}W.\]

Utilizamos la siguiente cota en el espacio producto. Si cambiar la coordenada de Bernoulli \(i\), de probabilidad \(p_i\), cambia \(f\) a lo sumo \(c_i\le L\), cualquiera de las dos colas a distancia \(t\) de su esperanza es a lo sumo \[\exp\left(-\frac{t^2}{2(V+Lt)}\right),\qquad V=\sum_i p_ic_i^2.\] Esta es una estimación de concentración de tipo Bernstein; véase {[}21, Capítulo 2{]} para el método de momentos exponenciales. Para completar el argumento, en la martingala de exposición un incremento tiene la forma \((X_i-p_i)d_i\), con \(|d_i|\le c_i\). Su varianza condicional es a lo sumo \(p_ic_i^2\) y sus momentos superiores están acotados por \(L^{k-2}\) veces esa varianza. Expandir la exponencial, multiplicar las cotas condicionales y aplicar la desigualdad de Markov da el denominador de Bernstein \(2(V+Lt/3)\) y, por tanto, la cota más débil mostrada. También se aplica tras fijar una colección disjunta de coordenadas.

Fijamos \(v\) y ponemos \(d=d_B(v)\). Sea \(R_v\) el número de triángulos retenidos por \(v\). Sumar las cargas de aristas da \(\mathbb ER_v\ge3\gamma d/8\), de donde \[\mathbb P(R_v<\gamma d/4)\le e^{-\gamma d/64}.\] Para cada triángulo \(U\) que no contiene \(v\), sea \(b_U\) el peso total de los triángulos por \(v\) que comparten una arista con \(U\), y sea \(Y_v=\sum_U b_UX_U\). Hay como máximo tres de esos triángulos por cada \(U\), por lo que \(b_U\le3\beta\). Recíprocamente, un triángulo exterior solo bloquea un triángulo por \(v\) en su arista opuesta. Intercambiar las sumas da \(\mathbb EY_v\le\gamma(1+\delta)^2d/2\le\gamma d\). Por tanto, \[\mathbb P(Y_v\ge\gamma d+t)
\le\exp\left(-\frac{t^2}{2(3\beta\gamma d+3\beta t)}\right).\] Sea \(Z_v\) el número de triángulos retenidos por \(v\) bloqueados desde fuera. Condicionalmente a las monedas exteriores, su esperanza bloqueada es a lo sumo \(\gamma Y_v\). Las monedas restantes son independientes. Con parámetro exponencial \(1/2\), la desigualdad \(e^{1/2}-1\le3/4\) da \[\mathbb P(Y_v\le y,\ Z_v\ge z)\le e^{-z/2+3\gamma y/4}.\]

Para \(u\ne v\), sea \(D_u\) el número de triángulos retenidos que contienen \(vu\), y ponemos \(F_v=\sum_u\binom{D_u}{2}\). Tenemos \(\mathbb ED_u\le\gamma(1+\delta)\le1/2\), de modo que \(\mathbb P(D_u\ge K)\le e^{-K/8}\) para \(K\ge1\). La independencia da \(\mathbb EF_v\le\gamma^2d\). Truncamos a \(\widehat F_v=\sum_u\binom{\min(D_u,K)}2\). Una moneda por \(v\) cambia dos sumandos a lo sumo \(K\) cada uno; las exteriores no influyen. Así, \[\mathbb P(\widehat F_v\ge\gamma^2d+t)
\le\exp\left(-\frac{t^2}{2(4K^2\gamma d+2Kt)}\right).\] Cuando todos los \(D_u<K\), el truncamiento coincide con \(F_v\). Un triángulo retenido por \(v\) que se descarta está bloqueado desde fuera o pertenece a un conflicto interno, lo que da las desigualdades deterministas \[\begin{gathered}
R_v\le d_{\mathcal P}(v)/2+Z_v+2F_v,\\
d_W(v)\le2Z_v+4F_v.
\end{gathered}\] Para \(d\ge\theta n\), utilizamos los umbrales \(\gamma d/4,d/32,\gamma d/16,\gamma d/16\) para \(R_v,Y_v,Z_v,\widehat F_v\). Para el descarte utilizamos \(2\gamma n,4\gamma^2n,4\gamma^2n\) para las tres últimas variables. Las respectivas desviaciones en las cotas de \(Y\) y \(\widehat F\) son \(d/64,\gamma d/24,\gamma n,2\gamma^2n\). Fuera de los sucesos adversos resultantes, \[\begin{gathered}
d_{\mathcal P}(v)\ge\gamma d/8\quad(d\ge\theta n),\\
d_W(v)\le24\gamma^2n.
\end{gathered}\] Su probabilidad de fallo por vértice es a lo sumo la suma de las ocho cantidades siguientes: \[\begin{gathered}
e^{-\gamma\theta n/64},\quad
e^{-\theta n/[24576\beta(\gamma+1/64)]},\quad
e^{-\gamma\theta n/128},\quad
e^{-\gamma\theta n/(4608K^2+96K)},\\
e^{-\gamma n/(12\beta)},\quad e^{-\gamma^2n/2},\quad
e^{-\gamma^3n/(2K^2+2K\gamma)},\quad ne^{-K/8}.
\end{gathered}\] Por ejemplo, el exponente de triángulos bloqueados es \(-\gamma d/32+3\gamma d/128=-\gamma d/128\). Sustituir \(t=\gamma d/24\) en la cota de conflictos truncados da el denominador \(4608K^2+96K\). Los demás términos se obtienen mediante las sustituciones indicadas.

Para una arista \(e\), condicionamos únicamente a que no haya sido tocada. Todas las monedas de triángulos por \(e\) son entonces cero y las demás conservan su ley producto. Para un triángulo \(i\) por \(e\), la suma de pesos de sus competidores exteriores es la suma de las cargas de sus otras dos aristas menos \(2w_i\), y pertenece por tanto a \([2(1-\delta-\beta),2(1+\delta)]\). Usando \(e^{-x-2x^2}\le1-x\le e^{-x}\) para \(0\le x\le1/4\), su probabilidad de supervivencia está entre \[e^{-2\gamma(1+\delta)-5\gamma^2\beta}
\quad\text{y}\quad e^{-2\gamma(1-\delta-\beta)}.\] Si \(f_e\) es la carga superviviente, sumar estas cotas y usar \(e^x\le1+2x\) para \(0\le x\le1\) da \[|e^{2\gamma}\mathbb E(f_e\mid e\text{ no tocada})-1|
\le(1+5\gamma)\delta+5\gamma\beta.\] Una moneda exterior cambia \(f_e\) a lo sumo \(3\beta\). Intercambiar las sumas de su influencia ponderada por la probabilidad de retención da como máximo \(2\gamma(1+\delta)^2\) y, por tanto, una cota de varianza de \(10\beta\gamma\). Como \(e^{2\gamma}\le2\), un error adicional mayor que \(\varepsilon\) tiene probabilidad condicional a lo sumo \[2\exp\left(-\frac{(\varepsilon/2)^2}
{2(10\beta\gamma+3\beta\varepsilon/2)}\right).\] Multiplicar por la probabilidad de que \(e\) no haya sido tocada solo disminuye esta cota. Aplicamos la cota de unión a las ocho estimaciones por vértice sobre \(n\) vértices y a esta estimación sobre como máximo \(n^2\) aristas. Así se demuestran simultáneamente las tres conclusiones de la ronda, sin condicionar a los sucesos favorables de los vértices.

\subsubsection{H.8. Iteración y umbral explícito}\label{h.8.-iteraciuxf3n-y-umbral-expluxedcito}

Ponemos \[\begin{gathered}
\zeta=\min\{\eta,1,\xi^2\},\quad \gamma=\zeta^2/2^{16},\quad
r=\left\lceil8\log(2/\zeta)/\gamma\right\rceil,\\
a=1+5\gamma,\quad c=(8ra^r)^{-1},\quad
\theta=\zeta/2,\quad\varepsilon=c/2,\quad B_{\max}=be^{2\gamma r}.
\end{gathered}\] Elegimos \(\tau(\zeta)=c\). Las desigualdades \(\log(2/\zeta)\le2/\zeta\) y \(\gamma r\le8\log(2/\zeta)+\gamma\) dan \[\begin{gathered}
r<2^{21}\zeta^{-3},\quad a^r\le2^{41}\zeta^{-40},\quad
c\ge\zeta^{43}/2^{65},\quad B_{\max}\le2^{17}(b+1)\zeta^{-16},\\
(1-\gamma/8)^r\le\zeta/2,\qquad24r\gamma^2\le\zeta/4.
\end{gathered}\] Para la última estimación, \(24r\gamma^2\le192\gamma\log(2/\zeta)+24\gamma^2\le384\zeta/2^{16}+24\zeta^4/2^{32}<\zeta/4\).

Tomamos \(K=\lfloor\sqrt{\lfloor\sqrt n\rfloor}\rfloor\) y \(\beta\le B_{\max}/n\). Entonces \(K^4\le n<16K^4\). Con \[\kappa=\frac{\zeta^{102}}{2^{160}(b+1)},\] cada exponencial de H.7 es a lo sumo \(e^{-\kappa K}\). Para la estimación de aristas, \(\varepsilon\ge\zeta^{43}/2^{66}\) y \(8(10B_{\max}\gamma+3B_{\max}\varepsilon/2)\le128B_{\max}\) dan una tasa al menos \(\zeta^{102}/[2^{156}(b+1)]\). Para los demás términos utilizamos las cotas mostradas y \(4608K^2+96K\le4704K^2\), \(2K^2+2K\gamma\le4K^2\). La cota de unión es a lo sumo \(2560K^8e^{-\kappa K}\).

Como \(e^{\kappa K}>(\kappa K)^9/9!\), esta cantidad es menor que uno si \(K>2560\cdot9!/\kappa^9\). Basta que \[n\ge\left\lceil2^{6500}(b+1)^{40}\zeta^{-4000}\right\rceil.\] En efecto, \(K>n^{1/4}/2\), esta última cota inferior es al menos \(2^{1624}(b+1)^{10}\zeta^{-1000}\) y el lado derecho requerido es menor que \(2^{1470}(b+1)^9\zeta^{-918}\). Este umbral también da \(\gamma B_{\max}/n\le1/4\) y \(5\gamma B_{\max}/n\le c/2\).

Para cargas iniciales exactas comenzamos en tiempo cero. En tiempo \(j\) multiplicamos los pesos base por \(e^{2\gamma j}\) sobre los triángulos supervivientes. Conservamos el conjunto activo \(B_j\) y todas las aristas tocadas descartadas \(W_j\). Inductivamente, el empaquetamiento evita ambos, cada arista no cubierta pertenece a su unión y \[\begin{gathered}
\begin{gathered}
\Delta(W_j)\le24\gamma^2nj,\\
d_{B_j}(v)\le\max\{\theta n,(1-\gamma/8)^jn\},\\
|\operatorname{load}_j(e)-1|\le cja^j.
\end{gathered}
\end{gathered}\] Para cargas iniciales próximas a uno, usamos en cambio pesos base \(e^{-2\gamma}w\) y comenzamos en tiempo uno sin empaquetamiento. Los pesos efectivos son entonces \(w\). La separación de grado da \((1-\xi)n\le(1-\gamma/8)n\) y el error inicial \(c\) es a lo sumo \(ca\). En cada paso posterior \(cja^j\le cra^r=1/8\) y la cota de peso es \(B_{\max}/n\). H.7 aumenta el error a lo sumo a \(a(cja^j)+c\le c(j+1)a^{j+1}\), el grado descartado a lo sumo en \(24\gamma^2n\) y disminuye cada estrella activa grande en su fracción \(\gamma/8\). Las pequeñas solo disminuyen. La identidad de residuo disjunto conserva el invariante de empaquetamiento. En tiempo \(r\), el grado máximo no cubierto es a lo sumo \(\zeta n/2+\zeta n/4\le\eta n\). Esto demuestra G.2, incluida la versión separada de cargas exactas, obtenida con \(\zeta=\min\{\eta,1\}\) y tiempo cero.

\subsubsection{H.9. Selección de triángulos completos para la reserva}\label{h.9.-selecciuxf3n-de-triuxe1ngulos-completos-para-la-reserva}

Seleccionamos cada triángulo del empaquetamiento dado independientemente con probabilidad \(\rho\). En un vértice hay como máximo \(n/2\) bloques, por lo que el grado seleccionado excede \(2\rho n\) con probabilidad a lo sumo \(e^{-\rho n/8}\). Para una consulta permitida de dos raíces, al menos \(3\omega n/4\) testigos tienen ambas aristas cubiertas por el empaquetamiento: cada estrella no cubierta elimina como máximo \(\omega n/8\). Un testigo sobrevive con probabilidad \(\rho\) o \(\rho^2\), según sus aristas pertenezcan a uno o dos bloques.

Un bloque influye en como máximo dos testigos. Si su influencia es \(c_i\), entonces \(c_i\le2\), \(\sum_i c_i\le2n\) y \(\sum_i c_i^2\le4n\). La cota de varianza de exposición es a lo sumo \(4\rho n\). La desigualdad de Bernstein con \(t=\rho^2\omega n/2\) muestra que tener menos de \(\rho^2\omega n/4\) testigos supervivientes tiene probabilidad a lo sumo \(e^{-\rho^3\omega^2n/48}\). Así se contabiliza explícitamente la dependencia por bloques compartidos.

Hay como máximo \(4n^2\) consultas y \(n\) condiciones de grado. Ponemos \(c_0=\rho^3\omega^2/48\). La probabilidad de fallo es a lo sumo \(5n^2e^{-c_0n}\). Si \(n\ge2^{27}\rho^{-9}\omega^{-6}\), entonces \(c_0^3n>30\) y \(e^{c_0n}>(c_0n)^3/6\) hace esta cota menor que uno. El grafo seleccionado es descomponible y tiene la cota de grado requerida; su oferta de consultas supera incluso \(\rho^2\omega n/12\). Junto con las reducciones explícitas de G.3 y G.5, esto completa las demostraciones auxiliares.

Copyright © 2026 Juan Pablo Traverso Gianini. Este documento se distribuye bajo \href{https://creativecommons.org/licenses/by/4.0/}{CC BY 4.0}.

\end{document}
